% !TeX program  = pdflatex
% !TeX encoding = UTF-8
% !TeX spellcheck = de_DE
% !TeX root     = Thesis.tex
%
% ===========================================================================
% E-Portfolio Hochschule Bremen
% Bachelorthesis, Omar Ziyada und Mohamad Alaskari
% Hauptdatei: Praeambel, Layout und Reihenfolge der Bestandteile
% ===========================================================================
%
% Kompilieren in TeXstudio:
%   F5  (pdflatex)             -> Text und Struktur
%   F8  (Bibliographie, biber) -> Literaturverzeichnis
%   F5, F5                     -> Seitenzahlen in den Verzeichnissen
% Oder in einem Schritt auf der Kommandozeile:
%   latexmk -pdf Thesis.tex
% ===========================================================================

\documentclass[
a4paper,
12pt,
oneside,
parskip=false,       % Absaetze durch Einzug, nicht durch Leerzeile
numbers=noenddot,    % 1.1 statt 1.1.
chapterprefix=false, % Ueberschrift "1 Einleitung", nicht "Kapitel 1"
listof=totoc,        % Abbildungs- und Tabellenverzeichnis ins Inhaltsverzeichnis
bibliography=totoc   % Literaturverzeichnis ebenso
]{scrreprt}

% ---------------------------------------------------------------------------
% Angaben fuer Deckblatt und PDF-Metadaten
% Quelle: Genehmigung des Pruefungsamtes vom 17.06.2026 und
% Confluence-Seite "Eckdaten" (BEH, Seiten-ID 2129939)
% ---------------------------------------------------------------------------
\newcommand{\ThesisTitel}{Entwicklung eines studentischen E-Portfolio-Portals für die Hochschule Bremen}
\newcommand{\ThesisArt}{Bachelorthesis}
\newcommand{\Hochschule}{Hochschule Bremen}
\newcommand{\Fakultaet}{Fakultät 4: Elektrotechnik und Informatik}
\newcommand{\Abteilung}{Abteilung: Informatik und Medieninformatik}
\newcommand{\Studiengang}{Internationaler Studiengang Medieninformatik (ISMI)}
\newcommand{\AutorEins}{Omar Ziyada}
\newcommand{\MatrikelEins}{5244029}
\newcommand{\AutorZwei}{Mohamad Alaskari}
\newcommand{\MatrikelZwei}{5244256}
\newcommand{\Erstpruefer}{Prof. Andreas Teufel}
\newcommand{\Zweitpruefer}{M.Sc. Jendrik Bulk}
\newcommand{\AbgabeOrt}{Bremen}
\newcommand{\AbgabeDatum}{6. Oktober 2026}

% ---------------------------------------------------------------------------
% Sprache, Zeichensatz, Schrift
% ---------------------------------------------------------------------------
\usepackage[utf8]{inputenc}
\usepackage[T1]{fontenc}
\usepackage[ngerman]{babel}
\usepackage{mathptmx}                 % Times als Brotschrift, wie in der Beispielarbeit
\usepackage[scaled=0.90]{helvet}      % Helvetica als serifenlose Schrift
\usepackage{courier}                  % Courier fuer technische Bezeichner
\usepackage{microtype}                % ruhigerer Blocksatz

% ---------------------------------------------------------------------------
% Satzspiegel und Zeilenabstand
% ---------------------------------------------------------------------------
\usepackage[
a4paper,
left=3.0cm,
right=2.5cm,
top=3.0cm,
bottom=3.0cm,
headsep=0.8cm,
footskip=1.2cm
]{geometry}

\usepackage{setspace}
\onehalfspacing

\setlength{\emergencystretch}{3em}    % verhindert ueberstehende Zeilen im Blocksatz
\clubpenalty=10000                    % keine Schusterjungen
\widowpenalty=10000                   % keine Hurenkinder
\displaywidowpenalty=10000

% ---------------------------------------------------------------------------
% Weitere Pakete
% ---------------------------------------------------------------------------
\usepackage{amsmath}
\usepackage{amssymb}                  % \square fuer die Ankreuzfelder in Anhang D
\usepackage{graphicx}
\graphicspath{{bilder/}}              % alle Abbildungen liegen in bilder/
\usepackage{booktabs}
\usepackage{tabularx}
\usepackage{longtable}
\usepackage{array}
\usepackage{csquotes}
\usepackage{enumitem}
\usepackage{xcolor}
\usepackage{caption}
\captionsetup{font=small, labelfont=bf, justification=raggedright, singlelinecheck=false}

% Tabellen: etwas Luft zwischen den Zeilen, linksbuendige Umbruchspalten
\renewcommand{\arraystretch}{1.15}
\newcolumntype{L}[1]{>{\raggedright\arraybackslash}p{#1}}
\newcolumntype{Y}{>{\raggedright\arraybackslash}X}
\newcolumntype{R}{>{\raggedleft\arraybackslash}X}

% ---------------------------------------------------------------------------
% Literaturverwaltung: biblatex im APA-Stil, Backend biber
% ---------------------------------------------------------------------------
\usepackage[
backend=biber,
style=apa,
language=ngerman,
sortlocale=de_DE,
uniquename=init,
backref=false
]{biblatex}
\addbibresource{literatur/literatur.bib}
\DeclareLanguageMapping{ngerman}{german-apa}
% "o. J." wie im Text und auf der Confluence-Seite Literaturverzeichnis,
% nicht das "o. D." von biblatex-apa
\DefineBibliographyStrings{ngerman}{nodate = {o.\,J\adddot}}
% Literaturverzeichnis linksbuendig, damit lange URLs nicht mit Leerraum
% im Blocksatz auseinandergezogen werden
\AtBeginBibliography{\raggedright}

% Abbildungen duerfen mehr Platz auf Textseiten einnehmen, damit seltener
% eine Seite mit nur einer Abbildung oder eine halbleere Seite entsteht
\renewcommand{\topfraction}{0.85}
\renewcommand{\bottomfraction}{0.7}
\renewcommand{\textfraction}{0.12}
\renewcommand{\floatpagefraction}{0.8}

% ---------------------------------------------------------------------------
% Ueberschriften: serif und fett, wie in der Beispielarbeit
% ---------------------------------------------------------------------------
\setkomafont{disposition}{\normalfont\bfseries}
\setcounter{secnumdepth}{3}   % nummeriert bis 1.1.1.1
\setcounter{tocdepth}{2}      % Inhaltsverzeichnis bis 1.1.1

% ---------------------------------------------------------------------------
% Kopf- und Fusszeile
% ---------------------------------------------------------------------------
\usepackage[automark]{scrlayer-scrpage}
\clearpairofpagestyles
\automark{chapter}
\ihead{\headmark}
\cfoot*{\pagemark}
\KOMAoptions{headsepline=true}
\setkomafont{pageheadfoot}{\normalfont\small}
\setkomafont{pagenumber}{\normalfont}

% ---------------------------------------------------------------------------
% Eigene Befehle
% ---------------------------------------------------------------------------
% Sichtbarer Platzhalter fuer noch nicht geschriebene Stellen.
% Vor der Abgabe darf kein einziger mehr im PDF stehen.
% Zum Ausblenden: die naechste Zeile auskommentieren.
\newcommand{\platzhalterAn}{}
\newcommand{\platzhalter}[1]{%
	\ifdefined\platzhalterAn
	\par\medskip
	\noindent\fcolorbox{gray!50}{gray!10}{%
		\parbox{\dimexpr\linewidth-2\fboxsep-2\fboxrule\relax}{%
			\small\itshape\textbf{Platzhalter:} #1}}%
	\par\medskip
	\fi}

% Abbildung mit reserviertem Platz.
% Die Datei wird in bilder/ gesucht; fehlt sie, erscheint ein Rahmen
% mit dem erwarteten Dateinamen, und der Umbruch stimmt schon jetzt.
% #1 Dateiname   #2 Breite als Faktor von \textwidth   #3 Label
% #4 Kurztitel fuer das Abbildungsverzeichnis          #5 vollstaendige Bildunterschrift
\newcommand{\abbildung}[5]{%
	\begin{figure}[htbp]
		\centering
		\IfFileExists{bilder/#1}{%
			\includegraphics[width=#2\textwidth,height=0.62\textheight,keepaspectratio]{#1}%
		}{%
			\fcolorbox{gray!50}{gray!08}{%
				\parbox[c][4.2cm][c]{\dimexpr#2\textwidth-2\fboxsep-2\fboxrule\relax}{%
					\centering\small\itshape Abbildung noch nicht abgelegt\\[0.5em]
					\normalfont\small\texttt{bilder/#1}}}%
		}
		\caption[#4]{#5}
		\label{#3}
	\end{figure}}

% ---------------------------------------------------------------------------
% hyperref immer zuletzt laden
% ---------------------------------------------------------------------------
\usepackage[hidelinks]{hyperref}
\hypersetup{
	pdftitle={\ThesisTitel},
	pdfauthor={\AutorEins, \AutorZwei},
	pdfsubject={\ThesisArt, \Hochschule},
	pdfkeywords={E-Portfolio, Hochschule Bremen, Blockeditor, Redaktionsworkflow},
	pdflang={de-DE}
}

% ===========================================================================
\begin{document}

	% =======================================================================
	% VORSPANN
	% =======================================================================
	\hypersetup{pageanchor=false}
	% ---------------------------------------------------------------------------
% Deckblatt
% Aufbau nach der Beispielarbeit BAThesis_2025_ISMI_Erol.
% Alle Angaben kommen aus den Makros in Thesis.tex und werden nur dort geaendert.
% ---------------------------------------------------------------------------
\begin{titlepage}
	\centering
	\singlespacing

	\vspace*{0.2cm}

	\includegraphics[width=0.35\textwidth]{hsb-logo.png}

	\vspace{1.2cm}

	{\bfseries \Hochschule\par}
	\vspace{0.5cm}
	{\Fakultaet\par}
	\vspace{0.5cm}
	{\Abteilung\par}
	\vspace{0.5cm}
	{\Studiengang\par}

	\vspace{1.6cm}

	{\bfseries\large \ThesisTitel\par}

	\vspace{1.5cm}

	{\bfseries \ThesisArt\par}

	\vspace{1.5cm}

	{Vorgelegt von:\par}
	\vspace{0.5cm}
	{\AutorEins\par}
	{Matrikelnummer: \MatrikelEins\par}
	\vspace{0.5cm}
	{\AutorZwei\par}
	{Matrikelnummer: \MatrikelZwei\par}

	\vspace{1.2cm}

	{Erstprüfer: \Erstpruefer\par}
	\vspace{0.5cm}
	{Zweitprüfer: \Zweitpruefer\par}

	\vfill

	{\AbgabeOrt, \AbgabeDatum\par}

	\vspace{0.5cm}
\end{titlepage}

	\hypersetup{pageanchor=true}

	\pagenumbering{roman}
	\pagestyle{scrheadings}

	% ---------------------------------------------------------------------------
% Eigenstaendigkeitserklaerung, KI-Angabe und Nachweis der Einzelbeitraege
% Uebernommen aus der Confluence-Seite "Eigenstaendigkeitserklaerung,
% KI-Angabe und Einzelbeitraege (Entwurf)" (Space BEH, Seiten-ID 101482498),
% Stand 05.10.2026, mit drei Korrekturen:
%   - Titel nach der Genehmigung des Pruefungsamtes vom 17.06.2026 (\ThesisTitel)
%   - KI-Angabe auf Wunsch der Verfassenden (05.10.2026) auf zwei Bereiche
%     beschraenkt: Softwareentwicklung nach eigenen Regeln, Skills und Agents
%     sowie Ausformulieren und grammatische Korrektur beim Schreiben
%   - Einzelbeitraege nach Angabe der Verfassenden (05.10.2026), ohne
%     Verweise auf Jira, Confluence, Commits oder Fassungen
% Blockierend: fehlt die Erklaerung bei der Abgabe, gilt die Thesis als nicht
% bestanden. Vor der Unterschrift gegen das Formblatt der Fakultaet 4 pruefen;
% die mit "zu bestätigen" markierten Angaben sind von den Verfassenden zu
% pruefen.
% ---------------------------------------------------------------------------
\addchap{Eigenständigkeitserklärung}

Hiermit versichern wir, \AutorEins{} (Matrikelnummer \MatrikelEins) und \AutorZwei{} (Matrikelnummer \MatrikelZwei), dass wir die vorliegende Bachelorarbeit mit dem Titel „\ThesisTitel“ selbstständig verfasst und keine anderen als die angegebenen Quellen und Hilfsmittel benutzt haben. Alle Stellen, die wörtlich oder sinngemäß aus veröffentlichten oder unveröffentlichten Schriften entnommen sind, haben wir als solche kenntlich gemacht. Die eingesetzten KI-gestützten Hilfsmittel sind mit Art und Umfang ihres Einsatzes im folgenden Abschnitt angegeben; für den Inhalt der Arbeit übernehmen wir die Verantwortung. Die Arbeit hat in gleicher oder ähnlicher Form noch keiner anderen Prüfungsbehörde vorgelegen. Die als Einzelleistung zu bewertenden Beiträge sind im Nachweis der Einzelbeiträge gekennzeichnet.

% Ort und Datum stehen gedruckt und brauchen keine Linie; unterschrieben
% wird allein auf der Linie ueber dem Namen. Liegt in thesis/unterschriften/
% eine eingescannte Unterschrift, steht sie auf der Linie; der Ordner ist
% per .gitignore vom Repository ausgenommen.
\newcommand{\unterschrift}[1]{%
	\IfFileExists{#1}{%
		\makebox[\linewidth][c]{\includegraphics[height=1.7cm,width=0.9\linewidth,keepaspectratio]{#1}}\\[-0.75cm]%
	}{}%
	\rule{\linewidth}{0.4pt}}
\vspace{1.5cm}

\noindent
\begin{minipage}[b]{0.45\textwidth}
	\AbgabeOrt, \AbgabeDatum
\end{minipage}
\hfill
\begin{minipage}[b]{0.45\textwidth}
	\unterschrift{unterschriften/unterschrift-omar.png}\\
	\AutorEins
\end{minipage}

\vspace{1.5cm}

\noindent
\begin{minipage}[b]{0.45\textwidth}
	\AbgabeOrt, \AbgabeDatum
\end{minipage}
\hfill
\begin{minipage}[b]{0.45\textwidth}
	\unterschrift{unterschriften/unterschrift-mohamad.png}\\
	\AutorZwei
\end{minipage}

% Erklaerung, KI-Angabe und Einzelbeitraege stehen je auf einer eigenen Seite.
\clearpage

\section*{Angabe der KI-gestützten Hilfsmittel}

Eingesetzt wurde das Sprachmodell Claude des Anbieters Anthropic in zwei Bereichen: bei der Softwareentwicklung über das Werkzeug Claude Code, nach selbst aufgestellten Regeln, Skills und Agents, die im Repository stehen (CLAUDE.md, Ordner .claude), und beim Schreiben der Arbeit. Die Texte haben die Verfassenden geschrieben und alle Entscheidungen selbst getroffen; das Sprachmodell hat beim Umformulieren und Korrigieren geholfen und Änderungen nach ihren Vorgaben eingearbeitet. Die Tabelle nennt je Bereich die Art des Einsatzes und das, was bei den Verfassenden lag.

\begingroup
\footnotesize
\setlength{\tabcolsep}{4pt}
\renewcommand{\arraystretch}{1.1}
\setlength{\LTcapwidth}{\textwidth}
\begin{longtable}{@{}L{2.6cm}L{6.2cm}L{5.8cm}@{}}
	\toprule
	\textbf{Bereich} & \textbf{Art des Einsatzes} & \textbf{Bei den Verfassenden} \\
	\midrule
	\endfirsthead
	\toprule
	\textbf{Bereich} & \textbf{Art des Einsatzes} & \textbf{Bei den Verfassenden} \\
	\midrule
	\endhead
	\bottomrule
	\endfoot
	Software\-entwicklung & Erzeugen und Überarbeiten von Programmcode nach selbst aufgestellten Regeln, Skills und Agents; KI-gestützte Prüfrollen vor dem Zusammenführen (Abschnitt~\ref{sec:qualitaetssicherung}) & Anforderungen, Regeln, Skills und Agents, fachliche Entscheidungen, Prüfung und Abnahme \\
	\zeilenlinie
	Schreiben der Arbeit & Umformulieren und Korrigieren der von den Verfassenden geschriebenen Texte sowie Einarbeiten von Änderungen nach deren Vorgaben & Schreiben der Texte, Gliederung, Forschungsfrage und Teilfragen, Inhalte, alle Entscheidungen, Prüfung und Freigabe der Texte \\
\end{longtable}
\endgroup

\clearpage

\section*{Nachweis der Einzelbeiträge}

Die Tabelle nennt je Teil der Arbeit, wer ihn im Schwerpunkt verfasst oder umgesetzt und wer ihn ergänzt hat.

\begingroup
\footnotesize
\setlength{\tabcolsep}{4pt}
\renewcommand{\arraystretch}{1.1}
\setlength{\LTcapwidth}{\textwidth}
\begin{longtable}{@{}L{3.6cm}L{5.5cm}L{5.5cm}@{}}
	\toprule
	\textbf{Teil} & \textbf{\AutorEins} & \textbf{\AutorZwei} \\
	\midrule
	\endfirsthead
	\toprule
	\textbf{Teil} & \textbf{\AutorEins} & \textbf{\AutorZwei} \\
	\midrule
	\endhead
	\bottomrule
	\endfoot
	Projektstart & gemeinsam: erste Struktur und erstes Grundgerüst & gemeinsam: erste Struktur und erstes Grundgerüst \\
	\zeilenlinie
	Recherche: Literaturauswertung und heuristische Analyse der Referenzplattform & Ergänzungen & Schwerpunkt \\
	\zeilenlinie
	Anforderungsanalyse: Anforderungen, User Stories, Use Cases & Schwerpunkt & Ergänzungen \\
	\zeilenlinie
	Konzeption: Architektur, Rechtemodell, Inhaltsmodell, Workflow, Datenmodell & gemeinsam & gemeinsam \\
	\zeilenlinie
	Design: Skizzen und Wireframes, Figma-Prototyp, Designsystem & gemeinsam & gemeinsam \\
	\zeilenlinie
	API-Contract & gemeinsam & gemeinsam \\
	\zeilenlinie
	Backend & Schwerpunkt & Mitarbeit \\
	\zeilenlinie
	Frontend & Mitarbeit & Schwerpunkt \\
	\zeilenlinie
	Arbeitsregeln, Skills und Agents & Schwerpunkt & Ergänzungen \\
	\zeilenlinie
	Kapitel 1 bis 5 & Verfasser & Ergänzungen \\
	\zeilenlinie
	Kapitel 6 bis 8 & Ergänzungen & Verfasser \\
	\zeilenlinie
	Heuristische Evaluation & erste Begutachtung & zweite Begutachtung \\
	\zeilenlinie
	Usability-Test & Moderation oder Protokoll im Wechsel & Testmaterial, Auswertung im Nachlauf, Moderation oder Protokoll im Wechsel \\
\end{longtable}
\endgroup

	% ---------------------------------------------------------------------------
% Zusammenfassung
% Wortgleich uebernommen aus der Confluence-Seite
% "Zusammenfassung / Abstract (Entwurf)" (Space BEH, Seiten-ID 101416961),
% Abschnitt "Zusammenfassung", Stand 04.10.2026.
% Auf Wunsch der Verfasser nur die deutsche Fassung; das englische Abstract
% der Confluence-Seite erscheint nicht in der Arbeit.
% ---------------------------------------------------------------------------
\addchap{Zusammenfassung}

An der Hochschule Bremen gibt es keinen Ort, an dem abgeschlossene studentische Projekte dauerhaft und öffentlich einsehbar sind. Die vorliegende Arbeit untersucht, welche funktionalen und gestalterischen Anforderungen ein hochschulisches E-Portfolio-Portal erfüllen muss, um solche Projekte strukturiert zu präsentieren, auffindbar zu machen und von den beteiligten Nutzergruppen akzeptiert zu werden. Aus der Literatur zu E-Portfolios, dem Technologieakzeptanzmodell und den Usability-Heuristiken nach Nielsen sowie aus einer heuristischen Analyse der Referenzplattform der HfG Schwäbisch Gmünd wurden 36 funktionale und 24 nicht-funktionale Anforderungen abgeleitet. Auf dieser Grundlage entstand ein Portal mit blockbasiertem Inhaltsmodell und einem redaktionellen Freigabeweg aus fünf Zuständen, dessen Regeln der Server durchsetzt. Geprüft wurde es mit lautem Denken in elf Sitzungen, einer heuristischen Evaluation durch zwei Begutachtende, der System Usability Scale (SUS) und einem Fragebogen nach dem Technologieakzeptanzmodell.

Der SUS-Mittelwert liegt mit 87,3 Punkten über dem vorab gesetzten Zielwert von 68, und alle elf Personen bewerten Nützlichkeit und Benutzerfreundlichkeit über dem Skalenmittelpunkt. Dem stehen 30 Usability-Probleme gegenüber, sieben davon schwerwiegend. Von den funktionalen Anforderungen sind 30 erfüllt und sechs teilweise, von den nicht-funktionalen zehn erfüllt, elf teilweise und drei nicht geprüft. Auffindbarkeit fördern verbindbare, einzeln rücknehmbare Ordnungsmerkmale und ein Zugang über die Person; Voraussetzung ist, dass Zustände und Handlungen beschriftet sichtbar sind und dasselbe überall gleich heißt. Das Inhaltsmodell lässt sich mit dem Redaktionsworkflow verbinden, indem die Bearbeitbarkeit an den Zustand des Projekts gebunden und auf dem Server durchgesetzt wird. Die Ergebnisse beruhen auf elf Personen und einer einzigen Testrunde.


	\tableofcontents
	\listoffigures
	\listoftables
	% ---------------------------------------------------------------------------
% Abkuerzungsverzeichnis
% Abgeglichen mit der Confluence-Seite "Abkuerzungsverzeichnis (Entwurf)"
% (Seiten-ID 101384193), Stand 05.10.2026. Aufgenommen ist nur, was in der
% Kurzfassung tatsaechlich vorkommt; ISO, JSON, SQL und UI sind deshalb entfallen.
% ---------------------------------------------------------------------------
\addchap{Abkürzungsverzeichnis}

\begin{description}[
	style=nextline,
	leftmargin=3.2cm,
	labelwidth=2.8cm,
	labelsep=0.4cm,
	font=\normalfont\bfseries,
	itemsep=0.2em]
	\item[A1 bis D7] Kennungen der 25 Testaufgaben in den Blöcken A bis D (Tabelle~\ref{tab:aufgaben})
	\item[API] Application Programming Interface, Programmierschnittstelle
	\item[API-Contract] Eigene, unveröffentlichte Schnittstellenbeschreibung des Projekts. Sie legt Datenmodell, Rechte, Zustandsautomat und Endpunkte fest, bevor sie umgesetzt werden; Verweise darauf nennen den Abschnitt, etwa §4.1.
	\item[BI] Behavioral Intention, Nutzungsabsicht (Konstrukt des TAM)
	\item[CSS] Cascading Style Sheets
	\item[FA] Funktionale Anforderung (FA-01 bis FA-44, Tabelle~\ref{tab:fa})
	\item[H1 bis H10] Die zehn Usability-Heuristiken nach Nielsen (2024) in den Bezeichnungen der Tabelle~\ref{tab:heuristiken}
	\item[HfG] Hochschule für Gestaltung Schwäbisch Gmünd
	\item[HSB] Hochschule Bremen
	\item[HTTP] Hypertext Transfer Protocol
	\item[ID] Kennung (Identifier)
	\item[NFA] Nicht-funktionale Anforderung (NFA-01 bis NFA-62, Tabelle~\ref{tab:nfa})
	\item[P1 bis P30] Kennungen der Usability-Probleme (Anhang~\ref{anh:problemliste})
	\item[PDF] Portable Document Format
	\item[PEOU] Perceived Ease of Use, wahrgenommene Benutzerfreundlichkeit (Konstrukt des TAM)
	\item[PU] Perceived Usefulness, wahrgenommene Nützlichkeit (Konstrukt des TAM)
	\item[REST] Representational State Transfer
	\item[SUS] System Usability Scale (Brooke, 1996)
	\item[TAM] Technology Acceptance Model, Technologieakzeptanzmodell (Davis, 1989)
	\item[TF] Teilfrage (TF 1 bis TF 3, Abschnitt~\ref{sec:ziele})
	\item[TN01 bis TN11] Pseudonyme der elf Teilnehmenden des Nutzertests
	\item[UC] Use Case, Anwendungsfall (UC-01 bis UC-05, Anhang~\ref{anh:user-stories})
	\item[US] User Story (US-01 bis US-42, Anhang~\ref{anh:user-stories})
	\item[W3C] World Wide Web Consortium
	\item[WCAG] Web Content Accessibility Guidelines (W3C, 2023)
\end{description}


	% =======================================================================
	% HAUPTTEIL
	% =======================================================================
	\clearpage
	\pagenumbering{arabic}

	% ---------------------------------------------------------------------------
% Kapitel 1: Einleitung
% Wortgleich uebernommen aus der Confluence-Seite
% "Kapitel 1: Einleitung (Entwurf)" (Space BEH, Seiten-ID 34504705),
% Stand 04.10.2026 (die vorherige LaTeX-Datei trug noch Version 7 vom
% 02.09.2026). Einzige Aenderung: Der Verweis auf den Befundstand der
% Referenzplattform zeigt auf die zusammengefuehrte Tabelle 3.1 der
% Kurzfassung statt auf die fruehere Tabelle 3.2.
% ---------------------------------------------------------------------------

\chapter{Einleitung}
\label{kap:einleitung}

\section{Problemfeld}
\label{sec:problemfeld}

An der Hochschule Bremen entstehen in jedem Semester studentische Projektarbeiten in großer Zahl. Ihr Verbleib nach der Bewertung ist ungeregelt: Ein Teil liegt auf den privaten Laufwerken der Bearbeitenden, ein anderer in einem Lernmanagementsystem, das auf den Kursbetrieb ausgelegt ist und nicht auf die Darstellung nach außen. Beide Ablagen dienen der Benotung und sind nicht darauf ausgelegt, die Arbeiten danach zu zeigen. Ein Ort, an dem abgeschlossene Projekte dauerhaft und öffentlich einsehbar wären, besteht nicht. Diese Ausgangslage beruht auf einer Rücksprache mit der Fakultät 4 (Elektrotechnik und Informatik), in der auch Zielgruppe und Anforderungen des Portals besprochen wurden (J. Egharevba, persönliche Mitteilung, Mai 2026).

Aus dieser Lage folgt ein Verlust auf mehreren Seiten. Studierenden fehlt ein Beleg der eigenen Arbeit, auf den sich in einer Bewerbung verweisen ließe. Der Hochschule entgeht eine Außendarstellung, die aus der Lehre selbst hervorginge und nicht eigens hergestellt werden müsste. Arbeitgebende und die fachlich interessierte Öffentlichkeit haben keinen niedrigschwelligen Zugang zu dem, was in den Studiengängen bearbeitet wird. Studierenden späterer Jahrgänge fehlt der Blick auf das, was vor ihnen entstanden ist, und damit ein Maßstab für die eigene Arbeit.

Lorenzo und Ittelson (2005) beschreiben ein E-Portfolio als digitalisierte Sammlung von Artefakten, die eine Einzelperson, eine Gruppe oder eine Institution repräsentiert, und sie unterscheiden studentische, lehrbezogene und institutionelle Portfolios; einschlägig ist die dritte Kategorie, denn Träger wäre die Hochschule selbst. Aus dem Anspruch, etwas zu repräsentieren, folgt für das vorliegende Vorhaben eine erste Festlegung: Eine geordnete Dateiablage genügt nicht. Ein Bestand, der die Einrichtung vertreten soll, muss gepflegt werden, bevor er dargestellt werden kann, und beides gehört in dasselbe System.

Dieselben Autoren beschreiben das institutionelle Portfolio zugleich als Zusammenführung studentischer und lehrbezogener Portfolios; ein Portal, das studentische Arbeiten sammelt und nach außen zeigt, ist danach beides. Es unterscheidet zwei Gegenstände: das einzelne Projekt, das eine Prüfung durchläuft und danach öffentlich wird, und die Profilseite einer Person, auf der die veröffentlichten Projekte zusammenlaufen und die über die Studienzeit hinweg wächst. Diese Trennung leitet Abschnitt~\ref{sec:eportfolio} aus einer eigenen Ordnung der Komponenten nach Mahasneh (2020) her.

Zur Sicht der Studierenden liegen empirische Daten vor. Mahasneh (2020) entwickelt aus einer Literaturdurchsicht und Interviews mit 20 Fachleuten ein Modell aus zehn Komponenten, das von der Biografie über Berichte und Projekte bis zu wissenschaftlichem Material und Proben studentischer Leistung reicht. 90 Studierende, die damit gearbeitet haben, beurteilen es mit einem Mittelwert von 4,58 auf einer fünfstufigen Skala positiv. Welche Bestandteile ein hochschulisches Portfolio enthalten sollte, lässt sich demnach aus vorhandener Forschung ableiten.

Eine positive Einstellung bedeutet noch nicht, dass ein System auch benutzt wird. Davis (1989) führt im Technologieakzeptanzmodell (TAM) die Akzeptanz einer Technologie auf zwei Wahrnehmungen zurück: die wahrgenommene Nützlichkeit und die wahrgenommene Benutzerfreundlichkeit. Shroff et al. (2011) haben das Modell auf E-Portfolio-Systeme übertragen. Für ein Portal, dessen Inhalte freiwillig eingestellt und vor der Veröffentlichung geprüft werden, ist Akzeptanz deshalb keine nachgelagerte Frage der Oberfläche. Bleibt das Einreichen mühsam oder sein Nutzen unklar, füllt sich der Bestand nicht, und die Sichtbarkeitslücke besteht fort.

\section{Ziele der Arbeit}
\label{sec:ziele}

Ziel der vorliegenden Arbeit ist die Konzeption und die prototypische Realisierung eines institutionellen E-Portfolio-Portals für die Hochschule Bremen, das studentische Projekte strukturiert darstellt und sie über Suche und Filter auffindbar macht. Studierende sollen einzelne Projekte in einem blockbasierten Editor aufbauen und zur Prüfung einreichen; über die Veröffentlichung entscheiden Redaktion und Administration. Daneben sollen sie eine Profilseite pflegen, auf der die veröffentlichten Projekte erscheinen, die sie verfasst oder an denen sie mitgewirkt haben. Für die Öffentlichkeit entsteht ein Zugang ohne Anmeldung. Der redaktionelle Freigabeprozess richtet sich dabei nicht gegen die Studierenden, sondern hält die Verantwortung für veröffentlichte Inhalte bei der Hochschule.

Das Portal tritt nicht an die Stelle des Lernmanagementsystems und bildet weder Kursbetrieb noch Bewertung ab; Gegenstand ist die Darstellung abgeschlossener Arbeiten. Ebenso wenig zielt die Untersuchung auf einen Marktvergleich verfügbarer E-Portfolio-Systeme, auf die Einbindung in die Verwaltungsinfrastruktur der Hochschule oder auf die organisatorische Einführung im Hochschulbetrieb. Der Bearbeitungszeitraum von neun Wochen begrenzt zusätzlich den Funktionsumfang: Umgesetzt wird, was zur Beantwortung der Teilfragen erforderlich ist, während weitergehende Funktionen dem Ausblick (Abschnitt~\ref{sec:ausblick}) vorbehalten bleiben.

Die übergeordnete Forschungsfrage dieser Arbeit lautet:

Welche funktionalen und gestalterischen Anforderungen muss ein hochschulisches E-Portfolio-Portal erfüllen, um studentische Projekte strukturiert zu präsentieren, auffindbar zu machen und von den beteiligten Nutzergruppen akzeptiert zu werden?

Diese Frage wird durch drei Teilfragen (TF) operationalisiert:

\begin{itemize}[topsep=0.4em, itemsep=0.4em]
	\item TF 1: Welche Anforderungen an Struktur, Inhalt und Workflow lassen sich aus der bestehenden Forschungsliteratur zu E-Portfolios ableiten?
	\item TF 2: Welche funktionalen und gestalterischen Merkmale fördern die Auffindbarkeit und Benutzerfreundlichkeit eines institutionellen Portfolio-Portals?
	\item TF 3: Wie lässt sich ein blockbasiertes Inhaltsmodell mit integriertem Redaktionsworkflow realisieren, das Studierenden eine flexible und medienreiche Projektpräsentation ermöglicht?
\end{itemize}

Beantwortet werden die Teilfragen auf vier Wegen: durch die Auswertung der Forschungsliteratur, die heuristische Analyse einer bestehenden Referenzplattform, die Entwicklung des Systems und eine eigene empirische Erhebung am gebauten System. Die Entwicklungsarbeit ist methodisches Mittel zur Überprüfung dessen, was aus der Literatur gewonnen wurde, und kein Selbstzweck; die Erhebung prüft anschließend, ob das Gebaute die vorab gesetzten Kriterien einlöst. TF 1 richtet sich auf vorhandenes Wissen und beugt der Gefahr vor, Anforderungen aus dem eigenen Vorverständnis zu schöpfen; wo jede Teilfrage bearbeitet und beantwortet wird, zeigt Abschnitt~\ref{sec:aufbau}.

Damit die Antworten überprüfbar bleiben, sind die Maßstäbe vorab festgelegt und nicht erst aus dem Ergebnis gewonnen. Die Akzeptanz misst sich an den beiden Konstrukten des TAM, der wahrgenommenen Nützlichkeit und der wahrgenommenen Benutzerfreundlichkeit (Davis, 1989). Erhoben werden sie mit den zwölf Items der beiden zugehörigen Skalen, ergänzt um die zwei Items zur Nutzungsabsicht, mit denen Davis (1989) in seiner zweiten Untersuchung die künftige Nutzung erfragt. Da Davis für diese Skalen keine Normwerte berichtet, dient als Orientierungspunkt der Mittelpunkt der siebenstufigen Antwortskala, also 4,00: Werte darüber gelten als eher zustimmend, Werte darunter als eher ablehnend, ein Wert von genau 4,00 als unentschieden, und zwar für alle drei Skalen. Dieser Punkt ist eine Festlegung der vorliegenden Arbeit und kein Literaturwert (Abschnitt~\ref{sec:testkonzept}). Den inhaltlichen Zuschnitt prüft das Komponentenmodell von Mahasneh (2020). Für die Gebrauchstauglichkeit stehen zwei Maßstäbe nebeneinander: das Urteil Begutachtender anhand der zehn Usability-Heuristiken (Nielsen, 2024) und der Summenwert der System Usability Scale (SUS) nach Brooke (1996), für den ein Zielwert von 68 Punkten gilt. Dieser Wert liegt auf der Notenskala aus 241 industriellen Untersuchungen (Sauro \& Lewis, 2016, zitiert nach Lewis \& Sauro, 2018) in der Mitte der Note C und entspricht damit dem Median der Studienmittelwerte. Die Festlegung vorab begrenzt den Spielraum, einen Maßstab nachträglich dem Ergebnis anzupassen; wo eine Anwendungsregel erst nach der Messung nötig wird, weist Abschnitt~\ref{sec:bewertungsrahmen} sie als solche aus.

\section{Lösungsansatz}
\label{sec:loesungsansatz}

Der Weg zur Beantwortung beginnt bei der Literatur und endet bei der Messung am gebauten System. Ein Teil der Anforderungen an Struktur und Freigabe der Inhalte leitet sich aus der Forschungsliteratur ab: Lorenzo und Ittelson (2005) liefern die institutionelle Perspektive, Mahasneh (2020) den Blick auf die Bestandteile; die übrigen stammen aus dem Anwendungsfall. Den zweiten Weg bildet die heuristische Analyse einer bestehenden Referenzplattform, des öffentlich zugänglichen Portfolio-Portals der Hochschule für Gestaltung Schwäbisch Gmünd, nach dem Verfahren von Nielsen und Molich (1990) und mit der heute kanonischen Liste der zehn Heuristiken (Nielsen, 2024) als Raster.

Untersucht wurde allein die Sicht ohne Anmeldung, weil nur der öffentliche Bereich zugänglich war. Nicht erfüllt ist allein die siebte Heuristik, Flexibilität und Effizienz der Nutzung: Ein Suchfeld fehlt ganz, wer Studiengang und Schlagwort zugleich einschränken möchte, findet dafür keinen Weg, und nach dem Semester lässt sich gar nicht auswählen. Teilweise erfüllt ist unter anderem die Sichtbarkeit des Systemzustands, weil eine Brotkrumennavigation fehlt und allein Seitenkopf und Bedienelement zeigen, welche Einschränkung wirkt; den vollständigen Befundstand führt Tabelle~\ref{tab:heuristiken} in Abschnitt~\ref{sec:verwandte-arbeiten}.

Aus diesen Befunden entsteht eine eigene Anforderung an die Auffindbarkeit, die über die Referenzplattform hinausgeht. Sie umfasst ein Suchfeld über Titel, Kurzbeschreibung, Personennamen und Schlagworte, kombinierbare Facettenfilter, eine sichtbare Anzeige der aktiven Auswahl und die Möglichkeit, eine Einschränkung einzeln oder die ganze Auswahl in einem Schritt zurückzunehmen. Übernommen werden dagegen das einheitliche Aufbauschema der Projektseiten und die Erschließung über Studiengang und Person; das Semester, das die Referenzplattform nur am Projekt nennt, wird im eigenen System zur auswählbaren Achse.

Bevor gebaut wird, entsteht der Entwurf zunächst als Handskizzen und Wireframes in Miro, anschließend als Prototyp in Figma, sodass Gestaltungsentscheidungen vor der ersten Codezeile verhandelbar werden. In der Qualitätssicherung der Realisierung kommen zusätzlich sprachmodellgestützte Prüfschritte zum Einsatz (Abschnitt~\ref{sec:qualitaetssicherung}). Das Portal unterscheidet vier Akteurgruppen: Studierende, Redaktion, Administration und externe Nutzende. Das unterschriebene Exposé nennt die zweite Gruppe Betreuende beziehungsweise Lehrende; die vorliegende Arbeit benennt sie nach ihrer Aufgabe im Portal und nicht nach einer Stellung im Lehrbetrieb. Drei Gruppen sind an ein Konto gebunden; die externen Nutzenden haben keines und sind zugleich die Hauptzielgruppe. Diese Gewichtung leitet die vorliegende Arbeit aus dem äußeren der beiden Zwecke ab, die Lorenzo und Ittelson (2005) dem institutionellen Portfolio zuschreiben (Abschnitt~\ref{sec:eportfolio}); sie deckt sich mit der Einschätzung der Fakultät 4 (Abschnitt~\ref{sec:methodisches-vorgehen}).

Die Projektinhalte entstehen in einem blockbasierten Editor, der eine Seite in einzeln anordenbare Bausteine zerlegt, darunter Text, Bilder, Galerien, Video und Audio. Der redaktionelle Ablauf kennt fünf Zustände. Der reguläre Weg führt vom Entwurf über die Einreichung und die Genehmigung zur Veröffentlichung; den fünften Zustand, die Überarbeitung, erreicht ein Projekt nur, wenn Redaktion oder Administration es zurückgeben. Bearbeiten lässt sich ein Projekt allein im Entwurf und in der Überarbeitung, damit ein geprüfter Stand nicht unbemerkt verändert wird.

Das fertige System wird von den Verfassern inspiziert und mit elf Personen erprobt, neun Studierenden und zwei Externen; auch die Aufgaben der Redaktion und der Administration bearbeiten Studierende. Die Erhebung verbindet vier Methoden, von denen jede eine Frage beantwortet, welche die übrigen offenlassen. Beim lauten Denken, für das Nielsen (1994b) die methodische Referenz ist, bearbeiten die Teilnehmenden Aufgaben entlang der Kernpfade des Portals und zeigen dabei Bedienprobleme und ihre Ursachen. Die heuristische Evaluation, von den beiden Verfassern vor dem Nutzertest getrennt vorgenommen, erfasst auch Bereiche, zu denen keine Aufgabe führt. Die SUS misst die wahrgenommene Gebrauchstauglichkeit (Brooke, 1996), der Fragebogen zum TAM die wahrgenommene Nützlichkeit, die wahrgenommene Benutzerfreundlichkeit und die Nutzungsabsicht; Shroff et al. (2011) haben das Modell für E-Portfolio-Systeme mit einem eigenen, aus früheren TAM-Skalen abgeleiteten Fragebogen geprüft.

Lautes Denken, Heuristiken und SUS liegen nach der Annahme der vorliegenden Arbeit auf der Seite des Aufwands: Sie zeigen vor allem, wie mühsam die Benutzung ist. Ob ihr Ergebnis den Beteiligten etwas wert ist, zeigt erst die Skala zur wahrgenommenen Nützlichkeit, denn ein Portal kann leicht zu bedienen und zugleich entbehrlich sein. Beide Fragebögen werden deshalb getrennt ausgewertet und nicht ineinander umgerechnet.

\section{Aufbau der Arbeit}
\label{sec:aufbau}

Kapitel~\ref{kap:anforderungsanalyse} bestimmt Akteurgruppen und Anforderungen mit ihrer Herkunft; es greift dabei auf Quellen vor, die Kapitel~\ref{kap:grundlagen} im Einzelnen auswertet. Kapitel~\ref{kap:grundlagen} legt die Grundlagen aus der E-Portfolio-, der Akzeptanz- und der Usability-Literatur und berichtet die Analyse der Referenzplattform. Beide Kapitel bearbeiten TF 1, Abschnitt~\ref{sec:usability} liefert zudem die Kriterien für TF 2. Kapitel~\ref{kap:konzeption} entwirft das System; die Abschnitte~\ref{sec:inhaltsmodell} und~\ref{sec:workflow} tragen TF 3, Abschnitt~\ref{sec:ui-design} den gestalterischen Teil von TF 2. Kapitel~\ref{kap:realisierung} erprobt den Entwurf an der Realisierung. Kapitel~\ref{kap:test} berichtet Methodik und Daten des Tests ohne Wertung; Kapitel~\ref{kap:evaluation} bewertet sie, beantwortet in Abschnitt~\ref{sec:bewertung-test} TF 2 und gleicht beide Anforderungskataloge mit dem gebauten Stand ab. Abschnitt~\ref{sec:gesamtzusammenfassung} beantwortet TF 1 und TF 3, verdichtet die Antwort auf TF 2, führt die Forschungsfrage zusammen und benennt die Grenzen der Arbeit; Abschnitt~\ref{sec:ausblick} gibt den Ausblick.

	% ---------------------------------------------------------------------------
% Kapitel 2: Anforderungsanalyse
% Kurzfassung vom 04.10.2026, neu geschrieben aus der Langfassung
% (thesis/langfassung/kapitel02.tex, Confluence-Seite 34701325).
% Ziel: Umfang etwa halbiert, Kapitel 1 bis 8 auf hoechstens 50 Seiten.
% Kapitel 1 bis 7 schliessen ohne eigene Zusammenfassung (Gliederungsseite).
% ---------------------------------------------------------------------------

\chapter{Anforderungsanalyse}
\label{kap:anforderungsanalyse}

Die 60 Anforderungen an das Portal stützen sich auf vier Quellen von ungleicher Beweiskraft, und 22 davon sind Festlegungen aus dem Zuschnitt des Vorhabens ohne externen Beleg. Jede Anforderung wird deshalb mit ihrer Herkunft geführt. Diese Angabe entscheidet, ob eine später erfüllte Anforderung etwas über die Forschungsliteratur aussagt oder allein über den Anwendungsfall Hochschule Bremen; sie trägt damit die erste Teilfrage.

\section{Stakeholder und Zielgruppen}
\label{sec:stakeholder}

Fakultät, Verwaltung und Prüfende prägen Anforderungen, treten im System aber nicht auf. Das Portal bedienen vier Akteurgruppen: Studierende, Redaktion, Administration und externe Nutzende. Zu den externen Nutzenden zählen Arbeitgebende, die fachlich interessierte Öffentlichkeit und Studierende späterer Jahrgänge, solange sie nicht angemeldet sind. Sie sind die Hauptzielgruppe und zugleich die einzige Gruppe ohne Konto. Da eine Rolle an ein Konto gebunden ist, bildet der öffentliche Zugang keine vierte Rolle, sondern eine Sichtbarkeitsebene, die aus dem Zustand „nicht angemeldet“ folgt.

Die drei kontogebundenen Gruppen sind ineinander enthalten. Studierende legen eigene Projekte an, reichen sie ein, ziehen sie zurück und pflegen ihre Profilseite. Die Redaktion darf all das und prüft zusätzlich die ihr zugewiesenen Einreichungen; die Administration darf, was die Redaktion darf, und verwaltet Konten, Stammdaten und die Auswahl der Startseite. Der Grund liegt in der Registrierung: Eine Person hat genau ein Konto, und wer hochgestuft wird, soll das eigene Portfolio nicht verlieren. Abbildung~\ref{abb:akteurgruppen} zeigt diese Schachtelung und die eine Regel, die sich ihr entzieht: Über die eigene Arbeit entscheidet niemand, auch die Administration nicht. Abschaltbar ist diese Sperre allein für den Betrieb mit nur einer verwaltenden Person (Abschnitt~\ref{sec:rechtemodell}).

\abbildung{abb-2-1-akteure.png}{0.8}{abb:akteurgruppen}%
{Die vier Akteurgruppen, ihre Schachtelung und die Grenze der Anmeldung}%
{Die vier Akteurgruppen, ihre Schachtelung und die Grenze der Anmeldung (eigene Darstellung)}

Die Ordnung, in die ein Projekt eingehängt wird, stammt aus der Hochschule selbst. Erhoben wurden die fünf Fakultäten mit ihren 75 Studiengängen (Hochschule Bremen, 2026a, 2026b, 2026c, 2026d, 2026e) und, soweit veröffentlicht, deren Modullisten. Fakultäten und Studiengänge werden zu Filterachsen, Module ordnen ein Projekt nur ein.

Zwischen den Gruppen bestehen Interessengegensätze, die sich nur ausgestalten lassen. Studierende möchten Arbeiten rasch zeigen, die Hochschule steht für Veröffentlichtes ein und verlangt eine Prüfung; das Portal antwortet mit einem mehrstufigen Zustandsmodell statt einer Direktveröffentlichung. Externe Nutzende möchten möglichst viel sehen, Studierende die Kontrolle über Unfertiges behalten; daraus folgt die scharfe Trennung zwischen veröffentlichten und internen Zuständen.

\section{Methodisches Vorgehen}
\label{sec:methodisches-vorgehen}

Die Anforderungen wurden quellengestützt hergeleitet: Jede erhält eine Kennung und eine Herkunft, sodass sie bis zu ihrem Beleg zurückverfolgbar bleibt. Die erste Quelle ist die Forschungsliteratur. Lorenzo und Ittelson (2005) liefern die Einordnung als institutionelles Portfolio, Mahasneh (2020) den inhaltlichen Zuschnitt, Davis (1989) und Shroff et al. (2011) die Begründung, warum Gebrauchstauglichkeit unter den Qualitätsanforderungen vorn steht (Kapitel~\ref{kap:grundlagen}).

Die zweite Quelle ist eigene Vorarbeit: die heuristische Analyse der Referenzplattform der Hochschule für Gestaltung Schwäbisch Gmünd nach Nielsen und Molich (1990) mit den zehn Heuristiken nach Nielsen (2024) als Raster (Abschnitt~\ref{sec:verwandte-arbeiten}). Jedes untersuchte Merkmal wurde einer von drei Klassen zugeordnet: übernehmen, gezielt verbessern oder eigenständig konzipieren. Übernommen werden das einheitliche Aufbauschema der Projektseiten und die Erschließung über Studiengang und Person. Verbessert wird die Auffindbarkeit, weil dort jede Suche fehlt und sich die Filter nicht verbinden lassen. Einreichung, Prüfung und Verwaltung sind eigenständig konzipiert, denn zugänglich war nur die Sicht ohne Anmeldung.

Die dritte Quelle ist eine Rücksprache mit der Fakultät 4 im Mai 2026 (J. Egharevba, persönliche Mitteilung, Mai 2026), aus der die Einschätzung stammt, dass der überwiegende Teil der Zielgruppe außerhalb der Hochschule liegt. Die vierte ist die erhobene Ordnung aus Fakultäten, Studiengängen und Modulen. In den Katalogen erscheinen diese drei als „Referenzanalyse“, „Fakultät“ und „Modulhandbücher“; „Kontext“ steht für Festlegungen ohne externen Beleg. Für die erste Teilfrage zählen allein Anforderungen mit der Herkunft Lorenzo und Ittelson (2005) oder Mahasneh (2020).

Das Vorgehen hat Grenzen. Vor der Entwicklung wurde keine Befragung künftiger Nutzender durchgeführt, und die Referenzanalyse stammt von zwei Begutachtenden statt der drei bis fünf, die Nielsen und Molich (1990) empfehlen. Die Anforderungen sind damit begründet, aber erst in den Kapiteln~\ref{kap:test} und~\ref{kap:evaluation} geprüft. Eine förmliche Priorisierung fehlt; begrenzt hat den Umfang faktisch der Bearbeitungszeitraum von neun Wochen.

\section{Funktionale Anforderungen}
\label{sec:funktionale-anforderungen}

Tabelle~\ref{tab:fa} führt 36 funktionale Anforderungen (FA) in Blöcken nach Zugang: FA-01 bis FA-06 rollenübergreifend, FA-10 bis FA-18 ohne Anmeldung, FA-20 bis FA-30 für Studierende, FA-31 bis FA-37 für die Redaktion und FA-40 bis FA-44 für die Administration. Ein Block nennt, wer eine Fähigkeit zuerst braucht, nicht, wer sie allein hat.

Der Zugang beruht auf zwei Adressen (FA-01): Eine private dient dauerhaft als Kennung, eine Hochschuladresse belegt einmalig die Zugehörigkeit. Ein Portfolio, das nach Lorenzo und Ittelson (2005) eine Person repräsentiert, soll auch nach der Exmatrikulation vorzeigbar bleiben, während Hochschulkonten mit ihr auslaufen. Die Rolle folgt aus der Endung der bestätigten Hochschuladresse (FA-02), weil sie das einzige maschinell prüfbare Merkmal der Registrierung ist.

Für die Hauptzielgruppe entscheidet die Auffindbarkeit über den Wert des Portals, und dort lag der einzige Negativbefund der Referenzanalyse. Die Anforderungen gehen deshalb über die Vorlage hinaus: verbindbare Filtergruppen mit Mehrfachauswahl (FA-12), ein Suchfeld für den Einstieg ohne Kenntnis der Systematik (FA-14) und ein sichtbarer, einzeln rücknehmbarer und als Adresse teilbarer Auswahlzustand (FA-15). Neben den Weg über das Projekt tritt einer über die Person: Aus der Biografie als einziger personenbezogener Komponente bei Mahasneh (2020) folgt die Profilseite (FA-16), auf der veröffentlichte Projekte ohne Zutun zusammenlaufen (FA-29).

Im redaktionellen Bereich sind Genehmigung (FA-32) und Veröffentlichung (FA-33) getrennt, damit geprüfte Arbeiten etwa gebündelt zum Semesterbeginn freigeschaltet werden können. Eine Rückgabe trägt eine Begründung (FA-34), die Verständigung läuft über Antwort, interne Notiz und Änderungsaufforderung (FA-36), und für die Freigabe gilt ein Vier-Augen-Prinzip, das an der Beteiligung hängt und nicht an der Rolle (FA-35). Der Zugriff der Administration auf fremde Inhalte soll lesend sein (FA-42), weil ein stiller Eingriff in einen fremden Text einem System widerspräche, das Autorschaft ausweist.

\begingroup
\footnotesize
\setlength{\tabcolsep}{4pt}
\renewcommand{\arraystretch}{1.05}
\setlength{\LTcapwidth}{\textwidth}
\begin{longtable}{@{}L{1.2cm}L{9.6cm}L{3.7cm}@{}}
	\caption[Funktionale Anforderungen mit Kennung und Herkunft]{Funktionale Anforderungen mit Kennung und Herkunft (eigene Darstellung)}
	\label{tab:fa}\\
	\toprule
	\textbf{ID} & \textbf{Anforderung} & \textbf{Herkunft} \\
	\midrule
	\endfirsthead
	\toprule
	\textbf{ID} & \textbf{Anforderung} & \textbf{Herkunft} \\
	\midrule
	\endhead
	\bottomrule
	\endfoot
	FA-01 & Selbstregistrierung mit privater Kennung und bestätigter Hochschuladresse & Lorenzo und Ittelson (2005) \\
	FA-02 & Rolle aus der bestätigten Hochschuladresse, Endungen konfigurierbar & Kontext \\
	FA-03 & Anmeldung, Passwortzurücksetzung, erzwungener Wechsel eines Initialpassworts & Kontext \\
	FA-04 & Zugriffssteuerung serverseitig auf Rollen- und Ressourcenebene & Kontext \\
	FA-05 & Reaktion angemeldeter Personen auf ein Projekt, höchstens eine je Person & Referenzanalyse \\
	FA-06 & Rückkanal für Rückfragen, auch ohne Anmeldung & Kontext \\
	FA-10 & Veröffentlichte Projekte ohne Anmeldung lesen & Lorenzo und Ittelson (2005) \\
	FA-11 & Projektdetailansicht mit allen Bausteinen und Angaben & Mahasneh (2020) \\
	FA-12 & Verbindbare Filtergruppen mit Mehrfachauswahl: Fakultät, Studiengang, Abschluss, Semester, Projekttyp, Kategorie & Referenzanalyse \\
	FA-14 & Suchfeld über Titel, Kurzbeschreibung, Personennamen und Schlagworte & Referenzanalyse \\
	FA-15 & Auswahlzustand sichtbar, einzeln und gesammelt rücknehmbar, als Adresse teilbar & Referenzanalyse \\
	FA-16 & Profilseite ohne Anmeldung mit Biografie und Studiengang & Mahasneh (2020); Modulhandbücher \\
	FA-17 & Kuratierte Auswahl veröffentlichter Projekte auf der Startseite & Fakultät \\
	FA-18 & Übersicht der Studiengänge und Module als Einordnungskontext & Modulhandbücher \\
	FA-20 & Blockbasierter Editor mit Bausteinen für Text und multimediale Inhalte & Mahasneh (2020) \\
	FA-21 & Projekt als Entwurf anlegen, zwischenspeichern, bearbeiten, löschen & Kontext \\
	FA-22 & Beschreibende Angaben pflegen, zuständige Redaktion wählen & Kontext \\
	FA-23 & Projekt zur redaktionellen Prüfung einreichen & Lorenzo und Ittelson (2005) \\
	FA-24 & Einreichung bis zur Entscheidung zurückziehen & Kontext \\
	FA-25 & Eigene Projekte mit ihrem Zustand überblicken & Nielsen (2024) \\
	FA-26 & Änderungsaufforderungen lesen, als erledigt kennzeichnen, erneut einreichen & Kontext \\
	FA-28 & Profilseite mit Biografie, Kontaktangaben und Bild pflegen & Mahasneh (2020) \\
	FA-29 & Veröffentlichte Projekte der Person erscheinen ohne Zutun auf der Profilseite & Mahasneh (2020) \\
	FA-30 & Gruppenarbeit: Mitwirkende mit Funktion aufnehmen, gemeinsam bearbeiten & Fakultät \\
	FA-31 & Arbeitsliste der zugewiesenen Einreichungen & Fakultät \\
	FA-32 & Projekt genehmigen & Lorenzo und Ittelson (2005) \\
	FA-33 & Genehmigte Arbeit in einem zweiten Schritt veröffentlichen & Fakultät \\
	FA-34 & Projekt mit Begründung zur Überarbeitung zurückgeben & Kontext \\
	FA-35 & Keine Freigabe der eigenen Arbeit, unabhängig von der Rolle & Fakultät \\
	FA-36 & Antwort, interne Notiz, Änderungsaufforderung mit Erledigungsstand & Kontext \\
	FA-37 & Veröffentlichtes Projekt für die Startseite vorschlagen & Fakultät \\
	FA-40 & Konten anlegen, Rollen ändern, deaktivieren, Passwörter zurücksetzen & Kontext \\
	FA-41 & Zustandsunabhängige Gesamtübersicht aller Projekte & Kontext \\
	FA-42 & Lesender Zugriff auf fremde Inhalte, Eingriff nur über Zustandsentscheidungen & Kontext \\
	FA-43 & Stammdaten pflegen: Kategorien, Fakultäten, Studiengänge, Module, Semester & Modulhandbücher \\
	FA-44 & Auswahl der Startseite kuratieren, mit unveränderlichem Protokoll & Fakultät \\
\end{longtable}
\endgroup

\section{Nicht-funktionale Anforderungen}
\label{sec:nichtfunktionale-anforderungen}

Ein Portal, dessen Inhalte freiwillig eingestellt werden, füllt sich nur, wenn das Einstellen als lohnend und mühelos wahrgenommen wird. Davis (1989) führt Nutzung auf genau diese beiden Wahrnehmungen zurück, und bei Shroff et al. (2011) ist der Pfad von der wahrgenommenen Benutzerfreundlichkeit zur wahrgenommenen Nützlichkeit mit 0,71 der stärkste abgesicherte Zusammenhang. Gebrauchstauglichkeit ist deshalb keine Verfeinerung, sondern Bedingung dafür, dass der Nutzen wahrgenommen wird.

Die nicht-funktionalen Anforderungen (NFA, Tabelle~\ref{tab:nfa}) stehen in Zehnerblöcken und tragen, wo möglich, eine messbare Schwelle mit vorab benanntem Instrument. NFA-01 bis NFA-05 prüfen lautes Denken und heuristische Evaluation (Nielsen, 2024). Für NFA-06 gilt ein Zielwert von 68 Punkten im Fragebogen nach Brooke (1996), der auf der Notenskala aus 241 Untersuchungen (Sauro \& Lewis, 2016, zitiert nach Lewis \& Sauro, 2018) in der Mitte der Note C liegt. NFA-10 und NFA-11 erhebt der Fragebogen zu den Konstrukten nach Davis (1989), die Barrierefreiheit messen die Erfolgskriterien des W3C (2023). Der Schutz nicht veröffentlichter Inhalte (NFA-61) ist die Bedingung dafür, dass Studierende überhaupt Unfertiges einstellen. Für Antwortzeiten und Tragfähigkeit (NFA-30 bis NFA-32) liegt kein Messverfahren vor; sie gelten für den Prototypbetrieb.

\begingroup
\footnotesize
\setlength{\tabcolsep}{4pt}
\renewcommand{\arraystretch}{1.05}
\setlength{\LTcapwidth}{\textwidth}
\begin{longtable}{@{}L{1.4cm}L{9.6cm}L{3.5cm}@{}}
	\caption[Nicht-funktionale Anforderungen mit Kennung und Herkunft]{Nicht-funktionale Anforderungen mit Kennung und Herkunft (eigene Darstellung)}
	\label{tab:nfa}\\
	\toprule
	\textbf{ID} & \textbf{Anforderung} & \textbf{Herkunft} \\
	\midrule
	\endfirsthead
	\toprule
	\textbf{ID} & \textbf{Anforderung} & \textbf{Herkunft} \\
	\midrule
	\endhead
	\bottomrule
	\endfoot
	NFA-01 & Erlernbarkeit: Editor ohne Schulung in wenigen Minuten bedienbar & Nielsen (1994b) \\
	NFA-02 & Einheitliche Bedienmuster, Begriffe und visuelle Sprache & Nielsen (2024) \\
	NFA-03 & Fehlervermeidung und Rückholbarkeit bei folgenreichen Aktionen & Nielsen (2024) \\
	NFA-04 & Zustand und Rückmeldung auf Systemaktionen jederzeit erkennbar & Nielsen (2024) \\
	NFA-05 & Wiedererkennen statt Erinnern & Nielsen (2024) \\
	NFA-06 & Mindestens 68 Punkte im SUS-Bogen & Brooke (1996), Lewis und Sauro (2018) \\
	NFA-10 & Erkennbarer Nutzen: eigene Arbeit wird extern sichtbar & Davis (1989) \\
	NFA-11 & Geringer wahrgenommener Aufwand bei Erstellung und Einreichung & Shroff et al. (2011) \\
	NFA-20 & Textkontrast mindestens 4,5:1, Bedienelemente mindestens 3:1 & W3C (2023) \\
	NFA-21 & Keine allein farbliche Codierung & W3C (2023) \\
	NFA-22 & Vollständige Tastaturbedienung mit sichtbarem Fokus & W3C (2023) \\
	NFA-23 & Textvergrößerung und Umbruch ohne Verlust & W3C (2023) \\
	NFA-24 & Mindestgröße interaktiver Ziele & W3C (2023) \\
	NFA-25 & Semantische Auszeichnung für Hilfsmittel & W3C (2023) \\
	NFA-26 & Dunkler Farbmodus mit denselben Kontrastschwellen & Kontext \\
	NFA-30 & Öffentliche Seiten laden unter zwei Sekunden & Kontext \\
	NFA-31 & Rückmeldung im Editor ohne spürbare Verzögerung & Kontext \\
	NFA-32 & Tragfähigkeit bei mehreren hundert Projekten & Kontext \\
	NFA-40 & Nutzbarkeit auf Rechner, Tablet und Telefon & Kontext \\
	NFA-41 & Korrekte Darstellung in aktuellen Browsern & Kontext \\
	NFA-50 & Farbvorgabe des Corporate Designs & Fakultät \\
	NFA-60 & Rechte serverseitig durchgesetzt & Kontext \\
	NFA-61 & Nicht veröffentlichte Inhalte und interne Notizen unzugänglich & Kontext \\
	NFA-62 & Sichere Ablage der Passwörter, sensible Felder verlassen den Server nicht & Kontext \\
\end{longtable}
\endgroup

\section{User Stories und Use Cases}
\label{sec:user-stories}

Ein Katalog sagt, was vorhanden sein muss, aber nicht, ob die Teile zusammenpassen. Die Anforderungen wurden deshalb zusätzlich als 30 User Stories und fünf mehrstufige Use Cases formuliert (Anhang~\ref{anh:user-stories}). Zwei Abläufe tragen das Portal. Im redaktionellen Weg legt eine studierende Person ein Projekt an, wählt die zuständige Redaktion und reicht ein (UC-01); die Redaktion genehmigt und veröffentlicht es oder gibt es mit begründeter Änderungsaufforderung zurück (UC-02), worauf die einreichende Seite die Punkte erledigt und erneut einreicht (UC-03). Die Prüfung ist damit eine Schleife und kein einmaliges Tor, und die Anforderungen an Rückmeldung und Nachverfolgbarkeit folgen daraus. Im Weg der Hauptzielgruppe grenzt eine externe Person die Projektübersicht über mehrere Filtergruppen ein, öffnet ein Projekt (UC-04) und gelangt von dort auf die Profilseite einer beteiligten Person (UC-05). Kein Schritt dieses Weges verlangt eine Anmeldung, und daran misst sich, ob das Portal seinen institutionellen Zweck erfüllt.

	% ---------------------------------------------------------------------------
% Kapitel 3: Grundlagen und verwandte Arbeiten
% Kurzfassung vom 04.10.2026, neu geschrieben aus der Langfassung
% (thesis/langfassung/kapitel03.tex, Confluence-Seite 34373645).
% Tabelle 3.1 fuehrt die frueheren Tabellen 3.1 und 3.2 zusammen.
% ---------------------------------------------------------------------------

\chapter{Grundlagen und verwandte Arbeiten}
\label{kap:grundlagen}

Drei Stränge der Literatur tragen diese Arbeit. Die E-Portfolio-Literatur sagt, was ein Portfolio ist und woraus es besteht, die Akzeptanzforschung, warum ein freiwillig gefülltes System an seiner Bedienbarkeit scheitern kann, und die Usability-Literatur, woran sich Bedienbarkeit bemisst. Diese Arbeitsteilung bestimmt, welche Anforderung sich auf einen Beleg stützen kann und welche eine eigene Ableitung bleibt. Das Kapitel bearbeitet die erste Teilfrage und liefert in Abschnitt~\ref{sec:usability} die Kriterien der zweiten.

\section{E-Portfolio: Begriffe und Kategorien}
\label{sec:eportfolio}

Lorenzo und Ittelson (2005) bestimmen ein E-Portfolio als digitalisierte Sammlung von Artefakten, die eine Einzelperson, eine Gruppe oder eine Institution repräsentiert. Es könne zugleich als Verwaltungswerkzeug dienen, das die Arbeiten ordnet und steuert, wer sie sehen darf. Sichtbarkeitssteuerung ist damit im Begriff angelegt und keine Zutat des Baus. Die Autoren unterscheiden studentische, lehrbezogene und institutionelle Portfolios und halten fest, dass sich die ersten beiden auf der breitesten Ebene zu einem institutionellen zusammenführen lassen. Die vorliegende Arbeit liest die dritte Kategorie deshalb als Aggregationsebene: Ein Portal, das studentische Arbeiten sammelt und zeigt, ist Ort ihrer Entstehung und zugleich Aggregat, das die Einrichtung vertritt.

Für das institutionelle Portfolio geben Lorenzo und Ittelson eine Bestimmung wieder, nach der es eine fokussierte Auswahl authentischer Arbeiten ist, die Rechenschaft der Einrichtung belegt (Kahn, 2001, zitiert nach Lorenzo \& Ittelson, 2005). Daraus folgen zwei Festlegungen. Eine fokussierte Auswahl verlangt einen Auswahlschritt, also eine redaktionelle Freigabe (Abschnitt~\ref{sec:workflow}). Und weil die Autoren die Leserschaft veröffentlichter Unterlagen einer Einrichtung als offen beschreiben, legt die vorliegende Arbeit den Zugang ohne Anmeldung fest; die Übertragung auf studentische Arbeiten ist ihr eigener Schritt. Ebenso eine Setzung ist die Einordnung als Präsentationsportfolio, das fertige Ergebnisse zeigt und keine prozessbegleitenden Artefakte.

Was ein hochschulisches E-Portfolio enthalten sollte, beantwortet Mahasneh (2020). Aus einer Literaturdurchsicht und Interviews mit 20 Fachleuten entsteht ein Modell aus zehn Komponenten: Biografie, Kursplan, Berichte und Forschungsarbeiten, Hausaufgaben, Projekte und Experimente, Aktivitäten, Zusammenfassungen, wissenschaftliches Material, Audio- und Videoclips sowie Leistungsproben. 90 Studierende bewerteten es mit einem Mittelwert von 4,58 auf einer fünfstufigen Skala; die Aussage, ein E-Portfolio helfe beim Veröffentlichen eigener Erfahrungen, erreicht 4,96 bei geringer Streuung. Die vorliegende Arbeit ordnet die Komponenten nach ihrem Bezugspunkt, was die Quelle nicht tut: Personenbezogen ist allein die Biografie, an eine Lehrveranstaltung gebunden sind Kursplan, Hausaufgaben und Aktivitäten, die übrigen sechs sind artefaktbezogen. Daraus folgt die Trennung von Profilseite und Projekt, und die drei veranstaltungsgebundenen Komponenten entfallen mit der Abgrenzung gegenüber dem Lernmanagementsystem.

Beide Quellen haben Grenzen. Lorenzo und Ittelson (2005) legen einen Übersichtsbericht mit Momentaufnahmen von 2005 vor, Mahasneh (2020) erhebt ohne Kontrollgruppe an einem einzigen College in Jordanien. Das Modell trägt deshalb als Inhaltsraster, nicht als Wirkungsnachweis. Aus einem Whitepaper übernehmen Lorenzo und Ittelson zudem die Frage, ob ein System benutzerfreundlich genug für eine leichte Übernahme sei (ePortConsortium, 2003, zitiert nach Lorenzo \& Ittelson, 2005); die E-Portfolio-Literatur stellt die Akzeptanzfrage also, beantwortet sie aber nicht.

\section{Technologieakzeptanzmodell (TAM)}
\label{sec:tam}

Davis (1989) erklärt die Annahme eines Systems mit zwei Größen. Die wahrgenommene Nützlichkeit (PU) ist der Grad, zu dem eine Person glaubt, dass das System ihre Arbeitsleistung steigert; die wahrgenommene Benutzerfreundlichkeit (PEOU) der Grad, zu dem sie die Benutzung für frei von Anstrengung hält. In zwei Untersuchungen mit 152 Personen erreichen die beiden Skalen aus je sechs Aussagen Reliabilitäten von 0,98 und 0,94. Die Nützlichkeit hängt stärker mit der Nutzung zusammen als die Benutzerfreundlichkeit, etwa mit 0,85 gegenüber 0,59 für die vorhergesagte Nutzung. Wird die Nützlichkeit kontrolliert, verschwindet der Beitrag der Benutzerfreundlichkeit fast vollständig. Davis liest das vorläufig als Kette: Benutzerfreundlichkeit wirkt über die Nützlichkeit auf die Nutzung. Normwerte berichtet er nicht.

\abbildung{abb-3-1-tam.png}{0.85}{abb:tam}%
{Die von Davis (1989) berichteten Zusammenhänge und die von Shroff et al. (2011) geprüften Pfade}%
{Die von Davis (1989) berichteten Zusammenhänge und die von Shroff et al. (2011) geprüften Pfade (eigene Darstellung nach Davis, 1989, und Shroff et al., 2011)}

Shroff et al. (2011) übertragen das Modell auf ein E-Portfolio-System, das 72 Studierende ein Semester lang genutzt hatten. Abbildung~\ref{abb:tam} stellt beide Untersuchungen nebeneinander: Abgesichert sind bei Shroff et al. nur die beiden Pfade, die von der Benutzerfreundlichkeit ausgehen, darunter der zur Nützlichkeit mit 0,71 als stärkster. Der gemeinsame Nenner ist schmal, liegt aber genau an der Stelle, auf die es hier ankommt. Stichprobe und Einrichtung sind klein, die Wirkungsrichtung setzen beide Modelle voraus, statt sie zu prüfen.

Für die vorliegende Arbeit folgt daraus eine Festlegung vor jeder Evaluation: Ist die Benutzerfreundlichkeit der Hebel auf die wahrgenommene Nützlichkeit, entscheidet sie darüber, ob sich ein freiwillig gefülltes Portal füllt, und ist eine Frage des Entwurfs, nicht der Nacharbeit. Das Modell liefert zudem das Messinstrument für Kapitel~\ref{kap:test}; welche Merkmale Benutzerfreundlichkeit herstellen, lässt Davis offen.

\section{Usability und Gestaltungsprinzipien}
\label{sec:usability}

Gebrauchstauglichkeit ist nach Brooke (1996) nur in Bezug auf einen Nutzungskontext bestimmbar; die Kriterien dieses Abschnitts gelten daher für ein institutionelles Portfolio-Portal. Das Verfahren, eine Oberfläche ohne Testpersonen zu beurteilen, geht auf Nielsen und Molich (1990) zurück: Mehrere Begutachtende prüfen unabhängig gegen eine kleine Liste von Prinzipien. In vier Versuchsreihen fand eine Person im Mittel 20 bis 51 Prozent der bekannten Probleme, fünf Personen rund zwei Drittel; daraus stammt die Empfehlung von drei bis fünf Begutachtenden und die Ergänzung durch nutzerzentrierte Verfahren wie das laute Denken (Nielsen, 1994b). Die heute geführte Liste der zehn Heuristiken geht auf eine Faktorenanalyse von 249 Problemen zurück (Nielsen, 1994a, 2024).

Die Heuristiken sind Faustregeln und keine Abnahmekriterien (Nielsen, 2024). Tabelle~\ref{tab:heuristiken} leitet je Heuristik ein Kriterium für ein Portfolio-Portal ab; diese Spalte ist eine Ableitung der vorliegenden Arbeit und der Maßstab für Abschnitt~\ref{sec:bewertung-test}. Die rechte Spalte enthält bereits das Ergebnis der Referenzanalyse aus Abschnitt~\ref{sec:verwandte-arbeiten}.

\begingroup
\footnotesize
\setlength{\tabcolsep}{4pt}
\renewcommand{\arraystretch}{1.05}
\setlength{\LTcapwidth}{\textwidth}
\begin{longtable}{@{}L{0.9cm}L{3.6cm}L{6.6cm}L{3.3cm}@{}}
	\caption[Die zehn Heuristiken, das abgeleitete Kriterium und der Befund der Referenzplattform]{Die zehn Heuristiken, das abgeleitete Kriterium und der Befund der Referenzplattform (eigene Darstellung nach Nielsen, 2024; heuristische Analyse Mai 2026)}
	\label{tab:heuristiken}\\
	\toprule
	\textbf{Nr.} & \textbf{Heuristik} & \textbf{Abgeleitetes Kriterium} & \textbf{Referenzplattform} \\
	\midrule
	\endfirsthead
	\toprule
	\textbf{Nr.} & \textbf{Heuristik} & \textbf{Abgeleitetes Kriterium} & \textbf{Referenzplattform} \\
	\midrule
	\endhead
	\bottomrule
	\endfoot
	H1 & Sichtbarkeit des Systemzustands & Zustand eines Vorgangs und wirkende Einschränkung sind ablesbar & teilweise: keine Brotkrume, Auswahl nur im Bedienelement \\
	H2 & Übereinstimmung mit der realen Welt & Benennungen folgen dem Vokabular der Hochschule & erfüllt \\
	H3 & Kontrolle und Freiheit & Einschränkung einzeln und vollständig rücknehmbar & teilweise: kein Zurücksetzen \\
	H4 & Konsistenz und Standards & Gleichartiges folgt demselben Aufbau & erfüllt \\
	H5 & Fehlervermeidung & Folgenreiche Schritte werden verhindert, bevor sie erklärt werden müssen & nicht beurteilbar \\
	H6 & Wiedererkennen statt Erinnern & Ordnungsmerkmale werden sichtbar angeboten & teilweise: keine Schlagwortübersicht \\
	H7 & Flexibilität und Effizienz & Arbeit ohne Weg über alle Ordnungsebenen erreichbar & nicht erfüllt: kein Suchfeld, Filter nicht kombinierbar \\
	H8 & Ästhetisches, minimalistisches Design & Gestaltung ordnet sich den Arbeiten unter & erfüllt \\
	H9 & Hilfe beim Beheben von Fehlern & Zurückgewiesene Eingabe nennt Grund und nächsten Schritt & nicht beurteilbar \\
	H10 & Hilfe und Dokumentation & Erklärung eines Vorgangs im System auffindbar & teilweise: keine Hilfeseite \\
\end{longtable}
\endgroup

Die Auffindbarkeit, nach der die zweite Teilfrage fragt, führen die Heuristiken nicht als eigenes Thema, und auch der SUS-Bogen erfasst sie nicht (Lewis \& Sauro, 2018). Die vorliegende Arbeit setzt sie deshalb aus drei Kriterien zusammen: Die Systematik des Bestandes liegt als auswählbare Merkmale vor (H6), jede Einschränkung ist einzeln und vollständig rücknehmbar (H3), und die wirkende Einschränkung bleibt erkennbar, weil sich ein leeres Ergebnis sonst nicht von einem leeren Bestand unterscheiden lässt (H1).

Daneben treten die Web Content Accessibility Guidelines 2.2 (W3C, 2023) als Standard mit teils rechnerisch prüfbaren Schwellen. Einschlägig sind für das Portal zwölf Erfolgskriterien, darunter Textkontrast von 4,5:1 (1.4.3), 3:1 für Bedienelemente (1.4.11), keine allein farbliche Information (1.4.1), Umbruch bei 320 CSS-Pixeln (1.4.10), sichtbarer Fokus (2.4.7), Zielgröße von 24 mal 24 CSS-Pixeln (2.5.8) sowie Tastaturbedienung, Struktur und Name, Rolle und Zustand (2.1.1, 1.3.1, 4.1.2). Zielniveau ist AA.

Als drittes Instrument dient die System Usability Scale (SUS) nach Brooke (1996): zehn abwechselnd gepolte Aussagen, die einen Wert von 0 bis 100 ergeben. Brooke nennt keinen Vergleichswert; diesen liefert eine Notenskala aus 241 Untersuchungen, auf der 68 die Mitte der Note C und den Median markiert (Sauro \& Lewis, 2016, zitiert nach Lewis \& Sauro, 2018). Ein unbequemer Befund gehört an diese Stelle: Davis (1989) fand, dass Flexibilität nicht immer mit Benutzerfreundlichkeit einhergeht, was ein Werkzeug mit gestalterischer Freiheit unmittelbar betrifft.

\section{Verwandte Arbeiten und Referenzplattformen}
\label{sec:verwandte-arbeiten}

Schon für ein Inhaltsmodell des studentischen E-Portfolios sieht Mahasneh (2020) keine Vorarbeit; für hochschulische Portale gilt das nach Einschätzung der vorliegenden Arbeit erst recht. Die Literatur ordnet dagegen die Wege zu einem solchen System: Lorenzo und Ittelson (2005) unterscheiden Eigenentwicklung, quelloffene Systeme, Einkauf und gängige Werkzeuge. Die Eigenentwicklung liefert genau das Gewünschte, verlangt aber Zeit, Fachkompetenz und dauerhaft gebundenes Personal, worauf Abschnitt~\ref{sec:ausblick} zurückkommt. Freie Werkzeuge erlauben freiere Portfolios, setzen aber Kenntnisse im Erstellen von Webseiten voraus, und Shroff et al. (2011) nennen bestehende Systeme starr in der Gestaltung ihrer Inhalte. Gestalterische Freiheit ohne solche Kenntnisse ist damit die Antwort auf ein benanntes Problem.

\abbildung{abb-3-2-hfg-filter.png}{0.8}{abb:hfg}%
{Studiengangsübersicht der Referenzplattform}%
{Studiengangsübersicht der Referenzplattform: eine Auswahlliste als vollständiger Filterapparat (Bildschirmaufnahme portfolio.hfg-gmuend.de, abgerufen am 05.09.2026)}

Als Referenz dient das Portfolio-Portal der Hochschule für Gestaltung Schwäbisch Gmünd (Hochschule für Gestaltung Schwäbisch Gmünd, o.\,J.), das Arbeiten aus vier Bachelor-Studiengängen ohne Anmeldung zeigt. Gewählt wurde es auf Empfehlung des Erstprüfers, ohne systematisches Auswahlverfahren; die Befunde stützen keine Aussage über Portfolio-Portale allgemein. Die beiden Verfasser begutachteten im Mai 2026 fünf Klassen von Seiten nach Nielsen und Molich (1990) und bewerteten vierstufig, mit „nicht beurteilbar“ für alles hinter der Anmeldung. Tabelle~\ref{tab:heuristiken} fasst das Ergebnis zusammen: drei Heuristiken erfüllt, vier teilweise, zwei nicht beurteilbar, H7 nicht erfüllt.

Abbildung~\ref{abb:hfg} zeigt den Kern des Negativbefunds. Unter dem Seitentitel steht eine einzige Auswahlliste, und sie ist der gesamte Filterapparat; ein Suchfeld fehlt, und welche Einschränkung wirkt, sagt allein das Bedienelement. Semester und Modul erscheinen an jedem Projekt, sind aber keine auswählbaren Achsen. Erschlossen ist der Bestand über Studiengang, Schlagwort und Person, wobei die Person eine Sammelseite ohne Biografie bleibt. Daraus folgen die eigenen Anforderungen an Suche, verbindbare Filter und eine Profilseite (Abschnitt~\ref{sec:funktionale-anforderungen}).

Die Reichweite ist begrenzt: zwei statt drei bis fünf Begutachtende, eine stichprobenartige Begehung, keine Messung von Ladezeiten oder Kontrasten. Da die Außensicht zugleich die Sicht der Hauptzielgruppe ist, mindert die fehlende Anmeldung die Aussage zur öffentlichen Auffindbarkeit nicht; für den Einreichungsweg fehlt dagegen jede Vorlage.

	% ---------------------------------------------------------------------------
% Kapitel 4: Konzeption
% Kurzfassung vom 04.10.2026, neu geschrieben aus der Langfassung
% (thesis/langfassung/kapitel04.tex, Confluence-Seite 51740676).
% Entfallen: Abbildungen Rechtemodell, Editor und Kommentarbereich (Figma),
% Handskizze, Wireframe Freischaltung und Wireframe Startseite.
% ---------------------------------------------------------------------------

\chapter{Konzeption}
\label{kap:konzeption}

Zwischen einer Anforderung und einer Zeile Programmtext liegt eine Entscheidung, und sie fällt hier. Die Abschnitte~\ref{sec:inhaltsmodell} und~\ref{sec:workflow} tragen die dritte Teilfrage, Abschnitt~\ref{sec:ui-design} den gestalterischen Teil der zweiten. Folgt eine Festlegung aus einer Anforderung, steht deren Kennung dabei.

\section{Systemarchitektur}
\label{sec:architektur}

Die Forderung, Rechte auf dem Server durchzusetzen (NFA-60) und nicht veröffentlichte Inhalte zu schützen (NFA-61), bestimmt den Schnitt: eine Oberfläche im Browser ohne eigene Wahrheit über Sichtbarkeit und ein Anwendungsserver, der jede Anfrage gegen dieselbe Datenhaltung prüft. Beide sprechen über eine zustandslose Schnittstelle nach dem Muster von REST. Diese Schnittstelle wird vor ihrer Umsetzung in einer schriftlichen Beschreibung festgelegt, im Folgenden API-Contract genannt: Entitäten, Rechtematrix, Zustandsautomat und Endpunkte. Der API-Contract ist eine eigene, unveröffentlichte Projektdokumentation; Verweise darauf nennen seinen Abschnitt, etwa §4.1, und er erscheint deshalb nicht im Literaturverzeichnis. Zwei Personen bauen beide Seiten gleichzeitig; ohne vorab entschiedene Beschreibung entstünde die Schnittstelle zweimal.

\abbildung{abb-4-1-architektur.png}{0.85}{abb:architektur}%
{Schnitt des Systems, seine fachlichen Bereiche und die beiden austauschbaren Fremddienste}%
{Schnitt des Systems, seine fachlichen Bereiche und die beiden austauschbaren Fremddienste (eigene Darstellung)}

Abbildung~\ref{abb:architektur} zeigt den Aufbau. Der Server ist nach Fachlichkeit geschnitten, und zwei Festlegungen gelten für alle Bereiche: ein einheitliches Fehlerformat mit maschinenlesbarem Code, über den die Oberfläche entscheidet (NFA-02), und versionierte, unveränderliche Schemamigrationen. Dateiablage und Nachrichtenversand hängen als austauschbare Dienste außen an. In keiner Antwort erscheint eine Speicheradresse, weil die Sichtbarkeit einer Datei am Zustand ihres Projekts hängt, den eine öffentliche Adresse nicht kennt (FA-04, NFA-61).

\section{Benutzerrollen und Rechtemodell}
\label{sec:rechtemodell}

Das Rechtemodell setzt die Schachtelung aus Abschnitt~\ref{sec:stakeholder} um. Drei Rollen hängen an einem Konto: \texttt{STUDENT}, \texttt{EDITOR} und \texttt{ADMIN}. Für externe Nutzende gibt es keinen Datensatz; geprüft wird nur, ob eine Anmeldung vorliegt, sonst gelten die öffentlichen Leserechte (FA-10, FA-16). Die Rolle folgt aus der Endung der bestätigten Hochschuladresse (FA-02); eine Rollenauswahl im Formular gäbe es sonst jeder Person frei. Die verwaltende Rolle entsteht nie durch Registrierung, sondern durch Hochstufung (FA-40).

Autorisiert wird auf zwei Ebenen (FA-04): Die Rollenebene prüft, ob eine Rolle eine Handlung überhaupt ausführen darf, die Ressourcenebene, ob diese Person bei diesem Projekt handeln darf, und wertet Autorschaft, Mitwirkung und zugewiesene Redaktion aus. Rolle und Kontozustand werden bei jedem Zugriff aus der Datenbank gelesen, damit eine Deaktivierung sofort wirkt. Quer zu den Rollen liegt die Freigabesperre (FA-35): Über die eigene Arbeit entscheidet niemand, auch die Administration nicht. Abschaltbar bleibt sie, weil bei einer einzigen verwaltenden Person sonst niemand deren Projekt freigeben könnte. Der Zugriff der Administration auf fremde Inhalte ist im Grundsatz lesend (FA-41, FA-42); ausgenommen sind Befugnisse, die der API-Contract ihr in jedem Zustand einräumt, darunter das Löschen von Projekten und Dateien und das Zusammenführen doppelter Kategorien. Abschnitt~\ref{sec:pruefung-fa} prüft diese Abweichung.

Die Fehlerpolitik unterscheidet fehlende Anmeldung, fehlende Berechtigung und nicht vorhandenen Gegenstand und maskiert keinen Fall. Getragen wird das, weil die Kennungen nicht erratbar sind (Abschnitt~\ref{sec:datenmodell}).

\section{Blockbasiertes Inhaltsmodell}
\label{sec:inhaltsmodell}

Ein Projekt besteht aus beschreibenden Angaben und einer geordneten Liste von Bausteinen. Zu den Angaben gehören Titel, Kurzbeschreibung, Projekttyp, Fakultät, Studiengang, Modul, Abschluss, Kategorie, Semester und bis zu fünf Schlagwörter; sie bilden zugleich die Merkmale der Suche. Im Entwurf dürfen sie fehlen, beim Einreichen werden die Pflichtangaben geprüft (FA-21, FA-22), denn sonst könnte niemand schreiben, bevor alle Rahmendaten feststehen. Das Semester ist eine Stammdatenzeile und kein freier Text, weil drei Schreibweisen desselben Semesters drei Filterwerte wären.

Das Inhaltsmodell ist an das Dokumentmodell des Editors gebunden: Ein Baustein entspricht einem Knoten der obersten Ebene. Sein Inhalt wird unverändert als Dokument gespeichert, seine Art zusätzlich als eigene Spalte, weil sich Integritätsregeln, Filter und Auswertungen über eine Spalte wartbar formulieren lassen. Eine Schemafassung je Baustein hält ihn bei späteren Änderungen überführbar. Fünfzehn Bausteinarten stehen bereit, von Überschrift, Text, Liste und Zitat über Quelltext, Formel, Tabelle und Trennlinie bis zu Bild, Galerie, Video, Audio, Dokument, Einbettung und Inhaltsverzeichnis (FA-20). Die Bandbreite folgt dem Komponentenmodell von Mahasneh (2020); Aufgabenlisten fehlen bewusst, weil ein Präsentationsportfolio fertige Ergebnisse zeigt.

\abbildung{abb-4-3-blockmodell.png}{0.85}{abb:blockmodell}%
{Aufbau eines Projekts: Angaben, Bausteinliste und der Verweis auf getrennt abgelegte Dateien}%
{Aufbau eines Projekts: Angaben, Bausteinliste und der Verweis auf getrennt abgelegte Dateien (eigene Darstellung)}

Abbildung~\ref{abb:blockmodell} zeigt die zweite tragende Grenze: Hochgeladene Dateien liegen nie im Dokument, sondern sind eine eigene Entität, auf die ein Baustein nur über eine Kennung verweist. Sonst steckte jede Datei codiert in jeder Antwort. Erlaubte Dateitypen und Größen prüft der Server, Einbettungen sind auf geprüfte Anbieter beschränkt, und ausgeliefert wird jede Datei über einen eigenen Endpunkt, weil ihre Sichtbarkeit am Projektzustand hängt.

An einem Projekt arbeiten mehrere Personen (FA-30). Ein Versionsfeld meldet beim Speichern einen Konflikt, statt fremde Arbeit still zu überschreiben, und eine Bearbeitungssperre mit regelmäßigem Lebenszeichen legt die Reihenfolge vorher fest. Das automatische Zwischenspeichern folgt der Folgerung aus Abschnitt~\ref{sec:tam}: Verlorene Arbeit beschädigt die Bereitschaft zur weiteren Nutzung (NFA-11).

\section{Workflow- und Statussystem}
\label{sec:workflow}

Ein Projekt durchläuft fünf Zustände, im Datenmodell \texttt{DRAFT}, \texttt{SUBMITTED}, \texttt{APPROVED}, \texttt{PUBLISHED} und \texttt{REVISION}, in der Oberfläche Entwurf, Eingereicht, Genehmigt, Veröffentlicht und Überarbeitung. Dass überhaupt eine Prüfung vor der Veröffentlichung liegt, folgt aus der fokussierten Auswahl nach Kahn (2001, zitiert nach Lorenzo \& Ittelson, 2005). Genehmigung (FA-32) und Veröffentlichung (FA-33) sind getrennt. Die Menge zulässiger Übergänge ist abschließend, jeder andere Wechsel wird serverseitig abgewiesen, und bearbeiten lässt sich der Inhalt nur im Entwurf und in der Überarbeitung, damit eine Prüfung nicht nachträglich wertlos wird.

\abbildung{abb-4-4-zustandsautomat.png}{0.85}{abb:zustandsautomat}%
{Zustandsautomat des Projekts mit den zulässigen Übergängen und der handelnden Gruppe}%
{Zustandsautomat des Projekts mit den zulässigen Übergängen und der handelnden Gruppe (eigene Darstellung nach API-Contract, §4.1)}

Die Ausgestaltung folgt dem Akzeptanzrahmen: Wird der Freigabeweg als Hürde erlebt, sinkt die Bereitschaft, das Portfolio zu pflegen (Shroff et al., 2011). Der Entwurf verkürzt deshalb den wiederholt durchlaufenen Pfad. Eine Einreichung lässt sich bis zur Entscheidung zurückziehen (FA-24), und aus der Überarbeitung wird erneut eingereicht, ohne dass ein Projekt neu entsteht (FA-26). Abbildung~\ref{abb:zustandsautomat} zeigt, dass die Überarbeitung kein Endpunkt ist, sondern aus drei Zuständen erreichbar, auch aus der Veröffentlichung: Eine Veröffentlichung ist widerrufbar, ohne dass ein Projekt verloren geht. Gibt die zugewiesene Redaktion ihre Zuweisung ab, fällt ein eingereichtes oder zurückgegebenes Projekt in den Entwurf zurück, damit es nicht ohne Zuständige liegen bleibt.

Der Zustand ist an jeder Darstellung eines Projekts sichtbar (H1, NFA-04), damit die einreichende Seite erkennt, ob die Arbeit bei ihr oder bei der Redaktion liegt. Die Zustandshistorie sehen nur Beteiligte, Redaktion und Administration, denn wie oft eine Arbeit zurückgewiesen wurde, gehört nicht in ihre öffentliche Darstellung. Die redaktionelle Verständigung läuft über drei Kommentararten (FA-36); eine Änderungsaufforderung löst den Wechsel in die Überarbeitung aus. Eine öffentliche Kommentarfunktion entfällt, weil sie eine Moderation verlangte. Die Startseite zeigt eine kuratierte Auswahl von zehn Projekten: Die Redaktion schlägt vor (FA-37), die Administration entscheidet, und ein unveränderliches Protokoll hält jede Änderung fest (FA-44).

\section{Design des User Interfaces (Figma-Prototyp)}
\label{sec:ui-design}

Der Prototyp zeigt den Entwurfsstand, nicht das gebaute System. Seine Entscheidungen folgen den drei Klassen der Referenzanalyse; die gezielte Verbesserung betrifft fast ausschließlich die Auffindbarkeit. Sechs Merkmalsgruppen erschließen den Bestand: Fakultät, Studiengang, Abschluss, Semester, Projekttyp und Kategorie (FA-12). Werte innerhalb einer Gruppe verknüpfen sich als Oder, verschiedene Gruppen als Und, und gewählte Fakultäten schränken die angebotenen Studiengänge ein, damit keine zwangsläufig leere Kombination entsteht.

\abbildung{abb-4-5-projektuebersicht.png}{0.55}{abb:facettenfilter}%
{Facettenfilter der Projektübersicht mit Trefferzahl und Zurücksetzen}%
{Facettenfilter der Projektübersicht mit Trefferzahl und Zurücksetzen (eigene Darstellung, Ausschnitt aus dem Figma-Prototyp)}

Abbildung~\ref{abb:facettenfilter} zeigt, wie der Entwurf die drei Kriterien der Auffindbarkeit adressiert. Die Gruppen stehen dauerhaft in der Seitenleiste und sind selbst die Übersicht über die Ordnung des Bestandes (H6). Jede Einschränkung ist einzeln entfernbar, daneben steht ein vollständiges Zurücksetzen (H3, FA-15). Eine Ziffer am Gruppenkopf und die Trefferzahl über der Liste sollen die wirkende Einschränkung ablesbar machen (H1); ob das genügt, prüft Abschnitt~\ref{sec:bewertung-test}. Der Suchzustand steht in der Adresse und ist damit weitergebbar, etwa in einer Bewerbung. Als zweiter Zugang tritt die Profilseite daneben (FA-16), die nicht beantwortet, welche Arbeiten es zu einem Thema gibt, sondern was eine Person gemacht hat.

Der Editor zeigt das Projekt so, wie es später erscheint; ein Formular mit festen Feldern wäre eine engere Beschreibung desselben Gegenstands. Die Werkzeugleiste hängt am Baustein, Medienarten führen zu einer Einfügehilfe für Hochladen oder geprüfte Einbettung, und der Kommentarbereich der Prüfung steht neben dem Projekt, weil eine Rückfrage zu ihrem Gegenstand gehört. Das Farbkonzept kennt eine Markenfarbe, das Blau der Hochschule (NFA-50). Eine konsequent angewandte visuelle Identität stützt die Reputation einer Einrichtung (Van den Bosch et al., 2005, zitiert nach McKenney, 2020), und Blau gehört unter US-amerikanischen Hochschulen zu den häufigsten Primärfarben (McKenney, 2020). Fakultätsfarben entfallen, weil Information nicht allein über Farbe vermittelt werden darf (W3C, 2023). Die Gestaltungswerte sind zweistufig abgelegt, Rohwerte unten und Rollenbezeichnungen oben, was zugleich den dunklen Modus trägt (NFA-26). Interaktive Grundkomponenten werden bezogen und nicht selbst gebaut, denn Tastaturbedienung und Rollenauszeichnung sind Verhalten und nicht Farbe (NFA-22, NFA-25).

Vor dem Prototyp standen Handskizzen und Wireframes in Miro, und am stärksten verändert hat sich der Zugang. Die früheste Skizze nennt die Rollen noch „Student“ und „Betreuer“, also nach Personen im Lehrbetrieb statt nach Aufgaben im Portal. Abbildung~\ref{abb:wireframes} zeigt den nächsten Stand mit getrennten Registrierungswegen, auf denen die Rolle gewählt wird; Lehrende landeten danach in einem Wartezustand bis zur Freischaltung. Der spätere Entwurf leitet die Rolle aus der bestätigten Adresse ab und streicht beides (FA-02), denn eine gewählte Rolle ist eine Angabe, eine abgeleitete ein Nachweis. Bestätigt hat der Weg dagegen den öffentlichen Einstieg mit kuratierter Auswahl.

\abbildung{abb-4-9-wireframes-registrierung.png}{0.6}{abb:wireframes}%
{Wireframes der getrennten Registrierungswege für Studierende und Lehrende}%
{Wireframes der getrennten Registrierungswege für Studierende und Lehrende (eigene Darstellung, Wireframe Mai 2026)}

\section{Datenbankmodell}
\label{sec:datenmodell}

Das relationale Modell führt Personen mit ihrem Portfolio, Projekte mit Bausteinen, Kommentaren, Reaktionen und Zustandshistorie, die kuratierte Auswahl mit Protokoll, Dateien, Bestätigungscodes und Anfragen des Rückkanals. Daneben stehen die Stammdaten Fakultät, Studiengang, Modul, Kategorie, Semester und Schlagwort. Abbildung~\ref{abb:datenmodell} zeigt die zentrale Stellung des Projekts: Ordnungsdaten und Personen zeigen darauf, Bausteine, Kommentare und Historie gehen von ihm aus. Nur weil Fakultät und Studiengang eigene Zeilen mit festem Bezug sind, lässt sich die Studiengangsauswahl einschränken (FA-12). Zugewiesene Redaktion und Mitwirkende sind getrennt geführt; Mitwirkende tragen eine frei benannte Funktion und keine Systemrolle (FA-30, FA-31).

\abbildung{abb-4-12-datenmodell.png}{0.9}{abb:datenmodell}%
{Entitäten des Datenmodells und ihre Beziehungen}%
{Entitäten des Datenmodells und ihre Beziehungen (eigene Darstellung)}

Fachliche Invarianten sichert die Datenbank selbst. Primärschlüssel sind zufällige Kennungen, weil sie in öffentlichen Adressen erscheinen und fortlaufende Zahlen ein Durchprobieren erlaubten. Die Reihenfolge der Bausteine sichert eine Eindeutigkeitsbedingung mit Prüfung am Ende der Transaktion, damit Umsortieren gelingt, und Schlagwörter sind ohne Rücksicht auf die Schreibweise eindeutig. Zwei Abwesenheiten sind ebenso Entscheidungen: Die Projektliste eines Portfolios wird bei jedem Zugriff ermittelt statt gespeichert (FA-29), und das Löschen eines Kontos ist eine Deaktivierung, weil fremde Projekte an der Person hängen.

	% ---------------------------------------------------------------------------
% Kapitel 5: Realisierung
% Kurzfassung vom 04.10.2026, neu geschrieben aus der Langfassung
% (thesis/langfassung/kapitel05.tex, Confluence-Seite 51642372).
% Entfallen: Abbildungen Blocktabelle, Einfuegeleiste, Projektseite,
% Projektliste und Administration; Abschnitt 5.2.6 ist in 5.2.4 aufgegangen.
% ---------------------------------------------------------------------------

\chapter{Realisierung}
\label{kap:realisierung}

Ob sich ein blockbasiertes Inhaltsmodell mit integriertem Redaktionsworkflow realisieren lässt, entscheidet sich daran, ob der Entwurf einer Umsetzung standhält. Die Realisierung dient deshalb als Werkzeug zur Prüfung der Annahmen aus Kapitel~\ref{kap:konzeption}. Dargestellt ist der eingefrorene Stand, auf dem auch der Test lief: 22 Endpunktklassen, 28 Schemamigrationen und 35 Routen der Oberfläche.

\section{Wahl der Realisierungsplattform}
\label{sec:plattform}

Drei Forderungen des Entwurfs bestimmen die Wahl: serverseitig erzwingbare Regeln, ein unverändert ablegbares Dokumentmodell mit abfragbarer Projektion und austauschbare Fremddienste. Das Backend läuft auf Java 21 mit Spring Boot 3.5, weil Steuerung, Sicherheit, Persistenz und Eingabeprüfung aus einem Ökosystem kommen. Die Rollenebene der Autorisierung setzen Pfadregeln der Sicherheitskonfiguration durch, die Ressourcenebene bleibt eigener Programmtext, weil Autorschaft, Mitwirkung und Redaktionszuweisung sich nicht als Pfadregel ausdrücken lassen. Angemeldet wird mit signierten Token ohne Sitzung auf dem Server; der bekannte Preis ist, dass sich ein ausgegebenes Token nicht zurückziehen lässt.

PostgreSQL 16 löst den scheinbaren Widerspruch aus Abschnitt~\ref{sec:datenmodell}: relationale Integrität für die Struktur und der Spaltentyp \texttt{jsonb} für die Knoten des Editors. Schemaänderungen laufen über versionierte Migrationen, deren Ergebnis beim Start gegen die Entitäten geprüft wird. Die Oberfläche ist eine Einseitenanwendung aus React 19, TypeScript und Vite; ihre Antworttypen werden aus der Schnittstellenbeschreibung des Servers erzeugt. Die Grundkomponenten stammen aus shadcn/ui auf der Verhaltensbibliothek Base UI, der Editor beruht auf Tiptap 3 und damit auf ProseMirror, an dessen Dokumentmodell das Inhaltsmodell gebunden ist. Quer zu beiden Seiten gilt: Jede Fehlerantwort trägt denselben Umschlag mit Fehlercode, und fehlt ein Pflichtwert der Konfiguration, startet die Anwendung nicht.

\section{Ausgewählte Realisierungsaspekte}
\label{sec:realisierungsaspekte}

\subsection{Der Baustein in der Datenbank}
\label{subsec:baustein-db}

Ein Baustein ist eine Zeile, die den Knoten unverändert als \texttt{jsonb} und seine Art als eigene Spalte trägt. Was in Abschnitt~\ref{sec:inhaltsmodell} Behauptung war, erweist sich hier als Bedingung: Über die Spalte lässt sich die Aufzählung der fünfzehn zulässigen Arten als Datenbankbedingung formulieren, über das Dokument nur als schwer wartbarer Ausdruck. Die Reihenfolge sichert eine Eindeutigkeitsbedingung über Projekt und Sortierindex mit aufgeschobener Prüfung, sodass der Tausch zweier Bausteine ohne Platzhalterwerte gelingt. Die von der Art abhängigen Regeln, etwa Pflichtattribute und die Liste zulässiger Einbettungsanbieter, liegen in einem eigenen Dienst. Die Schemafassung wird beim Lesen überführt, sodass ein alter Bestand ohne Migration weiterläuft.

\subsection{Autosave, Versionsfeld und Bearbeitungssperre}
\label{subsec:autosave}

Das automatische Speichern ersetzt die Bausteinliste in einem Aufruf als Verschmelzung: Bekannte Kennungen werden aktualisiert, unbekannte angelegt, fehlende gelöscht, und die Antwort liefert die neuen Kennungen zurück, sonst legte der nächste Takt dieselben Bausteine doppelt an. Der Takt ist entprellt und nach oben auf zehn Sekunden begrenzt, denn eine reine Entprellung speichert bei ununterbrochenem Tippen nie. Weicht die mitgeschickte Fassung ab, meldet der Server einen Konflikt, der nie selbsttätig wiederholt wird, weil ein zweiter Versuch genau die fremde Änderung überschriebe. Zwei Befunde wurden erst im Betrieb sichtbar: Inhalt und Kopfangaben erhöhen dieselbe Fassungsnummer und laufen deshalb nacheinander, und ein Takt ohne Änderung schreibt nichts mehr, nachdem er zuvor in jeder anderen Sitzung grundlos Konflikte erzeugt hatte. Die Bearbeitungssperre liegt im Arbeitsspeicher, weil sie Sekunden lebt; ihr Preis, dass sie nur je Instanz gilt, ist im API-Contract festgehalten und wird vom Versionsfeld aufgefangen.

\abbildung{abb-5-2-editor.jpg}{1.0}{abb:editor}%
{Editor mit Werkzeugleiste, Bausteinen und dem Kommentarbereich der Prüfung}%
{Editor mit Werkzeugleiste, Bausteinen und dem Kommentarbereich der Prüfung (Bildschirmaufnahme der laufenden Anwendung, Ausschnitt)}

\subsection{Der Editor als zweite Hälfte des Inhaltsmodells}
\label{subsec:editor}

Abbildung~\ref{abb:editor} zeigt einen Ausschnitt des gebauten Arbeitsplatzes: oben die Werkzeugleiste, darunter ein Bildbaustein, Titel und Kurzbeschreibung sowie das Inhaltsverzeichnis als eigener Baustein; rechts steht der Zustand Überarbeitung, daneben der Kommentarbereich mit zwei offenen Änderungsaufforderungen und einer Antwort. Zwischen Editor und Speicherung genügt eine schmale Zuordnung von Knoten zu Bausteinzeilen. Was das Ökosystem des Editors mitbringt, wird bezogen statt nachgebaut, darunter der Ziehgriff, der Platzhalter und der Vorschlagsmechanismus hinter dem Schrägstrichbefehl; am getesteten Stand öffnet der getippte Schrägstrich das Menü allerdings nicht (P30, Abschnitt~\ref{sec:durchfuehrung}). Eigene Erweiterungen entstehen nur, wo Fachwissen dahintersteht, etwa für Medien, Formel und Inhaltsverzeichnis. Öffentliche Detailseite, Prüfansicht und Druckansicht rendern gespeicherte Bausteine mit demselben schreibgeschützten Renderer, sodass dieselbe Struktur nicht dreimal getrennt gebaut werden muss.

\subsection{Der Workflow als serverseitige Schranke}
\label{subsec:workflow-schranke}

Der Freigabeweg liegt als Schranke im Schreibpfad: Inhalt ist nur im Entwurf und in der Überarbeitung bearbeitbar, jeder andere Zustand ergibt einen Konflikt. Dieselbe Schranke gilt für Hochladen und Löschen von Dateien, sonst entstünde an einem veröffentlichten Projekt am Freigabeweg vorbei eine öffentlich abrufbare Datei; ausgenommen ist allein die Administration. Zwei Schranken hat erst die Umsetzung verbindlich gemacht: Eingereicht wird nur mit zugewiesener Redaktion, weil ein Projekt sonst in keiner Arbeitsliste erscheint, und erneut eingereicht erst, wenn keine Änderungsaufforderung mehr offen ist.

\abbildung{abb-5-5-pruefansicht.jpg}{1.0}{abb:pruefansicht}%
{Prüfansicht mit den drei Kommentararten, dem Statusmenü und der Vorschau}%
{Prüfansicht mit den drei Kommentararten, dem Statusmenü und der Schaltfläche zur Vorschau (Bildschirmaufnahme der laufenden Anwendung, Ausschnitt)}

Abbildung~\ref{abb:pruefansicht} zeigt die Prüfansicht. Der Zustand steht als Schaltfläche über dem Kommentarbereich und ist zugleich der Ort, an dem der Übergang ausgelöst wird. Diese Bündelung erfüllt das Kriterium zur Sichtbarkeit des Systemzustands, hat sich im Test aber als schwer auffindbar erwiesen (P6, Abschnitt~\ref{sec:bewertung-test}). Umgekehrt lief die Umsetzung an einer Stelle dem Entwurf voraus: Die Hilfeseite hatte der Prototyp nicht vorgesehen, und ihre Vorlage wurde nachträglich ergänzt.

\subsection{Vom Designsystem zur Tokenschicht}
\label{subsec:tokenschicht}

Die zweistufige Tokenschicht aus Abschnitt~\ref{sec:ui-design} ist technisch erzwungen: Die mitgelieferten Farbpaletten von Tailwind CSS sind abgeschaltet, Klassen für Rohfarben existieren nicht. Die Sperre hat Nebenwirkungen, die kein Fehler und kein Test meldet. Sie entfernt auch den Wert „durchsichtig“, sodass bezogene Komponenten sichtbare Rahmen erhielten; sechs Radien kollidierten mit einem gleichnamigen Namensraum und verwiesen auf sich selbst; und ein Fokusring in der Markenfarbe war auf einer Schaltfläche derselben Farbe unsichtbar, bis ein Versatz in Hintergrundfarbe ergänzt wurde (W3C, 2023). Gefunden wurden alle drei erst durch den Abgleich jeder Klasse gegen das gebaute Stylesheet: Eine Regel, die eine andere erzwingt, verlangt eine eigene Prüfung am Erzeugnis.

\section{Qualitätssicherung}
\label{sec:qualitaetssicherung}

Eine sicherheitsrelevante Zusage des API-Contracts gilt erst, wenn sie ausführbar geprüft ist. Tests belegen etwa, dass der Passworthash in keiner Antwort erscheint und die Hochschuladresse den Server nur über zwei benannte Wege verlässt. Die 71 Testklassen des Servers laufen gegen dieselbe PostgreSQL-Datenbank wie die Anwendung, weil aufgeschobene Bedingungen und Ausdrucksindizes in einer eingebetteten Ersatzdatenbank nicht existieren. Die Schemaprüfung beim Start sichert die Deckung über alle 28 Migrationen, und die Fehlercodes werden mit der Tabelle des API-Contracts deckungsgleich geführt.

Auf der Seite der Oberfläche stellen 140 Testdateien Serverantworten über abgefangene HTTP-Aufrufe. Strukturelle Tests halten fest, dass kein Farbwert außerhalb der Wertedateien steht und beide Sprachdateien denselben Schlüsselbestand tragen. Ein eigener Lauf prüft die gebaute Anwendung in 37 Prüffällen auf Barrierefreiheit; protokolliert ist er zuletzt an dem Stand, den der erste Begutachtende prüfte, nicht am getesteten Stand (Abschnitt~\ref{sec:pruefung-nfa}). Er prüft je Route regelbasiert und mit eigenen Prüfungen für Überschriftenstruktur, Tastaturdurchgang, Fokussichtbarkeit und die Darstellung bei 320 Bildpunkten, die meisten in hellem und dunklem Modus. Eine Sammlung von 101 Aufrufen fährt die Endpunkte vom öffentlichen Lesen bis zur Administration ab. Vor jedem Zusammenführen prüfen zusätzlich drei im Repository festgelegte, durch Sprachmodelle gestützte Prüfrollen gegen API-Contract und Konventionen; sie ersetzen weder Tests noch Lesen.

Grenzen bleiben: Ein Durchgang mit Vorlesesoftware steht aus, geprüft ist nur der Ausgangszustand jeder Seite, und Kontraste von Bedienelementen und über Bildflächen misst der Lauf nicht. Diese Lücken greift Abschnitt~\ref{sec:pruefung-nfa} wieder auf.

	% ---------------------------------------------------------------------------
% Kapitel 6: Test
% Kurzfassung vom 04.10.2026, neu geschrieben aus der Langfassung
% (thesis/langfassung/kapitel06.tex, Confluence-Seite 51642395).
% Tabelle 6.3 fuehrt Aufgabenkatalog und Kennzahlen je Aufgabe zusammen;
% die Verteilung der Schweregrade und die Befunde je Heuristik stehen im Text.
% Entfallen: Abbildungen Testaufbau und Sitzungsablauf.
% ---------------------------------------------------------------------------

\chapter{Test}
\label{kap:test}

Die Kriterien aus Abschnitt~\ref{sec:usability} und der Entwurf aus Abschnitt~\ref{sec:ui-design} sind begründet, aber nicht geprüft. Die zweite Teilfrage verlangt diese Prüfung am gebauten Portal, der Akzeptanzteil der Forschungsfrage zusätzlich das Urteil künftiger Nutzender. Das Kapitel berichtet deskriptiv; bewertet wird in Kapitel~\ref{kap:evaluation}.

\section{Testkonzept und Methodik}
\label{sec:testkonzept}

Gebrauchstauglichkeit und Akzeptanz werden mit vier Methoden erhoben, die sich gegenseitig ergänzen, weil jede bei den Fragen der anderen blind ist (Tabelle~\ref{tab:methoden}). Die ersten drei erfassen den Aufwand der Benutzung, erst die Skala der wahrgenommenen Nützlichkeit ihren Wert. Ein Problem, das lautes Denken und heuristische Evaluation gemeinsam melden, wiegt in der Bewertung mehr als eines aus nur einer Quelle.

\begin{table}[htbp]
	\centering
	\footnotesize
	\caption[Vier Erhebungsmethoden, ihre Frage und ihr blinder Fleck]{Vier Erhebungsmethoden, ihre Frage und ihr blinder Fleck (eigene Darstellung)}
	\label{tab:methoden}
	\begin{tabularx}{\textwidth}{@{}L{2.4cm}YYY@{}}
		\toprule
		\textbf{Methode} & \textbf{beantwortet} & \textbf{beantwortet nicht} & \textbf{Ergebnis} \\
		\midrule
		Lautes Denken & woran eine Person scheitert & ob ein Problem repräsentativ ist & Probleme mit Ursache, Wertung, Zeit \\
		Heuristische Evaluation & Regelverstöße, auch ohne Aufgabe & ob Nutzende darüber stolpern & Befunde mit Schweregrad \\
		SUS-Bogen & wahrgenommene Gebrauchstauglichkeit & welche Stelle den Wert senkt & ein Wert je Person, 0 bis 100 \\
		TAM-Bogen & Nützlichkeit, Benutzerfreundlichkeit, Nutzungsabsicht & ob aus der Absicht Nutzung wird & drei Werte je Person, 1,00 bis 7,00 \\
		\bottomrule
	\end{tabularx}
\end{table}

Die heuristische Evaluation läuft vor den Sitzungen, damit ein Verstoß, der jede Person aufhielte, nicht alle Sitzungen belastet. Sie folgt Nielsen und Molich (1990) mit den zehn Heuristiken aus Tabelle~\ref{tab:heuristiken}. Die beiden Verfasser prüfen getrennt in zwei Durchgängen, erst im Überblick, dann Bildschirm für Bildschirm, und vergeben je Heuristik zusätzlich einen Erfüllungsstatus, damit der Vergleich mit der Referenzplattform möglich wird. Zusammengeführt wird, was am selben Bildschirm dieselbe Ursache hat; weichen die Grade ab, gilt der höhere. Den Schweregrad setzt Nielsen (1994c) aus Häufigkeit, Auswirkung und Persistenz zusammen. Damit zwei Begutachtende zum selben Wert kommen, übersetzt die vorliegende Arbeit das in die Regel der Tabelle~\ref{tab:schweregrad}; kehrt ein Problem unumgehbar wieder, steigt der Grad um eine Stufe.

\begin{table}[htbp]
	\centering
	\footnotesize
	\caption[Schweregrad aus Auswirkung und Häufigkeit]{Schweregrad aus Auswirkung und Häufigkeit; 1 kosmetisch, 2 gering, 3 schwerwiegend, 4 katastrophal (eigene Darstellung nach Nielsen, 1994c)}
	\label{tab:schweregrad}
	\begin{tabularx}{\textwidth}{@{}L{4.6cm}YYY@{}}
		\toprule
		\textbf{Auswirkung auf die Aufgabe} & \textbf{genau eine Person} & \textbf{mehr als eine, höchstens die Hälfte} & \textbf{mehr als die Hälfte} \\
		\midrule
		hoch: Aufgabe scheitert & 3 & 3 & 4 \\
		mittel: deutlicher Umweg & 2 & 3 & 3 \\
		gering: kurzes Stocken & 1 & 2 & 2 \\
		\bottomrule
	\end{tabularx}
\end{table}

Für das laute Denken ist Nielsen (1994b) die methodische Referenz. Getestet wird moderiert und in Präsenz, mit vorgelesenem Moderationstext und getrennten Rollen für Moderation und Protokoll, damit alle Sitzungen dasselbe messen. Der SUS-Bogen nach Brooke (1996) liegt wie der TAM-Bogen als deutsche Übersetzung neben dem Original vor (beide in Anhang~\ref{anh:testunterlagen}); sein Zielwert von 68 ist als NFA-06 vorab gesetzt. Der TAM-Bogen erhebt nach Davis (1989) auf einer siebenstufigen Skala je sechs Items zu PU und PEOU und zwei Items zur Nutzungsabsicht (BI). Die Items sind übersetzt und nur im Bezug angepasst („das Portal“ statt des Systemnamens, „diese Aufgaben“ statt der Erwerbsarbeit). Kodiert wird nach Position: Das Feld „äußerst wahrscheinlich“ zählt 7, „äußerst unwahrscheinlich“ 1, sodass hohe Werte Zustimmung bedeuten; bei BI2 sind die Endpunkte vertauscht gedruckt, der Wert geht deshalb als 8 minus Rohwert ein. Da Davis keine Normwerte berichtet, dient der Skalenmittelpunkt 4,00 als festgelegter Orientierungspunkt; die Reliabilität wird als Cronbachs Alpha mit der Untergrenze 0,80 nach Davis (1989) berichtet.

Block A bearbeiten alle elf Personen, Block B fünf, die Blöcke C und D je zwei. Nach Nielsen (2000) finden fünf Personen bei einer typischen Aufdeckungsrate von 31 Prozent rund 85 Prozent der Probleme, und für verschiedene Nutzergruppen empfiehlt er mindestens drei je Gruppe. Block B erreicht diese Untergrenze, C und D nicht. Ausgewertet wird deshalb deskriptiv, Erfolgsquoten stets mit dem Nenner ihres Blocks, Rangkorrelationen nach Spearman gelten als explorativ. Weitere Grenzen sind die Befangenheit der Begutachtenden, die das System selbst gebaut haben, eine einzige Testrunde statt mehrerer kleiner Runden, nicht eigens validierte Übersetzungen der Bögen und das Wissen der Teilnehmenden, dass die Testleitenden das Portal entwickelt haben.

\section{Testdesign (Teilnehmende, Aufgaben, Ablauf)}
\label{sec:testdesign}

Geplant waren fünf Studierende für Block B, zwei Lehrende in der Rolle der Redaktion für Block C, ersatzweise Studierende in einer Rollenaufgabe, zwei Hochschulangehörige für Block D und zwei Externe ohne Hochschulbezug. Ausgeschlossen war, wer am Projekt beteiligt ist, den Prototyp kennt oder die Arbeit betreut. Jede Person beginnt mit Block A ohne Anmeldung, weil sich Auffindbarkeit nur aus der Sicht der Hauptzielgruppe prüfen lässt; so entstehen elf Beobachtungen dazu statt zwei.

Jede Aufgabe hat eine vorgelesene Karte, einen vorab hergestellten Startzustand, ein Erfolgskriterium und ein Zeitbudget (Tabelle~\ref{tab:aufgaben}; Wortlaut in Anhang~\ref{anh:testunterlagen}). Die Karte nennt ein Ziel statt eines Weges und meidet das Vokabular der Oberfläche: Zu B6 heißt es „Sorge dafür, dass du wieder weiterschreiben kannst“, nicht „Ziehe die Einreichung zurück“. Die Karten der Blöcke C und D erklären dagegen den Ablauf mit den Wörtern der Oberfläche. A3 verlangt, Abschluss, Fakultät und Semester zugleich einzuschränken, A4, eine Einschränkung ohne Umweg über die Startseite zu entfernen und dann eine Profilseite zu öffnen; beide prüfen zwei der drei Kriterien der Auffindbarkeit.

Eine Vollsitzung dauert 76 Minuten. Davon entfallen 48 auf die beiden Aufgabenblöcke: 18 Minuten für Block A und 30 für den Rollenblock, darin 43 Minuten Zeitbudget der Aufgaben aus Tabelle~\ref{tab:aufgaben} und 5 Minuten Übergänge zwischen den Aufgaben. Externe Personen bearbeiten in 44 Minuten nur Block A. Beide Bögen folgen auf die letzte Aufgabe, getrennt durch eine Minute Pause. Die Moderation greift nach einer festen Leiter ein, die zugleich die Wertung bestimmt: Nachfrage bei Schweigen („erfolgreich“), Verweis auf die Karte nach etwa einer Minute Ratlosigkeit („mit Mühe“), Hinweis auf den Bereich bei überschrittenem Budget oder Bitte um Hilfe („nicht gelöst“), Abbruch bei Budget plus 50 Prozent. Getestet wird auf einer eigenen Testinstanz mit 48 veröffentlichten Projekten über mehrere Fakultäten, Semester und alle drei Abschlussarten, mit frischem Browserprofil und Testkonto je Sitzung. Die Unterlagen führen allein die Pseudonyme TN01 bis TN11.

\section{Durchführung des Usability-Tests}
\label{sec:durchfuehrung}

Alle elf Sitzungen liefen vollständig auf demselben Stand. Gegenüber dem Stand, den der erste Begutachtende prüfte, kamen ein dreistufiges Metadatenformular, vereinheitlichte Feldnamen sowie eine überarbeitete Filterzeile hinzu; der zweite prüfte überwiegend am getesteten Stand. Keiner der 21 heuristischen Befunde ist vor den Sitzungen als behoben geführt; die drei mit Grad 3 liefen dokumentiert mit.

Teilgenommen haben neun Studierende, fünf aus dem Bauingenieurwesen, zwei aus der Medieninformatik, je eine aus Pharmazie und Wirtschaftsinformatik, sowie zwei externe Berufstätige. Ihre Erfahrung mit Webanwendungen lag im Median bei 4 auf einer Skala von 1 bis 5, sechs hatten schon ein Portfolio-Portal benutzt. Lehrende nahmen nicht teil; Block C beruht deshalb auf Rollenaufgaben zweier Studierender. Die Moderation wechselte nicht wie geplant: Ein Verfasser moderierte sieben Sitzungen, der andere vier, nie zugleich mit dem Protokoll.

\begin{table}[htbp]
	\centering
	\footnotesize
	\setlength{\tabcolsep}{3.5pt}
	\renewcommand{\arraystretch}{1.02}
	\caption[Testaufgaben mit Zeitbudget und Kennzahlen]{Testaufgaben mit Zeitbudget, Wertungen und Zeiten in Minuten und Sekunden; bei $n=2$ ist der Median das abgerundete Mittel (eigene Erhebung, elf Sitzungen)}
	\label{tab:aufgaben}
	\begin{tabular}{@{}L{0.8cm}L{5.4cm}rrrrrrr@{}}
	\toprule
	\textbf{ID} & \textbf{Aufgabe} & \textbf{Budget} & \textbf{n} & \textbf{erf.} & \textbf{Mühe} & \textbf{n.\,gel.} & \textbf{Median} & \textbf{Max} \\
	\midrule
	A1 & Orientierung (ohne Wertung) & 2 & 11 & & & & 01:30 & 02:47 \\
	A2 & Bestimmtes Projekt finden & 3 & 11 & 11 & 0 & 0 & 00:45 & 01:10 \\
	A3 & Facettiert einschränken & 4 & 11 & 11 & 0 & 0 & 00:59 & 01:50 \\
	A4 & Zurücknehmen und zum Profil & 4 & 11 & 10 & 1 & 0 & 00:55 & 01:47 \\
	A5 & Projekt als PDF sichern & 3 & 11 & 8 & 3 & 0 & 01:14 & 02:40 \\
	B1 & Anmelden & 2 & 5 & 5 & 0 & 0 & 00:40 & 00:46 \\
	B2 & Projekt anlegen & 3 & 5 & 4 & 0 & 1 & 02:33 & 03:04 \\
	B3 & Inhalt mit mehreren Bausteinarten & 8 & 5 & 4 & 1 & 0 & 02:48 & 03:11 \\
	B4 & Sieben Pflichtangaben pflegen & 6 & 5 & 5 & 0 & 0 & 02:32 & 08:00 \\
	B5 & Einreichen & 2 & 5 & 5 & 0 & 0 & 00:20 & 00:23 \\
	B6 & Wieder bearbeitbar machen & 3 & 5 & 5 & 0 & 0 & 00:25 & 01:27 \\
	B7 & Eigenes Profil ergänzen & 3 & 5 & 5 & 0 & 0 & 00:45 & 01:58 \\
	C1 & Anmelden & 2 & 2 & 2 & 0 & 0 & 00:41 & 00:44 \\
	C2 & Arbeitsvorrat überblicken & 2 & 2 & 1 & 1 & 0 & 02:03 & 02:30 \\
	C3 & Prüfen und zurückgeben & 8 & 2 & 2 & 0 & 0 & 01:06 & 01:33 \\
	C4 & Genehmigen und veröffentlichen & 6 & 2 & 1 & 0 & 1 & 01:50 & 03:00 \\
	C5 & Für die Top 10 vorschlagen & 5 & 2 & 1 & 1 & 0 & 01:27 & 01:50 \\
	C6 & Interne Notiz hinterlassen & 4 & 2 & 2 & 0 & 0 & 01:06 & 01:12 \\
	D1 & Anmelden & 2 & 2 & 2 & 0 & 0 & 00:42 & 00:45 \\
	D2 & Lagebild verschaffen & 3 & 2 & 2 & 0 & 0 & 01:20 & 01:30 \\
	D3 & Jemandem das Prüfen erlauben & 3 & 2 & 2 & 0 & 0 & 01:55 & 02:05 \\
	D4 & Doppelte Kategorie aufräumen & 5 & 2 & 1 & 1 & 0 & 02:55 & 03:20 \\
	D5 & Mehrere Projekte veröffentlichen & 4 & 2 & 1 & 1 & 0 & 02:52 & 03:05 \\
	D6 & Die Top 10 festlegen & 7 & 2 & 2 & 0 & 0 & 02:05 & 02:15 \\
	D7 & Eine Änderung nachvollziehen & 3 & 2 & 2 & 0 & 0 & 01:42 & 01:50 \\
	\bottomrule
	\end{tabular}
\end{table}

In Block A wurden alle 44 Bearbeitungen von A2 bis A5 gelöst, vier davon mit Mühe. Block B kommt auf 34 von 35, Block C auf 11 von 12, Block D auf 14 von 14. Nicht gelöst wurden C4 bei TN07 und B2 bei TN02; TN02 verstand B2 als Auftrag, ein Projekt anzuzeigen, was wie ein Missverständnis in A4 nicht als Usability-Problem zählt. Vier Wertungen weichen von den Regeln ab und stehen wie protokolliert: B2 bei TN05, B4 bei TN03 und C2 bei TN06 überschritten ihr Budget, ohne als nicht gelöst zu gelten, und A4 bei TN06 gilt trotz Umwegs über die Startseite als erfolgreich (P8).

Die elf SUS-Werte lauten aufsteigend 67,5; 72,5; 85,0; 85,0; 87,5; 90,0; 92,5; 95,0; 95,0; 95,0; 95,0. Der Mittelwert beträgt 87,3, der Median 90,0, die Standardabweichung 9,5. Tabelle~\ref{tab:tam} enthält die Kennwerte des TAM-Bogens; die Werte je Block stehen in Anhang~\ref{anh:tam-bloecke}. Davis (1989) berichtet Alphawerte von 0,98 für PU und 0,94 für PEOU, die Untergrenze liegt nach Davis (1989) bei 0,80. Die explorative Rangkorrelation beträgt für PU mit BI 0,84, für PEOU mit BI 0,41, für PEOU mit PU 0,39 und für SUS mit PEOU 0,73.

\begin{table}[htbp]
	\centering
	\footnotesize
	\caption[Kennwerte des TAM-Bogens]{Kennwerte des TAM-Bogens, 7 = äußerst wahrscheinlich, BI2 vor der Mittelung gedreht (eigene Erhebung, elf Sitzungen)}
	\label{tab:tam}
	\begin{tabular}{@{}lrrrrrrl@{}}
		\toprule
		\textbf{Konstrukt} & \textbf{n} & \textbf{Median} & \textbf{Mittelw.} & \textbf{Min} & \textbf{Max} & \textbf{SD} & \textbf{Cronbachs Alpha} \\
		\midrule
		PU & 11 & 6,00 & 5,97 & 4,50 & 7,00 & 0,91 & 0,92 \\
		PEOU & 11 & 6,17 & 5,92 & 4,17 & 7,00 & 0,85 & 0,89 \\
		BI & 11 & 7,00 & 6,41 & 4,00 & 7,00 & 0,92 & entfällt (zwei Items) \\
		\bottomrule
	\end{tabular}
\end{table}

Die Problemliste in Anhang~\ref{anh:problemliste} umfasst 30 Usability-Probleme: zwölf aus den Sitzungen, davon drei auch heuristisch gefunden, und 18 allein aus der heuristischen Evaluation. Keines trägt Grad 4, sieben tragen Grad 3, 16 Grad 2 und sieben Grad 1. Tabelle~\ref{tab:probleme} führt die zwölf Probleme aus den Sitzungen. Acht davon traten bei genau einer Person auf, sieben stehen unter H6. Protokolliert sind Äußerungen wie „Wo ist denn hier der Download?“ (TN09, P1) und „Hat das jetzt geklappt?“ (TN10, P12). Vorgeschlagen wurden das Zuschneiden von Bildern und eine Top 10 nach Fakultäten.

\begingroup
\footnotesize
\setlength{\tabcolsep}{3.5pt}
\renewcommand{\arraystretch}{1.02}
\setlength{\LTcapwidth}{\textwidth}
\begin{longtable}{@{}L{0.8cm}L{0.8cm}L{7.4cm}L{1.4cm}L{0.8cm}L{0.8cm}L{1.3cm}@{}}
	\caption[Usability-Probleme aus den Sitzungen]{Usability-Probleme aus den Sitzungen; Quelle LD lautes Denken, beide zusätzlich heuristisch (eigene Erhebung)}
	\label{tab:probleme}\\
	\toprule
	\textbf{ID} & \textbf{Aufg.} & \textbf{Beobachtung} & \textbf{Häufigk.} & \textbf{Grad} & \textbf{Heur.} & \textbf{Quelle} \\
	\midrule
	\endfirsthead
	\toprule
	\textbf{ID} & \textbf{Aufg.} & \textbf{Beobachtung} & \textbf{Häufigk.} & \textbf{Grad} & \textbf{Heur.} & \textbf{Quelle} \\
	\midrule
	\endhead
	\bottomrule
	\endfoot
	P1 & A5 & Export- und Druckfunktion nicht auf Anhieb gefunden & 5 von 11 & 3 & H6 & LD \\
	P2 & B4 & Übernahme eines Mitglieds über „+“ nicht erkannt & 3 von 5 & 2 & H6 & beide \\
	P3 & B4 & Pflichtfelder und Bedeutung der Sternchen nicht erkannt & 2 von 5 & 3 & H4 & beide \\
	P4 & B3 & Bausteinarten Überschrift und Aufzählung nicht rasch gefunden & 1 von 5 & 2 & H6 & LD \\
	P5 & C2 & Arbeitsvorrat zuerst unter „Mein Portfolio“ gesucht & 1 von 2 & 2 & H6 & LD \\
	P6 & C4 & Abfolge Genehmigen und Veröffentlichen unklar, Schaltfläche nicht gefunden & 1 von 2 & 3 & H6 & beide \\
	P7 & C5 & Nominiersymbol auf dem Projektbild nicht erkannt, kein Tooltip & 2 von 2 & 3 & H6 & LD \\
	P8 & A4 & Einschränkung über Umweg zur Startseite zurückgenommen & 1 von 11 & 2 & H3 & LD \\
	P9 & A1 & Fortsetzung der Startseite nach unten nicht bemerkt & 1 von 11 & 1 & H1 & LD \\
	P10 & D4 & Folgen für zugeordnete Projekte unklar & 1 von 2 & 2 & H5 & LD \\
	P11 & D4 & Zuerst Projektliste statt Kategorien geöffnet & 1 von 2 & 2 & H6 & LD \\
	P12 & D5 & Keine Rückmeldung nach der Sammelaktion & 1 von 2 & 2 & H1 & LD \\
\end{longtable}
\endgroup

Die 21 heuristischen Befunde verteilen sich auf H4 mit acht, H2 mit vier, H6 mit drei, H8 mit zwei sowie H1, H5, H7 und H9 mit je einem; H3 und H10 tragen keinen. Der Befund unter H1 ist P17: Die wirkende Einschränkung steht allein in der Filterspalte. Die Zusammenführung hatte ihn zunächst H6 zugeordnet; nach einer Durchsicht der Zuordnungen steht er unter H1, weil er die Sichtbarkeit des Zustands der Suche betrifft. Grad 3 tragen vor dem Test drei Befunde: Auf Projektkarten mit hellem Titelbild sinkt der Kontrast des Titels bis auf 1,05:1 (P16), ungespeicherte Metadaten gehen beim Verlassen des Formulars ohne Warnung verloren (P28), und der im Editor angekündigte Schrägstrich öffnet das Auswahlmenü nicht (P30). Mit der im Test gezählten Häufigkeit stiegen P3 und P6 auf Grad 3.

	% ---------------------------------------------------------------------------
% Kapitel 7: Evaluation
% Kurzfassung vom 04.10.2026, neu geschrieben aus der Langfassung
% (thesis/langfassung/kapitel07.tex, Confluence-Seite 51642418).
% Tabelle 7.4 ersetzt die frueheren Abgleichtabellen 7.4 und 7.6 und die
% Abbildung 7.4; Tabelle 7.3 nimmt die Befundzahlen der Abbildung 7.3 auf.
% Entfallen: Abbildungen TAM-Werte je Person und Heuristikvergleich.
% ---------------------------------------------------------------------------

\chapter{Evaluation}
\label{kap:evaluation}

Ob das Portal leistet, was es soll, entscheidet sich an Maßstäben, die seit den Abschnitten~\ref{sec:ziele}, \ref{sec:nichtfunktionale-anforderungen} und~\ref{sec:usability} feststehen.

\section{Bewertungsrahmen auf Basis der Forschungsliteratur}
\label{sec:bewertungsrahmen}

Ein Maßstab, der erst nach der Messung gewählt wird, lässt sich dem Ergebnis anpassen. Bewertet wird deshalb an vorab gesetzten Zielwerten und Kriterien (Tabelle~\ref{tab:bewertungsrahmen}), wobei Literaturwerte eine Einordnung über das eigene System hinaus erlauben und Festlegungen der Arbeit nur eine Richtung angeben. Für die TAM-Skalen gibt es keine Norm; die Rangkorrelationen werden an der Anordnung gemessen, die Davis (1989) und Shroff et al. (2011) berichten, nicht an deren Beträgen.

\begin{table}[htbp]
	\centering
	\footnotesize
	\caption[Bewertungsfragen, Maßstab und Herkunft]{Bewertungsfragen, Maßstab und Herkunft (eigene Darstellung)}
	\label{tab:bewertungsrahmen}
	\begin{tabularx}{\textwidth}{@{}L{3.6cm}YL{4.4cm}@{}}
		\toprule
		\textbf{Frage} & \textbf{Maßstab} & \textbf{Herkunft} \\
		\midrule
		Gebrauchstauglichkeit & SUS-Wert mindestens 68, Einordnung auf der Notenskala & Brooke (1996), Lewis und Sauro (2018) \\
		Nützlichkeit, Benutzerfreundlichkeit, Nutzungsabsicht & Skalenwert über 4,00 & Festlegung der Arbeit, Konstrukte nach Davis (1989) \\
		Auffindbarkeit & Systematik als Auswahl (H6), Rücknahme (H3), sichtbare Einschränkung (H1) & Ableitung aus Nielsen (2024) \\
		Heuristische Qualität im Vergleich & Erfüllungsstatus je Heuristik & Nielsen und Molich (1990), Nielsen (2024) \\
		Inhaltlicher Zuschnitt & zehn Komponenten & Mahasneh (2020) \\
		Funktionsumfang und Qualität & Wortlaut der Anforderung & Tabellen~\ref{tab:fa} und~\ref{tab:nfa} \\
		\bottomrule
	\end{tabularx}
\end{table}

Drei Anwendungsregeln sind nachträglich festgelegt und als solche ausgewiesen. Weichen die Erfüllungsurteile der beiden Begutachtenden ab, was bei H5, H7 und H9 der Fall ist, gilt wie bei den Schweregraden das strengere. Zusätzlich gilt eine Heuristik nur als erfüllt, wenn die zusammengeführte Liste keinen Befund zu ihr führt; nach dieser Regel steht H1 auf teilweise erfüllt, seit P17 dort und nicht mehr unter H6 geführt wird (Abschnitt~\ref{sec:durchfuehrung}). Probleme aus den Sitzungen gehen in diesen Status nicht ein, damit der Vergleich mit der allein begutachteten Referenzplattform methodengleich bleibt. NFA-10 und NFA-11 werden an Block B gemessen, weil allein diese Personen eigene Arbeiten angelegt und eingereicht haben. Eine Anforderung gilt als erfüllt, wenn jedes Glied ihres Wortlauts belegt ist, als teilweise erfüllt, wenn ein Glied fehlt oder nur für einen Teil belegt ist, und als nicht geprüft, wenn für sie kein Verfahren vorlag. Als Nachweis zählen Nutzertest, Fragebogen, heuristische Evaluation, automatisierte Tests und die Inspektion von Programmtext und API-Contract. Aussagen aus den Blöcken C und D gelten als Hinweis auf eine Stelle, nicht als Quote.

\section{Bewertung der Testergebnisse}
\label{sec:bewertung-test}

\subsection{Erfolgsquoten und Auffindbarkeit}
\label{subsec:erfolg-auffindbarkeit}

Vier Wertungen stehen anders, als die Regeln aus Abschnitt~\ref{sec:testdesign} verlangt hätten. Tabelle~\ref{tab:lesarten} zeigt, dass die Bewertung nicht an ihnen hängt: Auch in der strengen Lesart wurde jede Aufgabe mindestens einmal gelöst, und an keinem Kernpfad blieb eine Stelle, die niemand überwand. Die ungelösten Bearbeitungen in Block C, C4 und C2, decken sich mit P6 und P5.

\begin{table}[htbp]
	\centering
	\footnotesize
	\caption[Gelöste Bearbeitungen je Block in zwei Lesarten]{Gelöste Bearbeitungen je Block in zwei Lesarten; nach Regel zählen Budgetüberschreitungen und der Umweg in A4 als nicht gelöst (eigene Auswertung nach Tabelle~\ref{tab:aufgaben})}
	\label{tab:lesarten}
	\begin{tabular}{@{}lrrrr@{}}
		\toprule
		\textbf{Block} & \textbf{Personen} & \textbf{Bearbeitungen} & \textbf{gelöst wie protokolliert} & \textbf{gelöst nach Regel} \\
		\midrule
		A, ohne Anmeldung & 11 & 44 & 44 & 43 \\
		B, Studierende & 5 & 35 & 34 & 32 \\
		C, Redaktion & 2 & 12 & 11 & 10 \\
		D, Administration & 2 & 14 & 14 & 14 \\
		\midrule
		Summe & 11 & 105 & 103 & 99 \\
		\bottomrule
	\end{tabular}
\end{table}

Das erste Kriterium der Auffindbarkeit, die Systematik als auswählbare Merkmale (H6), trägt A3: Alle elf Personen schränkten Abschluss, Fakultät und Semester zugleich ein, im Median in 59 Sekunden bei einem Budget von vier Minuten, auch die beiden ohne Hochschulbezug. Das zweite, die Rücknahme (H3), trägt A4 mit elf Lösungen, eine davon über einen Umweg (P8); die vollständige Rücknahme ist über das Zurücksetzen belegt, und die heuristische Evaluation fand unter H3 keinen Befund. Beide Kriterien sind eingelöst.

Das dritte, die Sichtbarkeit der wirkenden Einschränkung (H1), trägt der Entwurf nur zum Teil. Die Zählmarke am Gruppenkopf zeigt, dass eingeschränkt ist, aber nicht worauf, und die gewählten Werte stehen allein in der Filterspalte, die in der schmalen Ansicht vollständig über der Liste liegt (P17). Wer die Treffer durchsieht, kann eine unerwartete Trefferliste nicht mehr auf die eigene Auswahl zurückführen. Mittelbar berühren zwei weitere Probleme die Auffindbarkeit: Titel über hellen Bildern sind kaum lesbar (P16), und die Sicherungsfunktion fanden fünf von elf nicht auf Anhieb (P1). Nicht gesondert belegt ist der Anteil des Suchfelds am Erfolg in A2, weil der gewählte Weg nicht protokolliert wurde. Die Profilseite als zweiter Zugang erreichten in A4 alle elf.

\subsection{Benutzerfreundlichkeit}
\label{subsec:benutzerfreundlichkeit}

Der SUS-Mittelwert von 87,3 liegt 19,3 Punkte über dem Zielwert, und der Median von 90,0 zeigt, dass nicht einzelne Höchstwerte das Ergebnis tragen. Abbildung~\ref{abb:sus} ordnet die Werte ein: Neun Werte und der Mittelwert liegen in der Note A+, die ab 84,1 beginnt und dem 96.\ bis 100.\ Perzentil der Vergleichsstudien entspricht. Die beiden niedrigsten Werte, 72,5 und 67,5, gehören zu Personen, deren Rollenblock Schwierigkeiten zeigte (TN07 mit C4, TN05 mit B2 und B3); die Gegenrichtung gilt nicht, sodass nur ein Zusammentreffen belegt ist, kein Gleichlauf.

\abbildung{abb-7-1-sus-notenskala.png}{0.85}{abb:sus}%
{SUS-Werte der elf Personen auf der Notenskala}%
{SUS-Werte der elf Personen auf der Notenskala (eigene Darstellung, Notenskala wiedergegeben nach Lewis \& Sauro, 2018)}

Bei PU und PEOU liegen alle elf Personen über dem Orientierungspunkt 4,00, bei BI zehn, TN05 genau darauf. Das Urteil fällt damit auf allen drei Skalen eher zustimmend aus. Die fünf Studierenden aus Block B, deren Projekte den Bestand füllen sollen, schätzen die Nützlichkeit mit 5,27 niedriger ein als alle elf (5,97); bei fünf Personen ist das ein Hinweis. Wie bei Davis (1989) hängt die Nutzungsabsicht stärker mit der Nützlichkeit (0,84) als mit der Benutzerfreundlichkeit (0,41) zusammen. Den Pfad von PEOU zu PU, den Shroff et al. (2011) stützen, zeigt diese Stichprobe dagegen nicht (0,39); bei elf Personen ist das ein Befund über die Stichprobe und kein Widerspruch zum Modell.

Gegen die hohen Urteile stehen 30 Probleme, sieben mit Grad 3, mit zwei Schwerpunkten. Unter H6 stehen sieben der zwölf Sitzungsprobleme, durchweg Handlungen hinter unbeschrifteten Symbolen oder an unerwarteten Orten: die Entscheidung der Redaktion hinter der Statusmarke (P6), das Nominiersymbol (P7), das Pluszeichen für Mitglieder (P2). Unter H4 stehen acht heuristische Befunde, etwa wechselnde Anrede und zwei Namen für dieselbe Zuweisung. Vier der sieben schweren Probleme gehen auf den Nutzertest zurück; P1 und P7 kannte die Befundliste nicht, P3 und P6 führte sie niedriger. Die drei übrigen (P16, P28, P30) fand allein die heuristische Evaluation. Die Methoden haben sich damit ergänzt, wie Abschnitt~\ref{sec:testkonzept} erwartet.

\subsection{Vergleich mit der Referenzplattform}
\label{subsec:vergleich}

Tabelle~\ref{tab:heuristikvergleich} stellt beide Systeme je Heuristik nebeneinander. Bei H7 lag der einzige Negativbefund der Referenzplattform; das eigene System führt dort im öffentlichen Bereich keinen Befund, sein einziger (P30) liegt im Editor. Suchfeld und verbindbare Filter sind vorhanden, das aus H7 abgeleitete Kriterium ist im öffentlichen Bereich eingelöst. H3 und H10 steigen auf erfüllt. H1 bleibt wie bei der Referenzplattform teilweise erfüllt: In beiden Systemen zeigt allein das Bedienelement, welche Einschränkung wirkt (P17). Umgekehrt verliert das eigene System bei H2 und H4: Es zeigt mehr Funktionen als die Referenzplattform, aber nicht überall mit demselben Wort und Muster. H8 ist allein wegen zweier Befunde in der Administration teilweise erfüllt.

\begin{table}[htbp]
	\centering
	\footnotesize
	\caption[Erfüllungsstatus je Heuristik: Referenzplattform und eigenes System]{Erfüllungsstatus je Heuristik: Referenzplattform und eigenes System, Befunde der heuristischen Evaluation und Probleme aus den Sitzungen (eigene Darstellung nach Tabelle~\ref{tab:heuristiken} und den Einzelurteilen der Begutachtenden)}
	\label{tab:heuristikvergleich}
	\begin{tabular}{@{}llllrr@{}}
		\toprule
		\textbf{Nr.} & \textbf{Referenzplattform} & \textbf{eigenes System} & \textbf{Befunde} & \textbf{davon öffentl.} & \textbf{Sitzungen} \\
		\midrule
		H1 & teilweise & teilweise & 1 & 1 & 2 \\
		H2 & erfüllt & teilweise & 4 & 3 & 0 \\
		H3 & teilweise & erfüllt & 0 & 0 & 1 \\
		H4 & erfüllt & teilweise & 8 & 5 & 1 \\
		H5 & nicht beurteilbar & teilweise & 1 & 0 & 1 \\
		H6 & teilweise & teilweise & 3 & 1 & 7 \\
		H7 & nicht erfüllt & teilweise & 1 & 0 & 0 \\
		H8 & erfüllt & teilweise & 2 & 0 & 0 \\
		H9 & nicht beurteilbar & teilweise & 1 & 1 & 0 \\
		H10 & teilweise & erfüllt & 0 & 0 & 0 \\
		\bottomrule
	\end{tabular}
\end{table}

Der Vergleich ist qualitativ zu lesen. Beide Prüfungen stammen von denselben zwei Personen, die Referenzplattform wurde vier Monate früher und stichprobenartig begutachtet, und beim eigenen System urteilten die Begutachtenden über ihren eigenen Entwurf, was eine Verzerrung zu dessen Gunsten wahrscheinlicher macht als die umgekehrte.

\subsection{Antwort auf die zweite Teilfrage}
\label{subsec:antwort-tf2}

Die zweite Teilfrage lautet: „Welche funktionalen und gestalterischen Merkmale fördern die Auffindbarkeit und Benutzerfreundlichkeit eines institutionellen Portfolio-Portals?“ Bewährt haben sich verbindbare Facettenfilter entlang der Ordnung der Hochschule, eine einzeln rücknehmbare Einschränkung, die Profilseite als zweiter Zugang, ein kurzer Einreichweg und ein ohne Einweisung bedienbarer Editor; Einreichen und Zurückziehen dauerten im Median 20 und 25 Sekunden. Nicht bestätigt hat sich die Annahme, Trefferzahl und Zählmarken zeigten den Zustand der Suche ausreichend: Die wirkende Einschränkung muss bei den Treffern stehen. Neu ist eine Bedingung, unter der die Merkmale wirken: Handlungen und Zustände müssen beschriftet sichtbar sein, und dasselbe muss überall gleich heißen. Die Antwort beruht auf elf Personen und einer Testrunde, in mehreren Gruppen auf weniger als drei Personen.

\section{Überprüfung funktionaler Anforderungen}
\label{sec:pruefung-fa}

Jede der 36 funktionalen Anforderungen wurde gegen Programmtext, API-Contract, automatisierte Tests und Testaufgaben gehalten; 23 kamen zusätzlich in Aufgaben vor. Tabelle~\ref{tab:status} fasst das Ergebnis nach Bereichen zusammen: 30 sind erfüllt, sechs teilweise, keine ist nicht erfüllt. Die sechs Einschränkungen haben vier Ursachen. Das Suchfeld durchsucht nur den Namen der Autorin oder des Autors, nicht den von Mitwirkenden (FA-14). Der Auswahlzustand ist nur in der Filterspalte sichtbar (FA-15, P17). Einen Studiengang der Person führt das Datenmodell nicht (FA-16), und eine eigene Übersicht der Studiengänge und Module fehlt (FA-18); in beiden Fällen hat die Umsetzung den Einordnungskontext auf das beschränkt, was der Filter braucht. Schließlich kann die Administration fremde Projekte und Dateien löschen und beim Zusammenführen die Kategorie veröffentlichter Projekte ändern, was der API-Contract einräumt, der Wortlaut von FA-42 aber ausschließt; beim Löschen eines Projekts verliert zudem das Protokoll der Startseitenauswahl seinen Verweis darauf (FA-44).

\begin{table}[htbp]
	\centering
	\footnotesize
	\caption[Status der Anforderungen nach Bereich]{Status der funktionalen und nicht-funktionalen Anforderungen nach Bereich; aufgeführt sind die nicht vollständig erfüllten Kennungen (eigene Auswertung)}
	\label{tab:status}
	\begin{tabularx}{\textwidth}{@{}L{4.2cm}rYY@{}}
		\toprule
		\textbf{Bereich} & \textbf{erfüllt} & \textbf{teilweise erfüllt} & \textbf{nicht geprüft} \\
		\midrule
		FA rollenübergreifend & 6 von 6 & & \\
		FA ohne Anmeldung & 4 von 8 & FA-14, FA-15, FA-16, FA-18 & \\
		FA Studierende & 10 von 10 & & \\
		FA Redaktion & 7 von 7 & & \\
		FA Administration & 3 von 5 & FA-42, FA-44 & \\
		\midrule
		NFA Gebrauchstauglichkeit & 3 von 6 & NFA-02, NFA-03, NFA-05 & \\
		NFA Akzeptanz & 2 von 2 & & \\
		NFA Barrierefreiheit & 1 von 7 & NFA-20, NFA-22 bis NFA-26 & \\
		NFA Antwortzeiten & 0 von 3 & & NFA-30, NFA-31, NFA-32 \\
		NFA Geräte und Browser & 0 von 2 & NFA-40, NFA-41 & \\
		NFA Corporate Design & 1 von 1 & & \\
		NFA Sicherheit & 3 von 3 & & \\
		\midrule
		Summe & 40 von 60 & 17 & 3 \\
		\bottomrule
	\end{tabularx}
\end{table}

Den inhaltlichen Zuschnitt prüft das Komponentenmodell von Mahasneh (2020). Sieben der zehn Komponenten deckt das Portal ab: die Biografie über die Profilseite, Berichte und Forschungsarbeiten über die Projekttypen „Paper“ und „Thesis“, Projekte als Grundeinheit, Zusammenfassungen über Kurzbeschreibung und Textbausteine, wissenschaftliches Material über Formel, Quelltext, Tabelle und Dokument, Audio- und Videoclips über eigene Bausteine sowie Leistungsproben über veröffentlichte Projekte mit Bild und Galerie. Kursplan, Hausaufgaben und Aktivitäten setzen eine laufende Lehrveranstaltung voraus und entfallen mit der Abgrenzung aus Abschnitt~\ref{sec:ziele}; die Auslassung ist eine Entscheidung und keine Lücke. Untertitel und Transkripte für Audio und Video fehlen allerdings.

\section{Überprüfung nicht-funktionaler Anforderungen}
\label{sec:pruefung-nfa}

Von den 24 nicht-funktionalen Anforderungen sind zehn erfüllt, elf teilweise und drei nicht geprüft (Tabelle~\ref{tab:status}). Erfüllt sind sie vor allem dort, wo ein Fragebogen einen vorab gesetzten Zielwert prüft oder ein Test eine Regel vollständig abdeckt. NFA-06 ist mit 87,3 gegenüber dem Literaturwert 68 erfüllt (Lewis \& Sauro, 2018), NFA-10 und NFA-11 mit PU 5,27 und PEOU 5,80 in Block B, wobei alle fünf Personen über dem Orientierungspunkt liegen. NFA-01 belegt Aufgabe B3, die alle fünf ohne Einweisung in höchstens 3:11 Minuten lösten. Die Sicherheitsanforderungen NFA-60 bis NFA-62 belegen Integrationstests, die unberechtigte Zugriffe abweisen, Entwürfe und interne Notizen verbergen und den Passworthash nie ausliefern.

Von den Anforderungen, die Regeln nach Nielsen (2024) übersetzen, sind Einheitlichkeit (NFA-02) und Wiedererkennen (NFA-05) genau an den beiden Schwerpunkten aus Abschnitt~\ref{subsec:benutzerfreundlichkeit} nur teilweise erfüllt, Fehlervermeidung (NFA-03) vor allem wegen des Datenverlusts beim Verlassen des Formulars (P28). Dem hohen SUS-Wert widerspricht das nicht, denn Brooke (1996) ordnet den Bogen als grobes Verfahren ein.

Die Barrierefreiheit nach W3C (2023) ist nur teilweise belegt. Der automatisierte Lauf prüft nur den Ausgangszustand jeder Seite, misst Kontraste von Bedienelementen und Text über Bildern nicht, kam ohne Vorlesesoftware aus und ist am getesteten Stand nicht protokolliert. Belegt sind zudem drei Verstöße: Kartentitel über hellen Bildern erreichen 1,05:1 statt 4,5:1 (P16), Ränder von Ankreuz- und Eingabefeldern 1,4:1 bis 1,8:1 statt 3:1, und verdeckte Bilder der Beteiligtengruppe bieten 16 statt 24 Pixel Zielbreite (NFA-24). Belegt ist damit der Textkontrast der Tokenschicht in beiden Modi, nicht die Barrierefreiheit insgesamt. Antwortzeiten und Tragfähigkeit (NFA-30 bis NFA-32) sind wie vorab festgehalten nicht geprüft; die Darstellung ist nur in einem Chromium-Browser, für die Desktopbreite und zwei schmale Breiten belegt (NFA-40, NFA-41). Was hier nicht geprüft oder nur teilweise belegt ist, geht als Grenze in Kapitel~\ref{kap:zusammenfassung} ein.

	% ---------------------------------------------------------------------------
% Kapitel 8: Zusammenfassung und Ausblick
% Kurzfassung vom 04.10.2026, neu geschrieben aus der Langfassung
% (thesis/langfassung/kapitel08.tex, Confluence-Seite 58097676).
% ---------------------------------------------------------------------------

\chapter{Zusammenfassung und Ausblick}
\label{kap:zusammenfassung}

\section{Zusammenfassung}
\label{sec:gesamtzusammenfassung}

Die vorliegende Arbeit hat ein institutionelles E-Portfolio-Portal für die Hochschule Bremen entworfen und gebaut. Aus Forschungsliteratur und heuristischer Analyse einer Referenzplattform entstand ein Katalog aus 36 funktionalen und 24 nicht-funktionalen Anforderungen, die Konzeption und Realisierung prüften, ob er einer Umsetzung standhält, und vier Methoden mit elf Personen prüften das Ergebnis an vorab gesetzten Maßstäben. 30 funktionale Anforderungen sind erfüllt und sechs teilweise, von den nicht-funktionalen zehn erfüllt, elf teilweise und drei ungeprüft. Der SUS-Mittelwert liegt mit 87,3 um 19,3 Punkte über dem Zielwert von 68 (Lewis \& Sauro, 2018); ihm stehen 30 Usability-Probleme gegenüber, sieben davon schwerwiegend.

\subsection{Antwort auf die erste Teilfrage}
\label{subsec:antwort-tf1}

Die erste Teilfrage lautet: „Welche Anforderungen an Struktur, Inhalt und Workflow lassen sich aus der bestehenden Forschungsliteratur zu E-Portfolios ableiten?“ Zur Struktur liefert die Literatur zwei Anforderungen. Aus der institutionellen Kategorie und der offenen Leserschaft nach Lorenzo und Ittelson (2005) folgen der Zugang ohne Anmeldung und das Gewicht der externen Nutzenden; die Übertragung auf studentische Arbeiten ist ein Schritt der vorliegenden Arbeit. Aus den nach ihrem Bezugspunkt geordneten Komponenten von Mahasneh (2020) folgt die Trennung von Profilseite und Projekt. Zum Inhalt liefert dasselbe Modell das Raster, von dem das Portal sieben Komponenten abdeckt und drei veranstaltungsgebundene bewusst auslässt. Zum Workflow trägt die Bestimmung des institutionellen Portfolios als fokussierte Auswahl bei (Kahn, 2001, zitiert nach Lorenzo \& Ittelson, 2005), aus der die redaktionelle Freigabe folgt. Quer dazu begründet die Akzeptanzforschung (Davis, 1989; Shroff et al., 2011), dass Gebrauchstauglichkeit eine Frage des Entwurfs ist. Von den neun funktionalen Anforderungen mit Herkunft aus der E-Portfolio-Literatur sind acht erfüllt; der neunten (FA-16) fehlt allein der Studiengang der Person. Nicht aus der Literatur folgen die fünf Zustände des Freigabewegs, die Ableitung der Rolle aus der Hochschuladresse, die Ordnung der Hochschule und die kuratierte Startseite.

\subsection{Antwort auf die zweite Teilfrage}
\label{subsec:antwort-tf2-kurz}

Die zweite Teilfrage lautet: „Welche funktionalen und gestalterischen Merkmale fördern die Auffindbarkeit und Benutzerfreundlichkeit eines institutionellen Portfolio-Portals?“ Die Auffindbarkeit fördern sichtbar angebotene, verbindbare Ordnungsmerkmale, eine einzeln rücknehmbare Einschränkung und ein Zugang über die Person: Alle elf Personen schränkten drei Merkmale zugleich ein, im Median in 59 Sekunden, und alle fanden die Profilseite. Die Benutzerfreundlichkeit fördern kurze Wege für wiederkehrende Vorgänge und ein ohne Einweisung bedienbarer Editor. Nicht bestätigt hat sich, dass Trefferzahl und Zählmarken den Zustand der Suche ausreichend zeigen. Hinzu kommt eine Bedingung: Handlungen und Zustände müssen beschriftet sichtbar sein (H6), und dasselbe muss überall gleich heißen (H4). Gegenüber der Referenzplattform schließt das Portal die Lücke bei H7 im öffentlichen Bereich und bleibt bei H2 und H4 zurück.

\subsection{Antwort auf die dritte Teilfrage}
\label{subsec:antwort-tf3}

Die dritte Teilfrage lautet: „Wie lässt sich ein blockbasiertes Inhaltsmodell mit integriertem Redaktionsworkflow realisieren, das Studierenden eine flexible und medienreiche Projektpräsentation ermöglicht?“ Die Realisierung bestätigt vier Entscheidungen. Das Inhaltsmodell ist an das Dokumentmodell des Editors gebunden, sodass eine schmale Zuordnung zwischen Editor und Speicherung genügt. Der Inhalt eines Bausteins wird unverändert gespeichert, seine Art zusätzlich als abfragbares Merkmal geführt. Dateien liegen außerhalb des Dokuments und werden über einen eigenen Weg ausgeliefert, weil ihre Sichtbarkeit am Projektzustand hängt. Und die Bearbeitbarkeit ist an diesen Zustand gebunden und auf dem Server durchgesetzt, auch für Dateien; darin besteht die Integration des Workflows in das Inhaltsmodell. Für mehrere Bearbeitende verlangt das automatische Speichern ein Versionsfeld, das Konflikte meldet, und eine Bearbeitungssperre, die sie selten macht. Die Flexibilität tragen fünfzehn Bausteinarten nach dem Komponentenmodell von Mahasneh (2020). Dass Flexibilität nicht immer mit Benutzerfreundlichkeit einhergeht (Davis, 1989), zeigt der Test an einer Stelle (P4), und der Schrägstrichbefehl im Editor ist ein Fehler der Umsetzung (P30).

\subsection{Antwort auf die Forschungsfrage}
\label{subsec:antwort-ff}

Die Forschungsfrage lautet: „Welche funktionalen und gestalterischen Anforderungen muss ein hochschulisches E-Portfolio-Portal erfüllen, um studentische Projekte strukturiert zu präsentieren, auffindbar zu machen und von den beteiligten Nutzergruppen akzeptiert zu werden?“ Strukturiert präsentiert ein Portal, das zwischen Projekt und Person trennt, Projekte aus frei anordenbaren Bausteinen aufbaut und vor der Veröffentlichung redaktionell freigibt. Auffindbar macht es seinen Bestand, wenn es die Ordnung der Hochschule als verbindbare, einzeln rücknehmbare Merkmale anbietet, die wirkende Einschränkung bei den Treffern zeigt und einen Zugang über die Person eröffnet. Akzeptiert wird es, wenn sein Nutzen erkennbar ist und Zustände wie Handlungen beschriftet sichtbar sind. Bei Nützlichkeit und Benutzerfreundlichkeit liegen alle elf Befragten über dem Orientierungspunkt, und wie bei Davis (1989) hängt die Nutzungsabsicht stärker mit der Nützlichkeit zusammen (0,84 gegenüber 0,41). Übertragbar erscheinen der Zugang ohne Anmeldung, die Freigabe vor der Veröffentlichung, die Trennung von Projekt und Person und die sichtbare Beschriftung von Zustand und Handlung. An die Hochschule Bremen gebunden bleiben die Ausgestaltung des Freigabewegs, die Ableitung der Rolle, die kuratierte Startseite und die Ordnung aus fünf Fakultäten mit 75 Studiengängen.

\subsection{Grenzen der Arbeit}
\label{subsec:grenzen}

Die empirische Grundlage ist schmal: elf Personen in einer einzigen Testrunde, Redaktion, Administration und externe Nutzende mit je zwei Personen unter der Mindestzahl nach Nielsen (2000), keine Lehrenden, und in den Rollenblöcken erklärten die Karten den Ablauf. Die heuristische Evaluation stammt von zwei Begutachtenden, die das System selbst gebaut haben (Nielsen \& Molich, 1990). Beide Bögen sind nicht eigens validiert, der TAM-Bogen erfragt die Absicht nach einer Sitzung statt nach einem Semester, und die Teilnehmenden kannten die Rolle der Testleitenden, sodass die hohen Werte soziale Erwünschtheit und einen Deckeneffekt enthalten können. Die Anforderungen stützen sich auf eine einzige, auf Empfehlung gewählte Referenzplattform und keine Vorabbefragung. Antwortzeiten und Tragfähigkeit blieben ungeprüft, die Barrierefreiheit ist nur automatisiert und im Ausgangszustand belegt, und getestet wurde ein Prototyp mit 48 vorbereiteten Projekten.

\section{Ausblick}
\label{sec:ausblick}

Die wirkende Einschränkung gehört über die Trefferliste (P17), Handlungen hinter Symbolen und Statusmarken brauchen beschriftete Schaltflächen (P1, P2, P6, P7), und Begriffe sind zu vereinheitlichen. Zu beheben sind außerdem die Kontraste über Bildern und an Feldrändern, die Zielgröße der Personenbilder, der Datenverlust beim Verlassen eines Formulars (P28) und der Schrägstrichbefehl (P30). Offen sind die Suche über Mitwirkende, der Studiengang auf der Profilseite, eine Übersicht der Studiengänge und Module, Untertitel und Transkripte sowie der Abgleich der Löschbefugnisse mit FA-42. Danach wäre nach Nielsen (2000) in mehreren kleinen Runden erneut zu testen, mit mindestens drei Personen je Gruppe und Lehrenden in der Redaktion.

Zwei Anregungen kamen aus der Hochschule selbst. Der Erstprüfer regte an, auf der Startseite zufällig gewählte Projekte aus jeder Fakultät zu zeigen (A. Teufel, persönliche Mitteilung, 8. September 2026), und eine Testperson schlug eine Top 10 nach Fakultäten vor. Ein Portal, das nach Lorenzo und Ittelson (2005) eine Institution repräsentiert, sollte deren Breite schon auf der ersten Seite zeigen; die vorliegende Arbeit schlägt eine kuratierte Top 10 je Fakultät vor, aus der oben wechselnd ein Projekt erscheint. Zweitens schätzten zwei Mitarbeitende der Zentralen Studienberatung in einem Gespräch außerhalb der Erhebung ein, dass Schülerinnen und Schüler vor der Studienwahl die Projekte nicht ohne Weiteres verstehen. Eine Ansicht in verständlicher Sprache, etwa mit Hilfe eines Sprachmodells erzeugt, wäre ein Weg dorthin; zu klären wäre ihr Verhältnis zum Freigabeweg, und weil Gebrauchstauglichkeit kontextgebunden ist (Brooke, 1996), verlangt diese Gruppe einen eigenen Test.

Am Inhaltsmodell wäre die gleichzeitige Bearbeitung durch mehrere Mitwirkende der nächste Schritt, und Kommentare an einzelnen Bausteinen setzen stabile Bausteinkennungen voraus. Für einen Produktivbetrieb fehlen eine serverseitig wirksame Abmeldung, Sperren und Anmeldebegrenzung über mehrere Server sowie die Anonymisierung der Autorschaft beim Löschen eines Kontos. Die Eigenentwicklung trägt zudem nach Lorenzo und Ittelson (2005) das Risiko, das Fachpersonal nicht zu binden; ein Betrieb setzt voraus, dass eine Stelle der Hochschule Wartung und Weiterentwicklung übernimmt, wofür API-Contract, Dokumentation und automatisierte Tests die Grundlage bilden. Offen bleibt schließlich, ob aus der gemessenen Nutzungsabsicht tatsächlich Nutzung wird; das klärt erst eine Erhebung im laufenden Betrieb.


	% =======================================================================
	% NACHSPANN
	% =======================================================================
	% Solange die Belege im Text als reiner Text stehen, sorgt \nocite dafuer,
	% dass alle Eintraege aus literatur.bib im Verzeichnis erscheinen.
	% Sobald im Text \textcite und \parencite benutzt werden, kann die Zeile weg.
	\nocite{*}
	\printbibliography[title={Literaturverzeichnis}]

	\appendix
	% ---------------------------------------------------------------------------
% Anhang
% Uebernommen aus der Confluence-Seite "Anhang (Entwurf)"
% (Space BEH, Seiten-ID 101220353), Stand 04.10.2026.
% Angepasst wurden allein Verweise: Tabellen- und Abschnittsnummern laufen
% ueber \ref, der Verweis auf die entfallene Abbildung 2.2 zeigt auf
% Abschnitt 2.5. In Tabelle B.1 sind Quelle und Haeufigkeit abgekuerzt.
% ---------------------------------------------------------------------------

\chapter{User Stories und Use Cases}
\label{anh:user-stories}

Jede User Story folgt dem Format „Als [Gruppe] möchte ich [Ziel], damit [Nutzen]“ und nennt die funktionale Anforderung aus Tabelle~\ref{tab:fa}, die sie einlöst. Die Nummern stehen in Blöcken je Akteurgruppe. 34 der 36 funktionalen Anforderungen tragen eine Story; ohne Story bleiben FA-04 und FA-35, weil beide eine Eigenschaft des Systems beschreiben und keine Handlung einer Person. User Stories und Use Cases beschreiben den Sollzustand aus Kapitel~\ref{kap:anforderungsanalyse}. Wo der gebaute Stand davon abweicht, führt Abschnitt~\ref{sec:pruefung-fa} es aus; das betrifft den Studiengang auf der Profilseite (US-04, US-16, UC-05) und die Löschbefugnisse der Administration, die über US-34 hinausgehen.

\newcommand{\usTabelle}[3]{%
\begingroup
\footnotesize
\setlength{\tabcolsep}{4pt}
\renewcommand{\arraystretch}{1.08}
\setlength{\LTcapwidth}{\textwidth}
\begin{longtable}{@{}L{1.0cm}L{11.3cm}L{2.2cm}@{}}
	\caption[#1]{#1 (eigene Darstellung)}
	\label{#2}\\
	\toprule
	\textbf{ID} & \textbf{User Story} & \textbf{Bezug} \\
	\midrule
	\endfirsthead
	\toprule
	\textbf{ID} & \textbf{User Story} & \textbf{Bezug} \\
	\midrule
	\endhead
	\bottomrule
	\endfoot
	#3
\end{longtable}
\endgroup}

\usTabelle{User Stories der externen Nutzenden}{tab:us-extern}{%
	US-01 & Als externe Person möchte ich veröffentlichte Projekte über mehrere verbindbare Filtergruppen und ein Suchfeld eingrenzen, meine Auswahl sehen, einzeln zurücknehmen und als Adresse weitergeben können, damit ich schnell zu den Arbeiten eines bestimmten Studiengangs finde und den Fund teilen kann. & FA-10, FA-12, FA-14, FA-15 \\
	US-02 & Als externe Person möchte ich ein Projekt mit allen seinen Bausteinen und beschreibenden Angaben ansehen, damit ich mir ein Bild von der studentischen Arbeit machen kann. & FA-11 \\
	US-03 & Als studieninteressierte Person möchte ich Studiengänge und ihre Module einsehen können, damit ich verstehe, in welchem Zusammenhang ein Projekt entstanden ist. & FA-18 \\
	US-04 & Als externe Person, etwa eine Arbeitgeberin, möchte ich die Profilseite einer studierenden Person mit Biografie, Studiengang und allen ihren veröffentlichten Projekten ansehen, damit ich einen Gesamteindruck bekomme, statt einzelne Projekte getrennt suchen zu müssen. & FA-16, FA-29 \\
	US-05 & Als externe Person möchte ich eine Rückfrage oder eine Meldung an die Hochschule senden können, ohne mich anzumelden, damit ich auf einen Fehler hinweisen oder Kontakt aufnehmen kann. & FA-06 \\
	US-06 & Als externe Person möchte ich auf der Startseite eine kuratierte Auswahl veröffentlichter Projekte finden, damit ich einen Einstieg habe, ohne die Systematik des Portals zu kennen. & FA-17 \\
}

\usTabelle{User Stories der Studierenden}{tab:us-studierende}{%
	US-10 & Als studierende Person möchte ich ein Projekt in einem blockbasierten Editor aus Bausteinen für Text und multimediale Inhalte aufbauen, damit ich meine Arbeit flexibel darstellen kann, ohne technisches Vorwissen zu brauchen. & FA-20 \\
	US-11 & Als studierende Person möchte ich mein Projekt als Entwurf anlegen, zwischenspeichern, weiter bearbeiten und wieder löschen können, damit ich in mehreren Sitzungen daran arbeiten kann, ohne etwas zu verlieren. & FA-21 \\
	US-12 & Als studierende Person möchte ich die beschreibenden Angaben meines Projekts pflegen, die zuständige Redaktion wählen und das Projekt anschließend zur Prüfung einreichen, damit es von den richtigen Personen geprüft wird. & FA-22, FA-23 \\
	US-13 & Als studierende Person möchte ich jederzeit den Zustand meiner Projekte sehen, damit ich weiß, ob ich noch handeln muss. & FA-25 \\
	US-14 & Als studierende Person möchte ich die Änderungsaufforderungen der Redaktion lesen, einzeln als erledigt kennzeichnen und danach erneut einreichen, damit nachvollziehbar bleibt, welcher Punkt behandelt ist. & FA-26 \\
	US-15 & Als studierende Person möchte ich mein Projekt bearbeiten können, solange es im Entwurf oder in der Überarbeitung ist, damit ein bereits geprüfter Stand nicht unbemerkt verändert wird. & FA-21 \\
	US-16 & Als studierende Person möchte ich meine Profilseite mit Biografie, Studiengang, Kontaktangaben und Bild pflegen, damit sich Besuchende einen persönlichen Eindruck verschaffen können. & FA-28 \\
	US-17 & Als studierende Person möchte ich, dass jedes veröffentlichte Projekt, an dem ich beteiligt bin, ohne mein Zutun auf meiner Profilseite erscheint, damit ich nichts einzeln verlinken oder pflegen muss. & FA-29 \\
	US-18 & Als studierende Person möchte ich eine Einreichung zurückziehen können, solange die Redaktion nicht entschieden hat, damit ich einen Fehler bemerken und beheben kann, ohne die Redaktion damit zu befassen. & FA-24 \\
	US-19 & Als studierende Person möchte ich eine Gruppenarbeit als gemeinsames Projekt führen, Mitwirkende aufnehmen und ihre Funktion benennen, damit die Arbeit im Portfolio aller Beteiligten erscheint und niemand sie doppelt anlegen muss. & FA-30 \\
}

\usTabelle{User Stories der Redaktion}{tab:us-redaktion}{%
	US-20 & Als Mitglied der Redaktion möchte ich in „Meine Reviews“ die Projekte sehen, für die ich als zuständige Redaktion eingetragen bin, damit ich meine Arbeit ordnen kann, ohne die Einreichungen anderer mitzuführen. & FA-31 \\
	US-21 & Als Mitglied der Redaktion möchte ich ein geprüftes Projekt genehmigen können, damit festgestellt ist, dass es inhaltlich in Ordnung geht. & FA-32 \\
	US-22 & Als Mitglied der Redaktion möchte ich ein Projekt mit einer Begründung zur Überarbeitung zurückgeben können, damit die einreichende Person weiß, was zu ändern ist. & FA-34 \\
	US-23 & Als Mitglied der Redaktion möchte ich ein genehmigtes Projekt in einem zweiten Schritt veröffentlichen können, damit sich geprüfte Arbeiten gebündelt freischalten lassen, etwa zum Semesterbeginn. & FA-33 \\
	US-24 & Als Mitglied der Redaktion möchte ich zwischen einer Antwort, einer für Studierende unsichtbaren internen Notiz und einer Änderungsaufforderung mit Erledigungsstand wählen, damit ich mich mit anderen Mitgliedern der Redaktion abstimmen kann, ohne die einreichende Person zu verwirren. & FA-36 \\
	US-25 & Als Mitglied der Redaktion möchte ich ein veröffentlichtes Projekt für die Top 10 der Startseite vorschlagen können, damit gute Arbeiten sichtbar werden, ohne dass ich selbst über die Startseite entscheide. & FA-37 \\
}

\usTabelle{User Stories der Administration}{tab:us-administration}{%
	US-30 & Als Mitglied der Administration möchte ich Konten anlegen, Rollen ändern, Konten deaktivieren und Passwörter zurücksetzen, damit ich Ausnahmefälle bearbeiten kann, für die die Selbstregistrierung nicht greift. & FA-40 \\
	US-31 & Als Mitglied der Administration möchte ich alle Projekte unabhängig von ihrem Zustand einsehen können, damit ich bei Rückfragen auskunftsfähig bin. & FA-41 \\
	US-32 & Als Mitglied der Administration möchte ich die Stammdaten pflegen, also Kategorien, Fakultäten, Studiengänge, Module und Semester, damit die Einordnung der Projekte aktuell bleibt. & FA-43 \\
	US-33 & Als Mitglied der Administration möchte ich die Top 10 der Startseite kuratieren, Vorschläge annehmen oder ablehnen und die Reihenfolge festlegen, wobei jede Änderung unveränderlich protokolliert wird, damit die Außendarstellung verantwortet und nachvollziehbar bleibt. & FA-44 \\
	US-34 & Als Mitglied der Administration möchte ich fremde Projektinhalte lesen, aber nicht verändern können, damit die ausgewiesene Autorschaft eines Textes belastbar bleibt und ein Eingriff nur über die regulären Zustandsentscheidungen möglich ist. & FA-42 \\
}

\usTabelle{User Stories aller angemeldeten Personen}{tab:us-angemeldet}{%
	US-40 & Als Mitglied der Hochschule möchte ich mich mit einer privaten Adresse als dauerhafter Kennung und einer bestätigten Hochschuladresse registrieren, damit mein Portfolio auch nach der Exmatrikulation erreichbar bleibt und meine Zugehörigkeit trotzdem belegt ist. & FA-01, FA-02 \\
	US-41 & Als angemeldete Person möchte ich mich anmelden, mein Passwort zurücksetzen und ein von der Administration gesetztes Initialpasswort beim ersten Zugang wechseln, damit mein Konto mir gehört und nicht der Person, die es angelegt hat. & FA-03 \\
	US-42 & Als angemeldete Person möchte ich auf ein veröffentlichtes Projekt reagieren können, höchstens einmal je Projekt und rücknehmbar, damit Zuspruch sichtbar wird, ohne dass eine Kommentarspalte moderiert werden müsste. & FA-05 \\
}

Die fünf Use Cases fassen zusammenhängende User Stories zu einem durchgängigen Ablauf zusammen; ihre Kennungen und Titel entsprechen Abschnitt~\ref{sec:user-stories}. Einstufige Vorgänge wie das Zurückziehen (US-18), die Reaktion (US-42), die Rückfrage ohne Anmeldung (US-05) und die Stammdatenpflege (US-32) tragen keinen eigenen Use Case. Die Kuratierung der Top 10 (US-25, US-33) ist mehrstufig, wird aber im Workflow (Abschnitt~\ref{sec:workflow}) behandelt und nicht als eigener Use Case geführt.

\newenvironment{usecase}[1]{%
	\par\medskip\noindent\textbf{#1}\par\smallskip
	\begin{description}[leftmargin=0pt, itemsep=0.15em, topsep=0.2em, font=\normalfont\itshape]}
	{\end{description}}

\begin{usecase}{UC-01: Projekt erstellen und einreichen}
	\item[Akteur:] studierende Person, ebenso eine mitwirkende Person einer Gruppenarbeit
	\item[Bezug:] US-10, US-11, US-12, US-19
	\item[Vorbedingung:] Die Person ist angemeldet.
	\item[Auslöser:] Die Person möchte eine abgeschlossene Arbeit veröffentlichen lassen.
	\item[Ablauf:] (1) Die Person legt ein neues Projekt an; es entsteht mit einem vorläufigen Titel im Zustand Entwurf (US-11). (2) Sie füllt die Projektseite im Editor mit Bausteinen (US-10) und speichert beliebig oft zwischen. (3) Sie pflegt die beschreibenden Angaben und wählt die zuständige Redaktion; bei einer Gruppenarbeit nimmt sie hier die Mitwirkenden auf und benennt deren Funktion (US-12, US-19). (4) Sie reicht ein. Der Zustand wechselt auf Eingereicht.
	\item[Nachbedingung:] Das Projekt erscheint in „Meine Reviews“ der gewählten Redaktion. Eine Benachrichtigung gibt es nicht; das System versendet nur Nachrichten rund um das Konto.
	\item[Alternativablauf:] Die Person bricht ab; das Projekt bleibt als Entwurf gespeichert.
	\item[Fehlerfall:] Fehlt eine Pflichtangabe oder ist keine Redaktion gewählt, weist der Server das Einreichen ab und benennt, was fehlt.
\end{usecase}

\begin{usecase}{UC-02: Projekt-Review durch Redaktion}
	\item[Akteur:] die dem Projekt zugewiesene Redaktion, ebenso die Administration, die dafür keine Zuweisung braucht
	\item[Bezug:] US-20, US-21, US-22, US-23, US-24
	\item[Vorbedingung:] Mindestens ein zugewiesenes Projekt steht auf Eingereicht.
	\item[Ablauf:] (1) Die Redaktion öffnet „Meine Reviews“ und wählt ein eingereichtes Projekt (US-20). Die Liste führt die zugewiesenen Projekte in vier Zuständen; Entwürfe erscheinen dort nicht. (2) Sie öffnet die Prüfansicht des Projekts und sieht Inhalt, beschreibende Angaben und die bisherige Kommunikation. (3) Sie entscheidet: Genehmigen führt auf Genehmigt (US-21), das Projekt ist damit noch nicht öffentlich. Veröffentlichen führt ein genehmigtes Projekt in einem zweiten, eigenen Schritt auf Veröffentlicht (US-23). Zur Überarbeitung zurückgeben verlangt eine Änderungsaufforderung und führt auf Überarbeitung (US-22, US-24); derselbe Übergang entsteht, wenn die Redaktion am Projekt eine Änderungsaufforderung hinterlässt.
	\item[Nachbedingung:] Der Zustandswechsel steht in der Historie des Projekts; eine Rückgabe hinterlässt zusätzlich die Änderungsaufforderung.
	\item[Einschränkung:] Das Vier-Augen-Prinzip sperrt Genehmigen und Veröffentlichen für jede Person, die an dem Projekt beteiligt ist, gleich welcher Rolle (FA-35). Die Rückgabe zur Überarbeitung ist davon nicht betroffen.
\end{usecase}

\begin{usecase}{UC-03: Überarbeitung und erneute Einreichung}
	\item[Akteur:] studierende Person, ebenso eine mitwirkende Person
	\item[Bezug:] US-14, US-15
	\item[Vorbedingung:] Das Projekt steht auf Überarbeitung und trägt mindestens eine offene Änderungsaufforderung.
	\item[Ablauf:] (1) Die Person öffnet das Projekt und liest die Änderungsaufforderungen der Redaktion (US-14). (2) Sie passt die betroffenen Stellen im Editor an (US-15). (3) Sie kennzeichnet jede umgesetzte Änderungsaufforderung als erledigt (US-14). (4) Sie reicht erneut ein; der Zustand wechselt zurück auf Eingereicht.
	\item[Nachbedingung:] Das Projekt durchläuft erneut UC-02.
	\item[Fehlerfall:] Solange eine Änderungsaufforderung offen ist, weist der Server das erneute Einreichen ab; Schritt 3 ist die Bedingung für Schritt 4.
	\item[Alternativablauf:] Die zuständige Redaktion oder die Administration kann ein Projekt selbst aus der Überarbeitung zurück in die Prüfung holen, etwa wenn sich die Rückgabe als gegenstandslos erweist.
\end{usecase}

\begin{usecase}{UC-04: Externe Projektsuche}
	\item[Akteur:] externe Person ohne Konto
	\item[Bezug:] US-01, US-02, US-06
	\item[Vorbedingung:] keine; eine Anmeldung ist an keiner Stelle nötig.
	\item[Ablauf:] (1) Die Person ruft die Startseite auf und findet dort die Top 10, oder sie geht unmittelbar zur Projektübersicht (US-06). (2) Sie grenzt die Übersicht über die sechs Filtergruppen Fakultät, Studiengang, Abschluss, Semester, Projekttyp und Kategorie ein. Innerhalb einer Gruppe wirken mehrere Werte als Oder, zwischen den Gruppen als Und. Daneben steht ein Suchfeld über Titel, Kurzbeschreibung, Personennamen und Schlagworte (US-01). (3) Die Auswahl bleibt sichtbar, lässt sich einzeln oder gesammelt zurücknehmen und steht in der Adresse (US-01). (4) Die Person öffnet ein Projekt und sieht alle seine Bausteine (US-02).
	\item[Nachbedingung:] Sichtbar sind ausschließlich veröffentlichte Projekte. Entwürfe, Einreichungen, genehmigte Projekte und Überarbeitungen bleiben ohne Anmeldung unerreichbar.
\end{usecase}

\begin{usecase}{UC-05: Portfolio-Profil pflegen und extern einsehen}
	\item[Akteure:] studierende Person für die Pflege, externe Person für die Einsicht
	\item[Bezug:] US-16, US-17, US-04
	\item[Vorbedingung:] für die Pflege ein Konto, für die Einsicht keine
	\item[Ablauf:] (1) Die studierende Person pflegt ihre Profilseite mit Biografie, Studiengang, Kontaktangaben und Bild (US-16). (2) Sie lässt im Lauf der Zeit Projekte über UC-01 und UC-02 veröffentlichen. (3) Jedes veröffentlichte Projekt, an dem sie beteiligt ist, erscheint ohne weiteres Zutun auf ihrer Profilseite; eine Gruppenarbeit steht im Portfolio aller Mitwirkenden (US-17). (4) Eine externe Person ruft die Profilseite auf, etwa von einer Projektseite aus, und sieht Biografie und alle veröffentlichten Projekte dieser Person (US-04).
	\item[Nachbedingung:] Die Profilseite bleibt aktuell, sobald ein weiteres Projekt veröffentlicht wird.
	\item[Abgrenzung:] Die Profilseite selbst durchläuft keinen Freigabeweg; nur die Projektliste darauf hängt vom Zustand der einzelnen Projekte ab.
\end{usecase}

\chapter{Problemliste}
\label{anh:problemliste}

Die Liste führt alle 30 Usability-Probleme aus Abschnitt~\ref{sec:durchfuehrung}. P1 bis P12 stammen aus den Sitzungen, P13 bis P30 allein aus der heuristischen Evaluation. Die Häufigkeit steht bei Problemen aus den Sitzungen mit dem Nenner des jeweiligen Blocks, bei Befunden der heuristischen Evaluation als Zahl der Begutachtenden, die den Befund fanden (1 Beg. oder 2 Beg.). Der Schweregrad folgt Tabelle~\ref{tab:schweregrad}, die Heuristik den Bezeichnungen aus Tabelle~\ref{tab:heuristiken}; die Quelle unterscheidet lautes Denken (LD), heuristische Evaluation (HE) und beide.

\begingroup
\scriptsize
\setlength{\tabcolsep}{3pt}
\renewcommand{\arraystretch}{1.1}
\setlength{\LTcapwidth}{\textwidth}
\begin{longtable}{@{}L{0.6cm}L{0.8cm}L{4.9cm}L{4.0cm}L{1.2cm}L{0.7cm}L{0.7cm}L{0.8cm}@{}}
	\caption[Usability-Probleme P1 bis P30]{Usability-Probleme P1 bis P30 (eigene Erhebung)}
	\label{tab:problemliste}\\
	\toprule
	\textbf{ID} & \textbf{Aufg.} & \textbf{Beobachtung} & \textbf{Auswirkung} & \textbf{Häufigk.} & \textbf{Grad} & \textbf{Heur.} & \textbf{Quelle} \\
	\midrule
	\endfirsthead
	\toprule
	\textbf{ID} & \textbf{Aufg.} & \textbf{Beobachtung} & \textbf{Auswirkung} & \textbf{Häufigk.} & \textbf{Grad} & \textbf{Heur.} & \textbf{Quelle} \\
	\midrule
	\endhead
	\bottomrule
	\endfoot
	P1 & A5 & Die Export- beziehungsweise Druckfunktion auf der Projektseite wurde nicht auf Anhieb gefunden; eine Person wollte mit Strg+P drucken, eine übersah die Schaltfläche auf dem Bild. & Die Aufgabe gelang bei drei Personen nur mit Mühe, bei zwei mit Umweg. & 5 von 11 & 3 & H6 & LD \\
	P2 & B4 & Nach Auswahl eines Mitglieds wurde nicht erkannt, dass es erst mit „+“ übernommen wird; eine Person erwartete das Verhalten der Redaktionsauswahl. & Kurzes Stocken beim Hinzufügen von Mitgliedern. & 3 von 5 & 2 & H6 & beide \\
	P3 & B4 & Welche Metadaten Pflichtfelder sind und was die Sternchen bedeuten, wurde nicht erkannt; eine Person fragte nach der Bedeutung der Sternchen. & Die Aufgabe dauerte bei einer Person 8:00 Minuten bei 6 Minuten Budget. & 2 von 5 & 3 & H4 & beide \\
	P4 & B3 & Die Bausteinarten Überschrift und Aufzählung wurden nicht schnell gefunden; zunächst wurde ein Bild für das ganze Projekt gesucht. & Die Aufgabe gelang nur mit Mühe. & 1 von 5 & 2 & H6 & LD \\
	P5 & C2 & Der Arbeitsvorrat wurde zunächst unter „Mein Portfolio“ gesucht, bevor „Meine Reviews“ geöffnet wurde. & Die Aufgabe gelang nur mit Mühe. & 1 von 2 & 2 & H6 & LD \\
	P6 & C4 & Dass ein Projekt erst genehmigt und dann veröffentlicht wird, wurde nicht verstanden; die Schaltfläche zum Genehmigen wurde nicht gefunden. & Die Aufgabe wurde nach 3:00 Minuten nicht gelöst. & 1 von 2 & 3 & H6 & beide \\
	P7 & C5 & Das Symbol zum Nominieren auf dem Projektbild wurde nicht erkannt; beim Überfahren erscheint kein Tooltip. & Eine Person brauchte einen Umweg über die Review-Seite, eine gelangte nur mit Mühe ans Ziel. & 2 von 2 & 3 & H6 & LD \\
	P8 & A4 & Zum Zurücknehmen einer Einschränkung wurde der Weg über die Startseite und erneut zu den Projekten gewählt. & Umweg bei einer an sich gelösten Aufgabe. & 1 von 11 & 2 & H3 & LD \\
	P9 & A1 & Dass die Startseite weiter nach unten reicht, wurde nicht bemerkt. & Inhalte unterhalb des ersten Bildschirms blieben zunächst unbeachtet. & 1 von 11 & 1 & H1 & LD \\
	P10 & D4 & Beim Aufräumen der doppelten Kategorie war unklar, was mit den zugeordneten Projekten geschieht; erwartet wurde ein Zusammenführen. & Die Aufgabe gelang mit zwei Fehlversuchen nur mit Mühe. & 1 von 2 & 2 & H5 & LD \\
	P11 & D4 & Für die Kategorien wurde zuerst die Projektliste geöffnet. & Ein Fehlversuch vor dem richtigen Bereich. & 1 von 2 & 2 & H6 & LD \\
	P12 & D5 & Nach der Sammelaktion erschien keine Rückmeldung, dass sie ausgeführt wurde. & Die Aufgabe gelang mit einem Fehlversuch nur mit Mühe. & 1 von 2 & 2 & H1 & LD \\
	P13 & keine & Startseite und Platzhalter des Editors duzen, die jeweils benachbarte Seite (Projektbibliothek, Metadatenformular) siezt; die deutsche Sprachdatei enthält 2 Stellen mit Du gegen 101 mit Sie. & Wer ein Projekt anlegt, wird in zwei aufeinanderfolgenden Schritten desselben Vorgangs verschieden angesprochen. & 1 Beg. & 2 & H4 & HE \\
	P14 & keine & Dasselbe Feld heißt im Formular „Typ“, im Einreichdialog „Art der Arbeit“ und im Filter „Projekttyp“; ebenso „Redakteur:in“ gegen „Redaktionszuweisung“. & Die Meldung beim Einreichen nennt Felder, die im Formular unter diesem Namen nicht stehen. & 1 Beg. & 2 & H4 & HE \\
	P15 & keine & Die Überschrift lautet „My Portfolio“, Navigation und Seitenpfad nennen den Bereich „Mein Portfolio“. & Auf der Seite steht ein anderer Name als auf dem Weg, der zu ihr führt. & 1 Beg. & 1 & H2 & HE \\
	P16 & keine & Titel und Studiengangszeile der Projektkarten stehen in heller Schrift auf dem Titelbild; auf hellen Bildern fällt der Kontrast bis auf 1,05:1 (Titel) und 1,22:1 (Studiengang), gefordert sind 4,5:1 beziehungsweise 3:1. & Auf Karten mit hellem Titelbild sind Titel und Studiengang nicht lesbar; die Karte ist nur über das Bild zu unterscheiden. & 2 Beg. & 3 & H4 & HE \\
	P17 & keine & Die wirkende Einschränkung steht nur in der Filterspalte; über der Trefferliste steht allein die Trefferzahl, in der schmalen Ansicht liegt die Filterspalte vollständig über der Liste. & Beim Durchsehen der Treffer ist nicht mehr zu sehen, worauf eingeschränkt wurde. & 1 Beg. & 2 & H1 & HE \\
	P18 & keine & Die Adresse einer gefilterten Liste führt vier der sechs Filtergruppen als technische Kennung, zwei als lesbaren Wert. & Ein geteilter Link lässt für vier Gruppen nicht erkennen, welche Einschränkung er enthält. & 1 Beg. & 1 & H2 & HE \\
	P19 & keine & Eine Liste einzuschränken geht über drei Bedienmuster: Ankreuzfelder in einer Spalte, eine Segmentleiste, Segmentleiste mit Auswahlliste. & Was an einer Liste gelernt wurde, gilt an der nächsten nicht. & 1 Beg. & 2 & H4 & HE \\
	P20 & keine & In der Kontenliste sind die E-Mail-Adressen abgeschnitten, rechts der Tabelle bleibt Fläche frei. & Die Adresse, über die ein Konto gesucht wird, ist in der Liste nicht vollständig lesbar. & 1 Beg. & 2 & H8 & HE \\
	P21 & keine & Unbesetzte Ränge der Top 10 erscheinen auch ohne Anmeldung als Kachel „PLATZ FREI“; im Bestand sind sechs von zehn Rängen so belegt. & Der Bestenliste sieht man den Verwaltungszustand der Redaktion an. & 1 Beg. & 2 & H2 & HE \\
	P22 & keine & Die Zahl der Reaktionen steht als Symbol mit Ziffer ohne Wort. & Beim Überfliegen bleibt offen, worauf sich die Zahl bezieht. & 1 Beg. & 1 & H6 & HE \\
	P23 & keine & Derselbe Leerzustand ist zweimal verschieden formuliert, einmal mit Sie, einmal im Infinitiv. & Zwei Listen beantworten dieselbe Lage mit zwei verschiedenen Sätzen. & 1 Beg. & 1 & H4 & HE \\
	P24 & keine & Projekttypen heißen in der deutschen Oberfläche teils englisch („Paper“, „Thesis“), teils deutsch („Reguläres Projekt“). & Wem die englischen Begriffe nicht geläufig sind, der muss den Projekttyp erschließen. & 1 Beg. & 2 & H2 & HE \\
	P25 & keine & Löschsymbole reagieren verschieden: auf der Metadaten- und Editorseite ohne rote Hervorhebung, anderswo rot; nur manche tragen einen Tooltip. & Dieselbe Aktion gibt verschiedene Rückmeldungen. & 1 Beg. & 1 & H4 & HE \\
	P26 & keine & Hinzufügen erscheint bei Kategorien und Semestern als alleinstehendes „+“, in der Studienstruktur als beschriftete englische Schaltfläche („Add faculty“). & Gleich aufgebaute Verwaltungsseiten zeigen dieselbe Aktion verschieden. & 1 Beg. & 2 & H4 & HE \\
	P27 & keine & Nach einer ungültigen Datumseingabe verschiebt die Fehlermeldung das Feld „To“ gegenüber den übrigen Filterelementen. & Die Filterzeile verliert im Fehlerzustand ihre Ausrichtung. & 1 Beg. & 1 & H8 & HE \\
	P28 & keine & Wer das Metadatenformular über den Seitenpfad oder die Zurück-Taste verlässt, erhält keine Warnung und verliert ungespeicherte Änderungen; über „Abbrechen“ erscheint eine Warnung. & Eingegebene Metadaten gehen unbeabsichtigt verloren, je nach Weg aus der Seite. & 1 Beg. & 3 & H5 & HE \\
	P29 & keine & Eine Projektadresse mit ungültiger Kennung zeigt „Die Eingaben sind unvollständig oder ungültig. Bitte prüfen Sie die markierten Felder.“, obwohl die Seite keine Felder hat. & Die Meldung sagt nicht, dass die Adresse zu keinem Projekt führt, und schlägt eine unmögliche Handlung vor. & 1 Beg. & 2 & H9 & HE \\
	P30 & keine & Der Platzhalter im Editor nennt den Schrägstrich als Weg zum Auswahlmenü der Bausteinarten; ein getippter Schrägstrich öffnet kein Menü, erst ein Umweg über ein Leerzeichen, und das Menü schließt sich von selbst. & Der angekündigte Weg zum Einfügen von Bausteinen ist nicht benutzbar. & 1 Beg. & 3 & H7 & HE \\
\end{longtable}
\endgroup

\chapter{Kennwerte des TAM-Bogens je Block}
\label{anh:tam-bloecke}

Der TAM-Bogen lag in vier Fassungen vor, eine je zuletzt bearbeiteten Block (Abschnitt~\ref{sec:testkonzept}). Block A steht hier für die beiden externen Personen, die allein Block A bearbeiteten; alle übrigen Personen füllten den Bogen nach ihrem Rollenblock aus. Die Werte beruhen auf zwei bis fünf Personen und dienen der Orientierung, nicht dem Vergleich zwischen Rollen. NFA-10 und NFA-11 werden an Block B gemessen (Abschnitt~\ref{sec:bewertungsrahmen}).

\begin{table}[htbp]
	\centering
	\footnotesize
	\caption[Kennwerte des TAM-Bogens je Block]{Kennwerte des TAM-Bogens je Block; SD ist die Stichproben-Standardabweichung (eigene Erhebung)}
	\label{tab:tam-bloecke}
	\begin{tabular}{@{}llrrrrrr@{}}
		\toprule
		\textbf{Block} & \textbf{Konstrukt} & \textbf{n} & \textbf{Median} & \textbf{Mittelwert} & \textbf{Min} & \textbf{Max} & \textbf{SD} \\
		\midrule
		A, nur extern & PU & 2 & 6,92 & 6,92 & 6,83 & 7,00 & 0,12 \\
		 & PEOU & 2 & 6,92 & 6,92 & 6,83 & 7,00 & 0,12 \\
		 & BI & 2 & 7,00 & 7,00 & 7,00 & 7,00 & 0,00 \\
		\midrule
		B, Studierende & PU & 5 & 5,17 & 5,27 & 4,50 & 6,33 & 0,81 \\
		 & PEOU & 5 & 6,17 & 5,80 & 4,17 & 6,67 & 1,02 \\
		 & BI & 5 & 6,00 & 5,90 & 4,00 & 7,00 & 1,14 \\
		\midrule
		C, Redaktion & PU & 2 & 6,50 & 6,50 & 6,00 & 7,00 & 0,71 \\
		 & PEOU & 2 & 5,75 & 5,75 & 5,33 & 6,17 & 0,59 \\
		 & BI & 2 & 7,00 & 7,00 & 7,00 & 7,00 & 0,00 \\
		\midrule
		D, Administration & PU & 2 & 6,25 & 6,25 & 6,00 & 6,50 & 0,35 \\
		 & PEOU & 2 & 5,42 & 5,42 & 5,17 & 5,67 & 0,35 \\
		 & BI & 2 & 6,50 & 6,50 & 6,00 & 7,00 & 0,71 \\
		\bottomrule
	\end{tabular}
\end{table}

	% ---------------------------------------------------------------------------
% Anhang D: Testunterlagen
% Erzeugt aus Testmaterial/_generator (aufgaben.py, werte.py, p4_druck.py,
% BEZUG aus make_docs.py) mit thesis/werkzeuge/anhang_d.py.
% Nicht von Hand aendern, sondern neu erzeugen.
% ---------------------------------------------------------------------------

\chapter{Testunterlagen}
\label{anh:testunterlagen}

Die Unterlagen stehen hier im Wortlaut der gedruckten Fassungen, die in den elf Sitzungen verwendet wurden (Abschnitt~\ref{sec:testdesign}). Jede Aufgabenkarte war eine eigene Seite, auf der allein der Kartentext stand; Startzustand und Erfolgskriterium kannte nur die Testleitung. Zeitbudget und Kennzahlen je Aufgabe nennt Tabelle~\ref{tab:aufgaben}.

\section*{Aufgabenkarten}

\begingroup
\scriptsize
\setlength{\tabcolsep}{3pt}
\renewcommand{\arraystretch}{1.12}
\setlength{\LTcapwidth}{\textwidth}
\begin{longtable}{@{}L{0.6cm}L{6.1cm}L{3.8cm}L{4.2cm}@{}}
	\caption[Aufgabenkarten Block A: ohne Anmeldung]{Aufgabenkarten Block A: ohne Anmeldung, Kartentext wortgleich mit Startzustand und Erfolgskriterium der Testleitung (eigene Darstellung)}
	\label{tab:karten-a}\\
	\toprule
	\textbf{ID} & \textbf{Kartentext} & \textbf{Startzustand} & \textbf{Erfolgskriterium} \\
	\midrule
	\endfirsthead
	\toprule
	\textbf{ID} & \textbf{Kartentext} & \textbf{Startzustand} & \textbf{Erfolgskriterium} \\
	\midrule
	\endhead
	\bottomrule
	\endfoot
	A1 & „Du bist gerade auf der Startseite dieses Portals gelandet, ohne dass dir jemand erklärt hat, was das ist. Schau dich um und sag laut, was du vermutest: Was ist das für eine Seite, wer betreibt sie, und was kann man hier tun?“ & Startseite, abgemeldet, frisches Browserprofil, kein Scrollen vorbereitet. Die Top 10 ist vollständig belegt. & Keine Erfolgswertung und keine Zeitnahme gegen ein Kriterium. Erhoben wird, ob Zweck, Betreiber und mindestens eine Handlungsmöglichkeit ohne Hilfe benannt werden. \\
	A2 & „Eine Bekannte hat dir erzählt, dass an der Hochschule ein Projekt zum Thema Baustoffbörse entstanden ist. Finde dieses Projekt.“ & Startseite, abgemeldet. & Detailseite des Zielprojekts ist geöffnet, Hilfestufe höchstens 2. \\
	A3 & „Jetzt möchtest du dir einen Überblick verschaffen. Zeig mir alle Bachelorarbeiten der Fakultät Elektrotechnik und Informatik aus dem Semester WiSe 2025/26.“ & Projektliste, keine Filter gesetzt, abgemeldet. & Alle drei Einschränkungen gleichzeitig aktiv (Abschluss, Fakultät, Semester), Trefferliste sichtbar. \\
	A4 & „Die Liste ist dir jetzt zu eng geworden. Erweitere sie wieder, ohne von vorn anzufangen. Öffne danach ein Projekt, das dich interessiert, und finde heraus, was dieselbe Person sonst noch veröffentlicht hat.“ & Ergebnis aus A3 unverändert stehen lassen. & Mindestens eine Einschränkung entfernt, ohne Neuladen der Seite und ohne Rückkehr zur Startseite; danach ist die Portfolio-Profilseite der Autorin oder des Autors geöffnet. \\
	A5 & „Du willst dieses Projekt jemandem zeigen, der gerade kein Internet hat. Sichere es dir so, dass du es weitergeben kannst.“ & Das Projekt „Mehrwegpfand in der Bremer Gastronomie“ (aus A2 oder A4) geöffnet stehen lassen, abgemeldet. & Eine PDF-Datei ist erzeugt und gespeichert, ohne Hilfestufe 3. \\
\end{longtable}
\endgroup

\begingroup
\scriptsize
\setlength{\tabcolsep}{3pt}
\renewcommand{\arraystretch}{1.12}
\setlength{\LTcapwidth}{\textwidth}
\begin{longtable}{@{}L{0.6cm}L{6.1cm}L{3.8cm}L{4.2cm}@{}}
	\caption[Aufgabenkarten Block B: Studierende]{Aufgabenkarten Block B: Studierende, Kartentext wortgleich mit Startzustand und Erfolgskriterium der Testleitung (eigene Darstellung)}
	\label{tab:karten-b}\\
	\toprule
	\textbf{ID} & \textbf{Kartentext} & \textbf{Startzustand} & \textbf{Erfolgskriterium} \\
	\midrule
	\endfirsthead
	\toprule
	\textbf{ID} & \textbf{Kartentext} & \textbf{Startzustand} & \textbf{Erfolgskriterium} \\
	\midrule
	\endhead
	\bottomrule
	\endfoot
	B1 & „Hier sind deine Zugangsdaten. Melde dich am Portal an.“ \newline \textit{Darunter je Sitzung: E-Mail und Passwort des Testkontos.} & Abgemeldet, Projektliste geöffnet. Auf der Karte stehen die Anmeldeadresse und das Passwort. & Angemeldet, höchstens zwei Fehlversuche. \\
	B2 & „Du möchtest deine Arbeit aus dem letzten Semester im Portal zeigen. Lege sie hier an. Sie heißt ‚Lastenrad-Verleih im Bremer Westen‘.“ & Angemeldet als Studierende, Startseite. & Ein Projekt im Status Entwurf existiert und trägt den vorgegebenen Titel. \\
	B3 & „Baue die erste Seite deines Projekts: eine Überschrift, einen Absatz Text, ein Bild und eine Aufzählung mit drei Punkten. Was genau dort steht, ist egal.“ & Leerer Editor des in B2 angelegten Projekts; eine Bilddatei liegt sichtbar auf dem Schreibtisch des Testrechners. & Überschrift, Absatz, Bild und Aufzählung sind im Dokument, die Speicherung ist bestätigt. \\
	B4 & „Dein Projekt soll zur Prüfung eingereicht werden. Trage alles ein, was dafür nötig ist.“ & Ergebnis aus B3 unverändert stehen lassen: Projekt mit Inhalt, keine Metadaten gesetzt. & Alle sieben Pflichtangaben gesetzt: Projekttyp, Fakultät, Studiengang, Abschluss, Kategorie, Semester und mindestens eine zugewiesene Redaktion. \\
	B5 & „Gib dein Projekt zur Prüfung ab.“ & Ergebnis aus B4 unverändert stehen lassen: Projekt vollständig, Status Entwurf. & Status Eingereicht. \\
	B6 & „Dir fällt ein, dass ein Absatz fehlt. Sorge dafür, dass du wieder weiterschreiben kannst.“ & Ergebnis aus B5 unverändert stehen lassen: Projekt im Status Eingereicht. & Projekt wieder im Status Entwurf, Editor geöffnet. \\
	B7 & „Zum Schluss: Ergänze in deinem eigenen Profil eine kurze Beschreibung von dir, damit Besucher wissen, wer hinter den Projekten steht.“ & Angemeldet, eigenes Profil ohne Kurzbiografie. & Kurzbiografie gespeichert. \\
\end{longtable}
\endgroup

\begingroup
\scriptsize
\setlength{\tabcolsep}{3pt}
\renewcommand{\arraystretch}{1.12}
\setlength{\LTcapwidth}{\textwidth}
\begin{longtable}{@{}L{0.6cm}L{6.1cm}L{3.8cm}L{4.2cm}@{}}
	\caption[Aufgabenkarten Block C: Redaktion]{Aufgabenkarten Block C: Redaktion, Kartentext wortgleich mit Startzustand und Erfolgskriterium der Testleitung (eigene Darstellung)}
	\label{tab:karten-c}\\
	\toprule
	\textbf{ID} & \textbf{Kartentext} & \textbf{Startzustand} & \textbf{Erfolgskriterium} \\
	\midrule
	\endfirsthead
	\toprule
	\textbf{ID} & \textbf{Kartentext} & \textbf{Startzustand} & \textbf{Erfolgskriterium} \\
	\midrule
	\endhead
	\bottomrule
	\endfoot
	C1 & „Hier sind deine Zugangsdaten. Melde dich am Portal an.“ \newline \textit{Darunter je Sitzung: E-Mail und Passwort des Testkontos.} & Abgemeldet, Projektliste geöffnet. Auf der Karte stehen die Anmeldeadresse und das Passwort. & Angemeldet, höchstens zwei Fehlversuche. \\
	C2 & „Du bist in der Redaktion für einige eingereichte Projekte zuständig. Öffne deine Reviews und verschaffe dir einen Überblick, was gerade auf dich wartet.“ & Angemeldet als Redaktion, Startseite; vier Projekte warten im Arbeitsvorrat. Nach C1 steht „Meine Reviews“ schon offen: die Moderation öffnet die Startseite, bevor sie die Karte reicht. & Die Liste der eigenen Reviews ist geöffnet, mindestens ein wartendes Projekt wird benannt. \\
	C3 & „In deinen Reviews kannst du an jedem Projekt drei Arten von Kommentaren schreiben. Eine Antwort ist eine Rückfrage, die alle lesen, die das Projekt sehen dürfen. Eine Notiz ist intern: nur Redaktion und Administration sehen sie. Eine Änderung fordert die Person, die das Projekt eingereicht hat, auf, etwas zu überarbeiten, und gibt ihr das Projekt dafür zurück. Deine Aufgabe: Öffne das Projekt ‚Zahlungsbereitschaft für regionale Milch in Bremen‘. Dir fällt auf, dass die Quellenangaben fehlen. Sorge dafür, dass der Autor das nachträgt.“ & Seed-Projekt „Zahlungsbereitschaft für regionale Milch in Bremen“ im Status Eingereicht, die Testperson ist ihm als Redaktion zugewiesen. & Status Überarbeitung, Begründung gespeichert und für den Autor sichtbar. \\
	C4 & „Auf dem Weg zur Veröffentlichung durchläuft ein Projekt vier Stände. Wer es schreibt, legt es als Entwurf an und reicht es ein, dann steht es auf Eingereicht. Die Redaktion prüft und genehmigt es, dann steht es auf Genehmigt. Erst wenn die Redaktion es danach veröffentlicht, steht es auf Veröffentlicht und ist für alle sichtbar. Deine Aufgabe: Das Projekt ‚Wiederverwendung von Ziegeln aus einem Abbruch in Walle‘ ist in Ordnung. Veröffentliche es.“ & Zweites Seed-Projekt „Wiederverwendung von Ziegeln aus einem Abbruch in Walle“ im Status Eingereicht, die Testperson ist ihm als Redaktion zugewiesen. & Status Veröffentlicht, ohne Hilfestufe 3. \\
	C5 & „Das Projekt, das du gerade veröffentlicht hast, gefällt dir besonders gut. Nominiere es für die Top 10 auf der Startseite.“ & Das in C4 veröffentlichte Projekt „Wiederverwendung von Ziegeln aus einem Abbruch in Walle“ ist noch in der Prüfansicht geöffnet; es trägt keine offene Nominierung und keinen Top-10-Platz. In der Top 10 sind sechs der zehn Plätze belegt. & Nominierung mit Stufe ist abgegeben. \\
	C6 & „Mit dem veröffentlichten Projekt bist du fertig. Öffne jetzt ein anderes Projekt, ‚Audiowalk zur Werftgeschichte in Gröpelingen‘, und halte daran eine Notiz fest, die nur die Redaktion sehen soll und nicht die Autorin.“ & Drittes Seed-Projekt „Audiowalk zur Werftgeschichte in Gröpelingen“ im Status Eingereicht, die Testperson ist ihm als Redaktion zugewiesen. Geöffnet ist noch das in C4 veröffentlichte und in C5 nominierte Projekt. & Eine Notiz ist angelegt, keine Antwort und keine Änderung. \\
\end{longtable}
\endgroup

\begingroup
\scriptsize
\setlength{\tabcolsep}{3pt}
\renewcommand{\arraystretch}{1.12}
\setlength{\LTcapwidth}{\textwidth}
\begin{longtable}{@{}L{0.6cm}L{6.1cm}L{3.8cm}L{4.2cm}@{}}
	\caption[Aufgabenkarten Block D: Administration]{Aufgabenkarten Block D: Administration, Kartentext wortgleich mit Startzustand und Erfolgskriterium der Testleitung (eigene Darstellung)}
	\label{tab:karten-d}\\
	\toprule
	\textbf{ID} & \textbf{Kartentext} & \textbf{Startzustand} & \textbf{Erfolgskriterium} \\
	\midrule
	\endfirsthead
	\toprule
	\textbf{ID} & \textbf{Kartentext} & \textbf{Startzustand} & \textbf{Erfolgskriterium} \\
	\midrule
	\endhead
	\bottomrule
	\endfoot
	D1 & „Hier sind deine Zugangsdaten. Melde dich am Portal an.“ \newline \textit{Darunter je Sitzung: E-Mail und Passwort des Testkontos.} & Abgemeldet, Projektliste geöffnet. Auf der Karte stehen die Anmeldeadresse und das Passwort. & Angemeldet, höchstens zwei Fehlversuche. \\
	D2 & „Du betreust dieses Portal. Öffne die Administrationsseite und verschaff dir dort einen Überblick, was gerade im Portal los ist.“ & Angemeldet als Administration, Startseite des Portals (nicht der Administrationsseite). Nach D1 steht die Administrationsseite schon offen: die Moderation öffnet die Startseite, bevor sie die Karte reicht. & Die Administrationsseite ist geöffnet, und mindestens zwei konkrete Aussagen über den Zustand des Portals werden genannt. \\
	D3 & „Nele Ostermann soll künftig Projekte prüfen dürfen. Ändere dafür die Rolle des Kontos auf Redakteur:in.“ & Administrationsseite offen; das Zielkonto trägt die Rolle STUDENT und eine studentische HSB-Adresse (@stud.hs-bremen.de). Die Hochstufung ist der Tutor-Fall aus Contract §3.1.1: eine studentische Adresse belegt den Studierendenstatus, nicht die redaktionelle Tätigkeit. & Das Konto trägt die Rolle EDITOR. \\
	D4 & „Die Kategorien ‚Studienprojekt‘ und ‚Studentisches Projekt‘ meinen dasselbe. Führe die beiden zu einer Kategorie zusammen, ohne dass Projekte verloren gehen.“ & Die zwei gleichbedeutenden Kategorien „Studienprojekt“ und „Studentisches Projekt“ mit je mindestens drei Projekten daran. Daneben stehen ähnlich klingende Stammkategorien wie „Semesterprojekt“; deshalb nennt die Karte beide Namen. & Nur noch eine der beiden Kategorien existiert, und alle betroffenen Projekte hängen daran. \\
	D5 & „Die fünf genehmigten Projekte sollen veröffentlicht werden. Erledige das.“ & Projektverwaltung offen, fünf Projekte im Status Genehmigt, über den Statusfilter auffindbar. & Alle fünf Projekte tragen den Status Veröffentlicht. \\
	D6 & „Die Redaktion hat Projekte für die Top 10 der Startseite vorgeschlagen. Nimm aus diesen Vorschlägen Projekte in die Top 10 auf, bis alle zehn Plätze belegt sind.“ & Mindestens zwölf offene Nominierungen, keine Top 10 gesetzt. & Alle zehn Plätze sind belegt; die Verwaltung zeigt „10 von 10 Plätzen belegt“. Einen eigenen Speicherschritt gibt es nicht. \\
	D7 & „Jemand behauptet, die Top 10 der Startseite sei vor deiner Arbeit daran schon einmal geändert worden. Finde heraus, wann das zuletzt war und wer es war.“ & Ergebnis aus D6 unverändert stehen lassen: Top 10 gesetzt, im Änderungsprotokoll stehen mehrere Einträge von verschiedenen Personen. & Das Änderungsprotokoll ist geöffnet; die eigenen Einträge aus D6 werden als solche erkannt, und der jüngste Eintrag davor wird mit Person und Zeitpunkt genannt. \\
\end{longtable}
\endgroup

\section*{Fragebogen zur Gebrauchstauglichkeit (SUS)}

Deutsche Übersetzung über dem englischen Original nach Brooke (1996). Einleitung auf dem Bogen: „Mit ‚dem System‘ ist in allen zehn Aussagen das Portal gemeint, das du gerade benutzt hast. Bitte kreuze in jeder Zeile genau ein Feld an. Es gibt kein richtig und kein falsch.“ Skala: 1 = stimme überhaupt nicht zu, 5 = stimme voll zu.

\begingroup
\footnotesize
\setlength{\tabcolsep}{4pt}
\renewcommand{\arraystretch}{1.15}
\setlength{\LTcapwidth}{\textwidth}
\begin{longtable}{@{}L{0.7cm}L{9.6cm}*{5}{>{\centering\arraybackslash}p{0.62cm}}@{}}
	\caption[SUS-Bogen]{SUS-Bogen in der eingesetzten deutschen Fassung mit englischem Original (nach Brooke, 1996)}
	\label{tab:sus-bogen}\\
	\toprule
	\textbf{Nr.} & \textbf{Aussage} & \textbf{1} & \textbf{2} & \textbf{3} & \textbf{4} & \textbf{5} \\
	\midrule
	\endfirsthead
	\toprule
	\textbf{Nr.} & \textbf{Aussage} & \textbf{1} & \textbf{2} & \textbf{3} & \textbf{4} & \textbf{5} \\
	\midrule
	\endhead
	\bottomrule
	\endfoot
	1 & Ich denke, dass ich das System gerne regelmäßig nutzen würde. \newline {\scriptsize\itshape I think that I would like to use this system frequently} & $\square$ & $\square$ & $\square$ & $\square$ & $\square$ \\
	2 & Ich fand das System unnötig komplex. \newline {\scriptsize\itshape I found the system unnecessarily complex} & $\square$ & $\square$ & $\square$ & $\square$ & $\square$ \\
	3 & Ich fand das System einfach zu benutzen. \newline {\scriptsize\itshape I thought the system was easy to use} & $\square$ & $\square$ & $\square$ & $\square$ & $\square$ \\
	4 & Ich glaube, ich würde die Unterstützung einer technisch versierten Person benötigen, um das System zu benutzen. \newline {\scriptsize\itshape I think that I would need the support of a technical person to use this system} & $\square$ & $\square$ & $\square$ & $\square$ & $\square$ \\
	5 & Ich fand, die verschiedenen Funktionen des Systems waren gut integriert. \newline {\scriptsize\itshape I found the various functions in this system were well integrated} & $\square$ & $\square$ & $\square$ & $\square$ & $\square$ \\
	6 & Ich denke, das System enthielt zu viele Inkonsistenzen. \newline {\scriptsize\itshape I thought there was too much inconsistency in this system} & $\square$ & $\square$ & $\square$ & $\square$ & $\square$ \\
	7 & Ich kann mir vorstellen, dass die meisten Menschen den Umgang mit diesem System sehr schnell lernen. \newline {\scriptsize\itshape I would imagine that most people would learn to use this system very quickly} & $\square$ & $\square$ & $\square$ & $\square$ & $\square$ \\
	8 & Ich fand das System sehr umständlich zu nutzen. \newline {\scriptsize\itshape I found the system very cumbersome to use} & $\square$ & $\square$ & $\square$ & $\square$ & $\square$ \\
	9 & Ich fühlte mich bei der Nutzung des Systems sehr sicher. \newline {\scriptsize\itshape I felt very confident using the system} & $\square$ & $\square$ & $\square$ & $\square$ & $\square$ \\
	10 & Ich musste eine Menge lernen, bevor ich mit dem System zurechtkam. \newline {\scriptsize\itshape I needed to learn a lot of things before I could get going with this system} & $\square$ & $\square$ & $\square$ & $\square$ & $\square$ \\
\end{longtable}
\endgroup

\section*{Fragebogen zur Einschätzung des Portals (TAM)}

Deutsche Übersetzung über dem englischen Original nach Davis (1989); [X] steht im Original für den Namen des Systems. Für BI1 und BI2 stand kein englisches Original auf dem Bogen. Der Bogen lag in vier Fassungen vor, die sich allein im Bezugssatz über den Aussagen unterscheiden:

\begin{itemize}[topsep=0.3em, itemsep=0.2em]
	\item Fassung A (Studentische Projekte suchen und ansehen): „Denk beim Beantworten an das Suchen und Ansehen studentischer Projekte, so wie du es eben getan hast.“
	\item Fassung B (Eigene Projektarbeiten erstellen und einreichen): „Denk beim Beantworten an das Erstellen und Einreichen eigener Projektarbeiten, so wie du es eben getan hast.“
	\item Fassung C (Eingereichte Projekte prüfen und veröffentlichen): „Denk beim Beantworten an das Prüfen und Veröffentlichen eingereichter Projekte, so wie du es eben getan hast.“
	\item Fassung D (Das Portal verwalten): „Denk beim Beantworten an die Verwaltung des Portals, so wie du es eben getan hast.“
\end{itemize}

Einleitung auf dem Bogen: „Vierzehn Aussagen. Kreuze in jeder Zeile genau ein Feld an, je nachdem, wie wahrscheinlich die Aussage für dich ist. Es gibt keine richtigen und keine falschen Antworten.“ Die Felder trugen auf dem Bogen keine Zahlen, sondern Beschriftungen; in Teil 1, Teil 2 und bei BI1 von links: äußerst wahrscheinlich, ziemlich wahrscheinlich, etwas wahrscheinlich, weder noch, etwas unwahrscheinlich, ziemlich unwahrscheinlich, äußerst unwahrscheinlich. Kodiert wurde nach Position, das linke Feld mit 7 und das rechte mit 1, sodass hohe Werte Zustimmung bedeuten; die Spaltenköpfe der Tabelle geben diese Kodierung wieder. Bei BI2 sind die Endpunkte vertauscht gedruckt (links nicht zu erwarten, rechts zu erwarten); der Wert ging deshalb als 8 minus Rohwert ein.

\begingroup
\footnotesize
\setlength{\tabcolsep}{3pt}
\renewcommand{\arraystretch}{1.15}
\setlength{\LTcapwidth}{\textwidth}
\begin{longtable}{@{}L{1.15cm}L{8.45cm}*{7}{>{\centering\arraybackslash}p{0.5cm}}@{}}
	\caption[TAM-Bogen]{TAM-Bogen in der eingesetzten deutschen Fassung mit englischem Original (nach Davis, 1989)}
	\label{tab:tam-bogen}\\
	\toprule
	\textbf{Item} & \textbf{Aussage} & \textbf{7} & \textbf{6} & \textbf{5} & \textbf{4} & \textbf{3} & \textbf{2} & \textbf{1} \\
	\midrule
	\endfirsthead
	\toprule
	\textbf{Item} & \textbf{Aussage} & \textbf{7} & \textbf{6} & \textbf{5} & \textbf{4} & \textbf{3} & \textbf{2} & \textbf{1} \\
	\midrule
	\endhead
	\bottomrule
	\endfoot
	\multicolumn{9}{@{}l}{\textbf{Teil 1: Nutzen für deine Aufgaben}} \\
	PU1 & Die Nutzung des Portals würde es mir ermöglichen, diese Aufgaben schneller zu erledigen. \newline {\scriptsize\itshape Using [X] in my job would enable me to accomplish tasks more quickly.} & $\square$ & $\square$ & $\square$ & $\square$ & $\square$ & $\square$ & $\square$ \\
	PU2 & Die Nutzung des Portals würde meine Leistung bei diesen Aufgaben verbessern. \newline {\scriptsize\itshape Using [X] would improve my job performance.} & $\square$ & $\square$ & $\square$ & $\square$ & $\square$ & $\square$ & $\square$ \\
	PU3 & Die Nutzung des Portals würde meine Produktivität bei diesen Aufgaben erhöhen. \newline {\scriptsize\itshape Using [X] in my job would increase my productivity.} & $\square$ & $\square$ & $\square$ & $\square$ & $\square$ & $\square$ & $\square$ \\
	PU4 & Die Nutzung des Portals würde meine Effektivität bei diesen Aufgaben steigern. \newline {\scriptsize\itshape Using [X] would enhance my effectiveness on the job.} & $\square$ & $\square$ & $\square$ & $\square$ & $\square$ & $\square$ & $\square$ \\
	PU5 & Die Nutzung des Portals würde mir diese Aufgaben erleichtern. \newline {\scriptsize\itshape Using [X] would make it easier to do my job.} & $\square$ & $\square$ & $\square$ & $\square$ & $\square$ & $\square$ & $\square$ \\
	PU6 & Ich würde das Portal für diese Aufgaben nützlich finden. \newline {\scriptsize\itshape I would find [X] useful in my job.} & $\square$ & $\square$ & $\square$ & $\square$ & $\square$ & $\square$ & $\square$ \\
	\midrule
	\multicolumn{9}{@{}l}{\textbf{Teil 2: Umgang mit dem Portal}} \\
	PEOU1 & Es würde mir leichtfallen, den Umgang mit dem Portal zu erlernen. \newline {\scriptsize\itshape Learning to operate [X] would be easy for me.} & $\square$ & $\square$ & $\square$ & $\square$ & $\square$ & $\square$ & $\square$ \\
	PEOU2 & Ich würde es leicht finden, das Portal dazu zu bringen, das zu tun, was ich möchte. \newline {\scriptsize\itshape I would find it easy to get [X] to do what I want it to do.} & $\square$ & $\square$ & $\square$ & $\square$ & $\square$ & $\square$ & $\square$ \\
	PEOU3 & Meine Interaktion mit dem Portal wäre klar und verständlich. \newline {\scriptsize\itshape My interaction with [X] would be clear and understandable.} & $\square$ & $\square$ & $\square$ & $\square$ & $\square$ & $\square$ & $\square$ \\
	PEOU4 & Ich würde den Umgang mit dem Portal als flexibel empfinden. \newline {\scriptsize\itshape I would find [X] to be flexible to interact with.} & $\square$ & $\square$ & $\square$ & $\square$ & $\square$ & $\square$ & $\square$ \\
	PEOU5 & Es würde mir leichtfallen, im Umgang mit dem Portal sicher zu werden. \newline {\scriptsize\itshape It would be easy for me to become skillful at using [X].} & $\square$ & $\square$ & $\square$ & $\square$ & $\square$ & $\square$ & $\square$ \\
	PEOU6 & Ich würde das Portal einfach zu bedienen finden. \newline {\scriptsize\itshape I would find [X] easy to use.} & $\square$ & $\square$ & $\square$ & $\square$ & $\square$ & $\square$ & $\square$ \\
	\midrule
	\multicolumn{9}{@{}l}{\textbf{Teil 3: Künftige Nutzung}} \\
	BI1 & Angenommen, das Portal stünde mir zur Verfügung: Ich sage voraus, dass ich es künftig regelmäßig nutzen werde. & $\square$ & $\square$ & $\square$ & $\square$ & $\square$ & $\square$ & $\square$ \\
	BI2 & Angenommen, das Portal stünde mir zur Verfügung: Ich sage voraus, dass ich es künftig regelmäßig nutzen werde. (Endpunkte vertauscht gedruckt, Wert = 8 minus Rohwert) & $\square$ & $\square$ & $\square$ & $\square$ & $\square$ & $\square$ & $\square$ \\
\end{longtable}
\endgroup


\end{document}
